\documentclass[11pt]{article}

\usepackage[T1]{fontenc}
\usepackage{lmodern}
\usepackage{amsmath,amssymb,amsthm,mathtools}
\usepackage{booktabs}
\usepackage[margin=1in]{geometry}
\usepackage[colorlinks=true,linkcolor=blue,citecolor=blue,urlcolor=blue]{hyperref}

\newtheorem{theorem}{Theorem}[section]
\newtheorem{proposition}[theorem]{Proposition}
\newtheorem{lemma}[theorem]{Lemma}
\newtheorem{corollary}[theorem]{Corollary}
\newtheorem{conjecture}[theorem]{Conjecture}
\theoremstyle{remark}
\newtheorem{remark}[theorem]{Remark}

\newcommand{\R}{\mathbb R}
\newcommand{\N}{\mathbb N}
\newcommand{\sbr}[1]{[#1]_{<}}
\newcommand{\Ad}{A_{\rho,1}^{(d)}}
\newcommand{\Aact}{A_{\rho,K}}
\newcommand{\dd}{\,\mathrm d}

\title{The Dirichlet P\'olya Problem for Spherical Shells in General
Dimension:\\Exact Dimension Lifting and the Shifted-Tail Bottleneck}
\author{Working manuscript}
\date{18 July 2026}

\begin{document}
\maketitle

\begin{abstract}
For a spherical shell in dimension $d\geq3$, the separated radial problem
has Bessel orders $\ell+(d-2)/2$ and spherical-harmonic multiplicities
$\kappa_{d,\ell}$.  We give an exact branching decomposition of the strict
shell-action proxy into positive integer- and half-integer shifted planar
tails.  We also prove a signed primitive comparison and retain its full
positive remainder.  This yields the exact identity
\[
 W_d-P_d^{<}=\mathcal B_d(A)+
 \sum_{m\geq0}\binom{m+d-3}{d-3}D_A\left(m+\frac{d-2}{2}\right).
\]
Consequently, one dimension-free positive shifted-tail inequality implies
the sharp Dirichlet P\'olya inequality for every spherical shell and every
$d\geq3$.  The dimension lift is proved here; the inner shifted-tail
inequality remains the central open lemma.  We prove the outer and terminal
tail regions, close the first interface-cell milestone by exact rational
wall estimates, and prove all shifted tails above an explicit fixed-ratio
high-energy threshold.  A wall-safe first-shelf reduction compresses every
remaining many-cell tail to one scalar coupled to its terminal
one-interface reserve.  In the thin sector it also improves the uniform
threshold to $K\geq125/[8(1-\rho)^2]$.  A scale-free shelf-curvature
bound and a fractional-top inverse-action estimate give an explicit
exact-angle lower bound for the terminal reserve, including the previously
unused zero-level triangle.  We also prove an automatic ultra-thin
 high-floor sector and exclude every noncompact escape in both the
 first-floor and high-floor first-drop branches.  A global endpoint-chord
 and terminal-payment argument first excludes every negative high-floor
 first drop with $\rho>99/100$, without a restriction on the frequency or
 on terminal walls.  A normalized positive-arccos refinement first extends
 this exclusion to $\rho\geq603/625$.  A terminal stratification on the
 zero-level face then pushes it to $\rho\geq477/500$.  The active-width
 relation removes the resulting zero-count obstruction, and a sharp
 terminal-cap estimate extends the exclusion to $\rho\geq93/100$.
 A refined absent-level chord and the automatic $q\geq5/2$ terminal
 start first extend it to $\rho\geq927/1000$.  Switching from the lower
 bound on the shelf length to its geometric upper bound when the common
 prefix changes sign extends the exclusion further to
 $\rho\geq1847/2000$.  Retaining instead the exact correlated negative
 prefix, sharpening the active width, and improving the small-level
 distance extends the exclusion to $\rho\geq23/25$.
 Combined with the fixed-ratio scalar estimate, this reduces every
 residual high-floor candidate to
 \[
  \frac1{3154914}<\rho<\frac{23}{25},\qquad
  \frac{21}{4}<K<1577457.
 \]
 This is a sharper compact reduction, not a proof of the residual
 fixed-ratio walls.  A directed-rounding
 certificate proves all $19602$ finite first-floor chambers with
 $r\geq3/2$, $s\in\{0,1\}$, $x<100$.  A second certificate handles all
 $398$ finite small-start chambers, while an analytic comparison closes
 their infinite rays.  Together with the continuous large-ray certificate,
 this proves the entire first-floor first-drop sector $s\in\{0,1\}$.  A scale-free
 true-interface argument, completed by a directed-rounding moving-boundary
 certificate, proves the complementary cone $s\geq2$ throughout
 $K-\mu\geq\pi$.  Combined with the no-mode theorem below the active wall,
 this closes the whole ordinary first-floor first-drop branch needed for the
 spectral theorem.  In the no-drop branch, directed-rounding certificates
 now prove the reduced scalar strictly positive on every floor $f\geq2$:
 the two low-floor rows use $2268$ phase chambers, and an eight-shard
 certificate covers every $f\geq4$ residual.  The zero floor is automatic.
 The first-floor no-drop case $f=1$, $p=n$, and the residual high-floor
 first-drop compact walls remain uncertified; no assertion in this
 manuscript treats the full residual terminal-shelf scalar as proved.
\end{abstract}

\section{Statement and status}

For $0<\rho<1$ let
\[
 A_{\rho,1}^{(d)}=\{x\in\R^d:\rho<|x|<1\},
 \qquad d\geq3,
\]
and use the strict counting convention
\[
 N_D(A_{\rho,1}^{(d)},K^2)
 :=\#\{j:\lambda_j^D(A_{\rho,1}^{(d)})<K^2\}.
\]
Write
\[
 \omega_d=\frac{\pi^{d/2}}{\Gamma(d/2+1)},
 \qquad L_d=\frac{\omega_d}{(2\pi)^d}.
\]
The desired estimate is
\begin{equation}
 N_D(A_{\rho,1}^{(d)},K^2)
 \leq L_d|A_{\rho,1}^{(d)}|K^d
 =\frac{1-\rho^d}{2^d\Gamma(d/2+1)^2}K^d.
 \label{eq:polya-target}
\end{equation}

For $a>0$, define the zero-extended ball action
\begin{equation}
 G_a(z)=
 \begin{cases}
 \displaystyle\frac1\pi\left(\sqrt{a^2-z^2}
 -z\arccos\frac za\right),&0\leq z\leq a,\\[1mm]
 0,&z>a,
 \end{cases}
 \label{eq:ball-action}
\end{equation}
put $\mu=\rho K$, and set
\begin{equation}
 A(z)=A_{\rho,K}(z):=G_K(z)-G_\mu(z).
 \label{eq:shell-action}
\end{equation}
For $x>0$, let
\[
 \sbr{x}:=\#\{n\in\N:n<x\}=\max\{0,\lceil x\rceil-1\}.
\]
In particular, $\sbr{n}=n-1$ at a positive integral wall.

The only unproved statement on the primary path is the following.

\begin{conjecture}[Positive shifted shell tails]
\label{conj:shifted-tail}
For $0<\rho<1$, $K>0$, and
\[
 r\in\tfrac12\N,
\]
one has
\begin{equation}
 T_A(r):=\sbr{A(r)+\tfrac14}
 +2\sum_{j\geq1}\sbr{A(r+j)+\tfrac14}
 \leq2\int_r^K A(z)\dd z.
 \label{eq:ST}
\end{equation}
The zero extension makes the sum finite and makes the assertion trivial
when $r\geq K$.
\end{conjecture}

Sections~\ref{sec:spectral}--\ref{sec:defect} prove that
Conjecture~\ref{conj:shifted-tail} implies \eqref{eq:polya-target} in every
dimension $d\geq3$.  Section~\ref{sec:tail-progress} records the proved
 tail regions, the completed first interface-cell theorem, fixed-ratio and
 thin high-energy theorems, the terminal reserve, and the certified no-drop
 floors $f\geq2$.  It also states explicitly which wall problems remain.
 The distinct first-floor no-drop case $f=1$, $p=n$, is not covered by those
 certificates.
The sharpest high-floor ratio certificate currently proved excludes
$\rho\geq23/25$ and leaves exactly
\[
 \frac1{3154914}<\rho<\frac{23}{25},\qquad
 \frac{21}{4}<K<1577457.
\]
The analytic walls in this residual box are not certified here.

\section{General-dimensional spectral proxy}
\label{sec:spectral}

Put
\[
 a_d=\frac{d-2}{2},\qquad \nu_{\ell,d}=\ell+a_d.
\]
The degree-$\ell$ spherical harmonics on $\mathbb S^{d-1}$ have
multiplicity
\begin{equation}
 \kappa_{d,\ell}
 =\binom{\ell+d-1}{d-1}-\binom{\ell+d-3}{d-1}
 =\frac{(2\ell+d-2)(\ell+d-3)!}{\ell!(d-2)!}.
 \label{eq:harmonic-multiplicity}
\end{equation}
Here and below a binomial coefficient with lower index larger than its
nonnegative upper index is zero.

Under the unitary radial substitution
\[
 u(r)=r^{(d-1)/2}R(r),
\]
the degree-$\ell$ block becomes
\begin{equation}
 L_{\ell,d}=-\frac{\mathrm d^2}{\mathrm dr^2}
 +\frac{\nu_{\ell,d}^2-\frac14}{r^2}
 \quad\hbox{on }(\rho,1),
 \qquad u(\rho)=u(1)=0.
 \label{eq:radial-operator}
\end{equation}
Its positive eigenfrequencies are the zeros of
\[
 D_{\nu,\rho}(k)
 =J_\nu(\rho k)Y_\nu(k)-Y_\nu(\rho k)J_\nu(k),
 \qquad \nu=\nu_{\ell,d}.
\]
With the standard continuous Bessel phase, let
\[
 \gamma_{\rho,K}(\nu)
 =\frac{\theta_\nu(K)-\theta_\nu(\rho K)}\pi.
\]
The determinant factorization and phase monotonicity give the exact strict
count
\begin{equation}
 N_D(A_{\rho,1}^{(d)},K^2)
 =\sum_{\ell\geq0}\kappa_{d,\ell}
 \sbr{\gamma_{\rho,K}(\nu_{\ell,d})}.
 \label{eq:exact-phase-count}
\end{equation}
The real-order phase-difference enclosure of
\cite{FLPSphase,FLPSannuli} gives, for $0\leq\nu\leq K$,
\begin{equation}
 \gamma_{\rho,K}(\nu)<A_{\rho,K}(\nu)+\frac14.
 \label{eq:phase-action}
\end{equation}
For $\nu\geq K$, both the exact count and the action proxy vanish.  Since
$x<y$ implies $\sbr{x}\leq\sbr{y}$, including when $y$ is integral,
\eqref{eq:exact-phase-count}--\eqref{eq:phase-action} imply the wall-safe
proxy bound
\begin{equation}
 N_D(A_{\rho,1}^{(d)},K^2)\leq P_d^{<}(\rho,K),
 \qquad
 P_d^{<}:=\sum_{\ell\geq0}\kappa_{d,\ell}q_\ell,
 \label{eq:strict-proxy}
\end{equation}
where
\begin{equation}
 q_\ell=\sbr{A(\ell+a_d)+\frac14}.
 \label{eq:q-def}
\end{equation}

The argument above is dimension-independent after the substitutions
$\nu=\ell+a_d$ and \eqref{eq:harmonic-multiplicity}.  In particular, the
published phase enclosure is already a real-order result: no interpolation
is introduced when $d$ changes parity.

\section{Exact harmonic branching}
\label{sec:branching}

Let $p=d-2$, so $a_d=p/2$, and define
\begin{equation}
 r_m=m+\frac p2,
 \qquad
 c_{d,m}=\binom{m+p-1}{p-1}=\binom{m+d-3}{d-3}.
 \label{eq:branching-weights}
\end{equation}

\begin{proposition}[Exact shifted-tail decomposition]
\label{prop:branching}
For every finitely supported sequence $(q_\ell)_{\ell\geq0}$,
\begin{equation}
 \sum_{\ell\geq0}\kappa_{d,\ell}q_\ell
 =\sum_{m\geq0}c_{d,m}
 \left(q_m+2\sum_{j\geq1}q_{m+j}\right).
 \label{eq:abstract-branching}
\end{equation}
Consequently, for \eqref{eq:q-def},
\begin{equation}
 \boxed{P_d^{<}=\sum_{m\geq0}c_{d,m}T_A(r_m).}
 \label{eq:proxy-tail-decomp}
\end{equation}
This identity is exact at every strict action wall.
\end{proposition}

\begin{proof}
The generating functions are
\[
 \sum_{\ell\geq0}\kappa_{d,\ell}s^\ell
 =\frac{1+s}{(1-s)^{d-1}},
 \qquad
 \sum_{m\geq0}c_{d,m}s^m=\frac1{(1-s)^{d-2}},
\]
and
\[
 1+2\sum_{j\geq1}s^j=\frac{1+s}{1-s}.
\]
Coefficient comparison gives
\[
 \kappa_{d,\ell}=c_{d,\ell}+2\sum_{m=0}^{\ell-1}c_{d,m}.
\]
Substitution and finite reindexing prove \eqref{eq:abstract-branching}.
For \eqref{eq:q-def}, one has
$q_{m+j}=\sbr{A(r_m+j)+1/4}$, which proves
\eqref{eq:proxy-tail-decomp}.  No floor has been approximated or moved.
\end{proof}

The shifts are
\[
 r_m\in
 \begin{cases}
 \N_0+\frac12,&d\text{ odd},\\
 \N,&d\text{ even}.
 \end{cases}
\]
Thus Conjecture~\ref{conj:shifted-tail} is asked only on the two grids that
actually occur.

\section{The signed Weyl primitive}
\label{sec:primitive}

Define the right-continuous branching staircase
\begin{equation}
 S_d(z)=\sum_{\substack{m\geq0\\r_m\leq z}}c_{d,m}
 \label{eq:S-def}
\end{equation}
and its discrepancy primitive
\begin{equation}
 \Delta_d(z)=\int_0^z
 \left(S_d(u)-\frac{u^p}{p!}\right)\dd u.
 \label{eq:Delta-def}
\end{equation}

\begin{proposition}[Signed branching primitive]
\label{prop:primitive-sign}
For every $d\geq3$ and $z\geq0$,
\begin{equation}
 \boxed{\Delta_d(z)\leq0.}
 \label{eq:Delta-sign}
\end{equation}
For $d=3$, equality at positive $z$ occurs precisely at the positive
integers.  For $d\geq4$, the inequality is strict for $z>0$.
\end{proposition}

\begin{proof}
For $0\leq z<r_0=p/2$, $S_d(z)=0$, so
\[
 \Delta_d(z)=-\frac{z^{p+1}}{(p+1)!}\leq0.
\]
Let $M\geq0$ and $r_M=M+p/2$.  The hockey-stick identity gives
\[
 \int_0^{r_M}S_d(u)\dd u
 =\sum_{m=0}^{M-1}(r_M-r_m)c_{d,m}
 =\binom{M+p}{p+1}.
\]
Therefore
\begin{equation}
 (p+1)!\Delta_d(r_M)
 =\prod_{j=0}^{p}(M+j)
 -\left(M+\frac p2\right)^{p+1}\leq0,
 \label{eq:grid-Delta}
\end{equation}
because the arithmetic mean of $M,M+1,\ldots,M+p$ is $M+p/2$.

On the half-open cell $[r_M,r_{M+1})$,
\[
 S_d(z)=\binom{M+p}{p},
 \qquad
 \Delta_d'(z)=\binom{M+p}{p}-\frac{z^p}{p!}.
\]
If the critical point lies in this cell, it is
\[
 z_M^*=\left(\prod_{j=1}^{p}(M+j)\right)^{1/p}.
\]
Direct substitution into the cell formula gives
\begin{equation}
 (p+1)!\Delta_d(z_M^*)
 =p(z_M^*)^p
 \left(z_M^*-M-\frac{p+1}{2}\right)\leq0,
 \label{eq:critical-Delta}
\end{equation}
since $z_M^*$ is the geometric mean of $M+1,\ldots,M+p$ and their
arithmetic mean is $M+(p+1)/2$.  If the critical point is outside the
cell, the cell maximum is at an endpoint, already covered by
\eqref{eq:grid-Delta}.  This proves \eqref{eq:Delta-sign}.

For $p=1$, the cell maximum is $z=M+1$ and has value zero.  For $p\geq2$,
the consecutive factors in both AM--GM applications are unequal, which
gives the stated strictness.
\end{proof}

\section{Exact defect identity and conditional closure}
\label{sec:defect}

Define
\begin{equation}
 D_A(r)=2\int_r^K A(z)\dd z-T_A(r)
 \label{eq:tail-defect}
\end{equation}
and the continuous payment
\begin{equation}
 W_d=\frac2{p!}\int_0^K z^pA(z)\dd z.
 \label{eq:W-def}
\end{equation}

\begin{lemma}[Exact Weyl normalization]
\label{lem:W-normalization}
For $p=d-2$,
\begin{equation}
 \int_0^a z^pG_a(z)\dd z
 =\frac{\Gamma((p+1)/2)}
 {4\sqrt\pi\,(p+2)\Gamma((p+4)/2)}a^{p+2}.
 \label{eq:G-moment}
\end{equation}
Consequently,
\begin{equation}
 \boxed{
 W_d=\frac{1-\rho^d}{2^d\Gamma(d/2+1)^2}K^d
 =L_d|A_{\rho,1}^{(d)}|K^d.}
 \label{eq:W-weyl}
\end{equation}
\end{lemma}

\begin{proof}
Scale $z=ax$ in \eqref{eq:ball-action}, integrate the term containing
$\arccos x$ by parts, and evaluate the resulting beta integrals.  This
gives \eqref{eq:G-moment}.  Multiplying the difference of
\eqref{eq:G-moment} at $a=K$ and $a=\rho K$ by $2/p!$, the duplication
formula for $\Gamma$ reduces the coefficient to
\[
 \frac1{2^{p+2}\Gamma((p+4)/2)^2}
 =\frac1{2^d\Gamma(d/2+1)^2}.
\]
The last equality in \eqref{eq:W-weyl} follows from
$|A_{\rho,1}^{(d)}|=\omega_d(1-\rho^d)$.
\end{proof}

\begin{proposition}[Branching-defect identity]
\label{prop:defect-identity}
Let
\begin{equation}
 \mathcal B_d(A)=2\int_0^K
 (-A'(z))(-\Delta_d(z))\dd z.
 \label{eq:B-def}
\end{equation}
Then $\mathcal B_d(A)\geq0$ and
\begin{equation}
 \boxed{
 W_d-P_d^{<}=\mathcal B_d(A)
 +\sum_{m\geq0}c_{d,m}D_A(r_m).}
 \label{eq:master-defect}
\end{equation}
\end{proposition}

\begin{proof}
The shell action is absolutely continuous, decreasing, and vanishes at
$K$.  Fubini gives
\[
 \sum_{m\geq0}c_{d,m}\,2\int_{r_m}^K A(z)\dd z
 =2\int_0^KS_d(z)A(z)\dd z.
\]
Since $\Delta_d'=S_d-z^p/p!$ almost everywhere and both endpoint terms
vanish, integration by parts gives
\[
 2\int_0^KS_dA
 =W_d+2\int_0^K A\Delta_d'
 =W_d-2\int_0^K A'\Delta_d
 =W_d-\mathcal B_d(A).
\]
Both $A'$ and $\Delta_d$ are nonpositive, so
$\mathcal B_d(A)\geq0$.  Subtract \eqref{eq:proxy-tail-decomp} and use
\eqref{eq:tail-defect}.
\end{proof}

\begin{theorem}[Conditional all-dimensional shell theorem]
\label{thm:conditional}
If Conjecture~\ref{conj:shifted-tail} holds, then
\eqref{eq:polya-target} holds for every $d\geq3$, $0<\rho<1$, and $K\geq0$.
It then also holds for every shell
$A_{r,R}^{(d)}=\{x\in\R^d:r<|x|<R\}$ and every spectral parameter.
\end{theorem}

\begin{proof}
The case $K=0$ is trivial, so assume $K>0$.
Every $r_m$ in \eqref{eq:branching-weights} lies in $\tfrac12\N$, so the
conjecture gives $D_A(r_m)\geq0$.  Equations
\eqref{eq:strict-proxy}, \eqref{eq:master-defect}, and
\eqref{eq:W-weyl} yield
\[
 N_D(A_{\rho,1}^{(d)},K^2)\leq P_d^{<}\leq W_d
 =L_d|A_{\rho,1}^{(d)}|K^d.
\]
For a shell with outer radius $R$, take $\rho=r/R$ and
$K=R\sqrt\Lambda$ and use Dirichlet dilation.  The non-strict counting
form follows by applying the strict estimate at $\Lambda+\varepsilon$ and
letting $\varepsilon\downarrow0$.
\end{proof}

\section{Proved tail regions and the first interface cell}
\label{sec:tail-progress}

The action derivative is
\begin{equation}
 -A'(z)=\frac1\pi
 \begin{cases}
 \displaystyle\arccos\frac zK-\arccos\frac z\mu,&0<z<\mu,\\[2mm]
 \displaystyle\arccos\frac zK,&\mu\leq z<K.
 \end{cases}
 \label{eq:action-slope}
\end{equation}
Thus $A$ is decreasing and concave on $(0,\mu)$, decreasing and convex on
$(\mu,K)$, and $A,A'$ match at $\mu$.

\begin{proposition}[Outer and terminal tails]
\label{prop:easy-tails}
Conjecture~\ref{conj:shifted-tail} holds in either of the following cases:
\begin{enumerate}
 \item $r\geq\mu$;
 \item $G_K(r)\leq3/4$.
\end{enumerate}
\end{proposition}

\begin{proof}
If $r\geq\mu$, then $A=G_K$ on $[r,K]$; the claim is trivial when
$r\geq K$.  Otherwise set $B=\lceil K-r\rceil$ and zero-extend
\[
 g(s)=G_K(r+s)\quad(0\leq s\leq K-r)
\]
to the integer interval $[0,B]$.  Since $G_K(K)=G_K'(K)=0$, this
extension is nonnegative, decreasing, convex, $1/2$-Lipschitz, and
vanishes at $B$.  The convex shifted-floor theorem of \cite{FLPSballs}
therefore bounds its ordinary-floor tail by $2\int_0^B g$.  At every
sample,
\[
 [g(j)+1/4]_<\leq\lfloor g(j)+1/4\rfloor.
\]
The strict tail is consequently bounded by
$2\int_0^Bg=2\int_r^KG_K$.

If $G_K(r)\leq3/4$, monotonicity and $A\leq G_K$ imply
$A(r+j)+1/4\leq1$ for every $j\geq0$.  Every strict bracket is therefore
zero.  Notice that equality at $1$ is harmless precisely because the
bracket is strict.
\end{proof}

We next sharpen the prescribed first milestone.  For the ball action set
\begin{equation}
 T_K(r)=\sbr{G_K(r)+\tfrac14}
 +2\sum_{j\geq1}\sbr{G_K(r+j)+\tfrac14},
 \qquad
 D_K(r)=2\int_r^KG_K(z)\dd z-T_K(r).
 \label{eq:ball-defect}
\end{equation}
The wall-safe translation in the preceding proof gives $D_K(r)\geq0$.

\begin{lemma}[A low-action $1/3$-Lipschitz reserve]
\label{lem:one-layer-reserve}
Let $g$ be strictly decreasing and convex on $[0,b]$, with $g(b)=0$ and
$\operatorname{Lip}(g)\leq1/3$.  If
\[
 \frac34<g(0)<1,
\]
then
\begin{equation}
 2\int_0^b g(z)\dd z-
 \left(\sbr{g(0)+\tfrac14}
 +2\sum_{j\geq1}\sbr{g(j)+\tfrac14}\right)>\frac{11}{16}.
 \label{eq:one-layer-reserve}
\end{equation}
The function is understood to be zero-extended after $b$.
\end{lemma}

\begin{proof}
Let $h=g(0)$ and let $x\in(0,b)$ be defined by $g(x)=3/4$.  Since
$g(j)<1$, a sample contributes exactly one if and only if $j<x$.
Therefore the bracketed count is
\[
 2\lceil x\rceil-1.
\]
The Lipschitz bound and $g(b)=0$ give
\[
 g(z)\geq\max\{h-z/3,0\},
 \qquad
 2\int_0^b g(z)\dd z\geq3h^2>\frac{27}{16}.
\]
If $x\leq1$, the count is one and \eqref{eq:one-layer-reserve} follows.

Suppose $x>1$ and put $a=(h-3/4)/x$.  Convexity implies that, for
$z\geq x$,
\[
 g(z)\geq\max\left\{\frac34-a(z-x),0\right\}.
\]
Hence
\[
 \int_0^b g(z)\dd z
 \geq\frac34x+\frac{(3/4)^2}{2a}
 =x\left(\frac34+\frac9{32(h-3/4)}\right)
 >\frac{15}{8}x.
\]
Since $2\lceil x\rceil-1<2x+1$, the defect is larger than
\[
 \frac{15}{4}x-(2x+1)=\frac74x-1>\frac34>\frac{11}{16}.
\]
\end{proof}

\begin{proposition}[Exact one-interface reduction]
\label{prop:one-interface}
Assume
\[
 0<\mu-r<1,
 \qquad r\in\tfrac12\N.
\]
Define
\begin{equation}
 n_0=\sbr{G_K(r)+\tfrac14}
 -\sbr{G_K(r)-G_\mu(r)+\tfrac14}\in\N_0
 \label{eq:n0-def}
\end{equation}
and
\[
 \delta_\mu(r)=\int_r^\mu G_\mu(z)\dd z.
\]
Then
\begin{equation}
 \boxed{D_A(r)=D_K(r)+n_0-2\delta_\mu(r).}
 \label{eq:one-interface-identity}
\end{equation}
Moreover,
\begin{equation}
 2\delta_\mu(r)
 <\frac{4(\mu-r)^{5/2}}{15\sqrt\mu}
 <\frac{4\sqrt2}{15}<1.
 \label{eq:cap-bound}
\end{equation}
Consequently, the shifted-tail inequality holds if any of the following is
true:
\begin{enumerate}
 \item $n_0\geq1$;
 \item $G_K(\max\{r,K/2\})\geq1$;
 \item $r\geq K/2$.
\end{enumerate}
The only remaining one-interface configurations satisfy
\begin{equation}
 \boxed{
 r<\frac K2,\qquad
 K<\frac1{\omega_0},\qquad
 n_0=0,\qquad
 \omega_0:=G_1(1/2)=\frac{\sqrt3}{2\pi}-\frac16.}
 \label{eq:one-interface-residual}
\end{equation}
In particular,
\[
 r\in\left\{\frac12,1,\frac32,2,\frac52,3,\frac72,4,\frac92\right\}.
\]
\end{proposition}

\begin{proof}
Since $r+j>\mu$ for every $j\geq1$, all samples after the first agree
with the ball samples.  Thus $T_A(r)=T_K(r)-n_0$.  Subtracting
$2\int_r^\mu G_\mu$ from the ball integral proves
\eqref{eq:one-interface-identity}.

The turning estimate
\[
 G_\mu(z)<\frac{(\mu-z)^{3/2}}{3\sqrt\mu}
 \qquad(0<z<\mu)
\]
integrates to \eqref{eq:cap-bound}, using $\mu>r\geq1/2$ and
$\mu-r<1$.  If $n_0\geq1$, then $D_K(r)\geq0$ and
\eqref{eq:one-interface-identity} is positive.

For $r<K$, use the same integer endpoint
$B=\lceil K-r\rceil$ and zero-extended function
$g(s)=G_K(r+s)$.  Put
\[
 s_*=\max\{0,K/2-r\}.
\]
Then $g$ is $1/3$-Lipschitz on $[s_*,B]$ and
$g(s_*)=G_K(\max\{r,K/2\})$.  The refined ordinary-floor theorem of
\cite[Theorem~3.4]{FLPSannuli}, followed by
$[\,\cdot\,]_<\leq\lfloor\,\cdot\,\rfloor$, gives
\begin{equation}
 D_K(r)\geq\frac12
 \left\lfloor G_K(\max\{r,K/2\})\right\rfloor.
 \label{eq:refined-ball-reserve}
\end{equation}
For $r\geq K$ the same inequality is trivial.  Hence the second condition
and \eqref{eq:cap-bound} prove the claim.

If $r\geq K/2$ and $G_K(r)<1$, the translated ball action is globally
$1/3$-Lipschitz.  Its tail is either zero or
Lemma~\ref{lem:one-layer-reserve} applies.  In the latter case
$D_K(r)>11/16>4\sqrt2/15$, so \eqref{eq:one-interface-identity} is again
positive.  This proves the third condition.

If none of the three conditions holds, then $r<K/2$ and
$G_K(K/2)=\omega_0K<1$, giving \eqref{eq:one-interface-residual}.  The
elementary rational bounds
\[
 \frac{265}{153}<\sqrt3<\frac{26}{15},
 \qquad
 \frac{333}{106}<\pi<\frac{355}{113}
\]
(the square-root bounds follow by squaring) give
$1/10<\omega_0<1/9$.  Consequently
\[
 \frac92<\frac1{2\omega_0}<5.
\]
Since $r<K/2<1/(2\omega_0)$ and $r\in\tfrac12\N$, only the nine displayed
shifts remain.
\end{proof}

\begin{remark}[The residual first-interface family]
Proposition~\ref{prop:one-interface} replaces the only unresolved part of
the first-interface problem by a bounded nine-shift family:
$2r<K<1/\omega_0$, $r<\mu<\min\{r+1,K\}$, and $n_0=0$.  The inequality to
be checked there is exactly
\[
 D_K(r)\geq2\delta_\mu(r).
\]
Proposition~\ref{prop:finite-wall} and
Lemma~\ref{lem:exact-wall-certificate} below close this family.
\end{remark}

\begin{proposition}[Finite ball-wall reduction]
\label{prop:finite-wall}
In the residual family of Proposition~\ref{prop:one-interface}, put
\[
 B=G_K(r)+\frac14,\qquad
 q=[B]_<,\qquad
 e=B-q\in(0,1],
\]
and set
\[
 E_r(K)=D_K(r)-e.
\]
Then
\begin{equation}
 2\delta_\mu(r)<e.
 \label{eq:delta-below-remainder}
\end{equation}
Consequently, the residual one-interface inequality follows from
$E_r(K)\geq0$.

On every open $K$-interval containing no ball-sample wall,
$E_r$ is strictly increasing.  At a wall of the first sample
$G_K(r)+1/4$, the two unit drops in $D_K(r)$ and $e$ cancel, so $E_r$ is
continuous.  At a wall of a later sample $G_K(r+j)+1/4$, $j\geq1$,
$E_r$ drops by two for each crossed level.  It is therefore enough to
check the right-hand limits at the finitely many later-sample walls,
together with the first positive-count wall and any lower endpoint
$K=2r$ at which the count is already positive.

More explicitly, all walls that can occur have the form
\begin{equation}
 G_K(r+j)=n-\frac14,
 \qquad
 j\geq0,\qquad n\in\{1,2,3\},\qquad r+j<10.
 \label{eq:finite-wall-list}
\end{equation}
\end{proposition}

\begin{proof}
Let $h=G_\mu(r)$.  If $q=0$, strict increase of $G_a(r)$ in the radius
$a$ gives $h<G_K(r)<B=e$.  If $q\geq1$, write $B=q+e$ and use the strict
bracket convention; then $[B-h]_<=[B]_<$ holds exactly when $h<e$.
Thus, in both cases,
\[
 n_0=0
 \quad\Longleftrightarrow\quad
 [B-h]_<=[B]_<
 \quad\Longleftrightarrow\quad
 h<e.
\]
When $q\geq1$, equality $h=e$ lands on the integral wall $q$ and causes a
strict-count drop.  The function $G_\mu$ is strictly convex on
$[r,\mu]$, with endpoint values $h$ and $0$.  It lies strictly below its
endpoint chord, and therefore
\[
 2\delta_\mu(r)<h(\mu-r)<e,
\]
which proves \eqref{eq:delta-below-remainder}.  If $T_K(r)=0$, then
$T_A(r)=0$ and the desired inequality is immediate.  Otherwise it is
enough, by \eqref{eq:one-interface-identity}, to prove
$D_K(r)\geq e$.

Away from sample walls the integer-valued tail and $q$ are constant.
Writing $x=r/K$, differentiation of the ball action and its tail integral
gives
\begin{equation}
 \pi E_r'(K)
 =
 K\left(\arccos x-x\sqrt{1-x^2}\right)-\sqrt{1-x^2},
 \label{eq:E-first-derivative}
\end{equation}
and
\begin{equation}
 \pi E_r''(K)
 =
 \arccos x+x\sqrt{1-x^2}
 -\frac{x^2}{K\sqrt{1-x^2}}.
 \label{eq:E-second-derivative}
\end{equation}
Here $x<1/2$ and $K>1$.  Hence
\[
 \pi E_r''(K)>
 \frac\pi3-\frac1{2\sqrt3}>0.
\]
For $r\geq1$, evaluation at $K=2r$ gives
\[
 \pi E_r'(2r)
 =
 \frac{2\pi r}{3}-\frac{\sqrt3}{2}(r+1)>0.
\]
For $r=1/2$, positivity of the tail count implies
$G_K(r)>3/4$.  Since $G_K$ is $1/2$-Lipschitz and vanishes at $K$, this
forces $K>2$.  Formula~\eqref{eq:E-first-derivative} is already positive
at $(r,K)=(1/2,2)$, since
\[
 \pi E_{1/2}'(2)
 =
 2\arccos\frac14-\frac{3\sqrt{15}}8
 >
 \frac{2\pi}{3}-\frac32>0.
\]
Equation~\eqref{eq:E-second-derivative} then keeps it positive.  Thus $E_r$ is
strictly increasing between walls.

The stated jump rules follow directly from the weights $1,2,2,\ldots$
in $T_K(r)$.  At a first-sample wall, $e$ equals $1$ and its right-hand
limit is $0$, exactly matching the unit right-hand drop of $D_K$.
Simultaneous later-sample crossings are combined at the same wall
abscissa.  Finally, $K<1/\omega_0<10$ and
$G_K(z)<K/\pi<10/3$.  Thus only $n=1,2,3$ and samples $r+j<10$ can occur,
and
\[
 \partial_KG_K(z)=\frac1\pi\sqrt{1-\frac{z^2}{K^2}}>0
 \qquad(K>z)
\]
shows that every pair $(j,n)$ produces at most one wall.  This proves the
finite reduction.
\end{proof}

\begin{lemma}[Exact residual wall certificate]
\label{lem:exact-wall-certificate}
In the residual family of Proposition~\ref{prop:finite-wall}, if
$T_K(r)>0$, then
\begin{equation}
 E_r(K)>0.
 \label{eq:exact-wall-positive}
\end{equation}
More precisely, $E_r(K)>3/20$ for $r\leq4$, while for $r=9/2$ one has
\[
 E_{9/2}(K)>\frac{1253}{555}.
\]
\end{lemma}

\begin{proof}
Put
\[
 H_r(K)=2\int_r^K G_K(z)\dd z-G_K(r)-\frac14,
 \qquad
 M_r(K)=\sum_{j\geq1}\sbr{G_K(r+j)+\frac14}.
\]
The first-sample bracket cancels from the definition of $E_r$, including
at its strict walls, and hence
\begin{equation}
 E_r(K)=H_r(K)-2M_r(K).
 \label{eq:E-as-H-minus-M}
\end{equation}
For $x=r/K$, direct integration gives
\begin{equation}
 H_r(K)=\frac A2-
 \frac{A\arcsin x+B\sqrt{1-x^2}}{\pi}-\frac14,
 \quad
 A=\frac{K^2}{2}+r^2+r,
 \quad
 B=K\left(\frac{3r}{2}+1\right).
 \label{eq:H-closed-form}
\end{equation}

We record rational envelopes sufficient to check every wall.  For
$0\leq x\leq13/20$, let
\[
 a_m=\frac{\binom{2m}{m}}{4^m(2m+1)},
 \qquad
 b_m=\frac{\binom{2m}{m}}{4^m(2m-1)},
\]
and set
\[
 \begin{aligned}
 L(x)&=\sum_{m=0}^{6}a_mx^{2m+1},&
 U(x)&=L(x)+\frac{x^{15}}{1-x^2},\\
 V(x)&=1-\sum_{m=1}^{6}b_mx^{2m},&
 W(x)&=V(x)-\frac{x^{14}}{1-x^2}.
 \end{aligned}
\]
The power series, whose omitted positive coefficients are smaller than
one, give
\[
 L\leq\arcsin x\leq U,
 \qquad
 W\leq\sqrt{1-x^2}\leq V.
\]
Together with $333/106<\pi<355/113$, they imply
\begin{equation}
 \begin{aligned}
 \underline G(K,z)
 &=K\left[\frac{113}{355}(W+xL)-\frac x2\right]
 \leq G_K(z)\\
 &\leq
 K\left[\frac{106}{333}(V+xU)-\frac x2\right]
 =\overline G(K,z),
 \end{aligned}
 \label{eq:rational-G-envelope}
\end{equation}
where now $x=z/K$, and
\begin{equation}
 \frac A2-\frac{106}{333}(AU+BV)-\frac14
 \leq H_r(K)\leq
 \frac A2-\frac{113}{355}(AL+BW)-\frac14.
 \label{eq:rational-H-envelope}
\end{equation}
Thus every comparison below is a rational-number comparison.

Let $k_n(z)$ denote the unique solution of
$G_K(z)=n-1/4$.  Substitution in
\eqref{eq:rational-G-envelope} gives the following pairwise disjoint
brackets:
\[
\begin{array}{c|c@{\qquad}c|c}
\text{wall}&\text{bracket}&\text{wall}&\text{bracket}\\ \hline
k_1(3/2)&(22/5,9/2)&k_1(2)&(5,26/5)\\
k_1(5/2)&(57/10,29/5)&k_1(3)&(63/10,32/5)\\
k_1(7/2)&(69/10,7)&k_1(4)&(15/2,38/5)\\
k_2(3/2)&(77/10,39/5)&k_1(9/2)&(81/10,41/5)\\
k_2(2)&(83/10,17/2)&k_1(5)&(87/10,44/5)\\
k_2(5/2)&(9,227/25)&&
\end{array}
\tag{\ref{eq:rational-G-envelope}a}
\]
These are all later-sample walls.  Indeed, the same rational bounds give
\[
 \frac{227}{25}<\frac1{\omega_0}<\frac{459}{50}
\]
and, at $K=459/50$,
\[
 \overline G(K,11/2)<\frac{18}{25}<\frac34,
 \quad
 \overline G(K,3)<\frac85<\frac74,
 \quad
 \overline G(K,3/2)<\frac94<\frac{11}{4}.
\]
Monotonicity in $K$ and $z$ excludes, respectively, every omitted
level-$3/4$, level-$7/4$, and level-$11/4$ or higher wall.

At a right-hand wall limit, the value of $M_r$ is the event number in
the following exact order:
\[
\begin{array}{c|l}
r&(z,n;M_r)\\ \hline
\frac12&(\frac32,1;1),(\frac52,1;2),(\frac72,1;3),
(\frac32,2;4),(\frac92,1;5),(\frac52,2;6)\\
1&(2,1;1),(3,1;2),(4,1;3),(2,2;4),(5,1;5)\\
\frac32&(\frac52,1;1),(\frac72,1;2),(\frac92,1;3),(\frac52,2;4)\\
2&(3,1;1),(4,1;2),(5,1;3)\\
\frac52&(\frac72,1;1),(\frac92,1;2)\\
3&(4,1;1),(5,1;2)\\
\frac72&(\frac92,1;1)\\
4&(5,1;1)\\
\frac92&\varnothing.
\end{array}
\]
Equations~\eqref{eq:H-closed-form}--\eqref{eq:rational-H-envelope}
then give the strict rational bounds
\[
\begin{array}{c|c|c|c|c}
r&\text{first-wall lower}&\text{first-wall upper}
 &\text{activation or }K=2r&\text{other-wall lower}\\ \hline
\frac12&>3/20&<8/25&>11/25&>13/20\\
1&>17/100&<1/2&>14/25&>7/10\\
\frac32&>9/25&<53/100&>69/100&>37/50\\
2&>2/5&<29/50&>37/50&>39/50\\
\frac52&>23/50&<63/100&>22/25&>81/100\\
3&>13/25&<69/100&>47/50&>17/20\\
\frac72&>29/50&<19/25&>109/100&--\\
4&>33/50&<83/100&>163/100&--.
\end{array}
\]
For each $r\leq4$, the two first-wall columns concern the right-hand value
at $G_K(r+1)=3/4$.  Its displayed upper bound is smaller than both
competing lower bounds, so this is the controlling wall.
Formula~\eqref{eq:E-first-derivative} shows that $H_r$, and hence $E_r$
between its later-sample walls, is increasing throughout the
positive-count residual region.  For $r\leq3$ the tail is zero before its
first-sample activation; for $r=7/2,4$ the displayed $K=2r$ value owns the
lower endpoint.  Finally, for $r=9/2$ there is no later wall and
\[
 E_{9/2}(9)=\frac{43}{2}-\frac{279\sqrt3}{8\pi}
 >\frac{1253}{555}.
\]
This proves \eqref{eq:exact-wall-positive}.  A dependency-free replay of
all $628$ rational comparisons is supplied in
\path{scripts/general_d_finite_wall_fraction_verify.py}; it is an audit of,
not a numerical premise for, the displayed estimates.
\end{proof}

\begin{theorem}[One-interface shifted tails]
\label{thm:one-interface-complete}
Conjecture~\ref{conj:shifted-tail} holds whenever
\[
 0<\mu-r<1.
\]
It also holds at the endpoint $r=\mu$ by
Proposition~\ref{prop:easy-tails}.
\end{theorem}

\begin{proof}
The three regions in Proposition~\ref{prop:one-interface} are already
proved there.  In its residual family, if $T_K(r)=0$, then $T_A(r)=0$ and
the claim is immediate.  Otherwise Lemma~\ref{lem:exact-wall-certificate}
and Proposition~\ref{prop:finite-wall} give
\[
 D_K(r)=E_r(K)+e>e>2\delta_\mu(r).
\]
The exact identity \eqref{eq:one-interface-identity} now gives
$D_A(r)>0$.
\end{proof}

\begin{proposition}[Fixed-ratio high-energy tails]
\label{prop:fixed-rho-high-energy}
For $0<\rho<1$, put
\[
 c_\rho=\frac{\arccos\rho}{\pi},
 \qquad d_\rho=1-2c_\rho=\frac{2\arcsin\rho}{\pi},
\]
\[
 \eta_\rho=G_1\left(\max\left\{\rho,\frac12\right\}\right),
 \qquad
 a_\rho=\frac{2\pi\rho}{1-\rho},
 \qquad C_*=\frac12+\frac{8\sqrt2}{15},
\]
and
\begin{equation}
 K_0(\rho)=
 \frac{a_\rho+2\eta_\rho C_*
 +\sqrt{a_\rho^2+4a_\rho\eta_\rho C_*+\eta_\rho^2}}
 {2\eta_\rho^2}.
 \label{eq:fixed-rho-threshold}
\end{equation}
If $K\geq K_0(\rho)$, then
Conjecture~\ref{conj:shifted-tail} holds for every
$r\in\tfrac12\N$.
\end{proposition}

\begin{proof}
Only $r<\mu$ requires proof.  Set
\[
 n=\lfloor\mu-r\rfloor,
 \qquad q=r+n,
 \qquad B=\lceil K-r\rceil.
\]
Thus $q\leq\mu<q+1$.  For integers $a<b$, write
\[
 \mathbf T(f;a,b)=\frac12\lfloor f(a)\rfloor
 +\sum_{j=a+1}^{b-1}\lfloor f(j)\rfloor
 +\frac12\lfloor f(b)\rfloor.
\]
The strict brackets in $T_A$ are bounded above by ordinary floors.
Additivity at the split point and $A\leq G_K$ therefore give, when
$n\geq1$,
\begin{equation}
 \frac{T_A(r)}2\leq
 \mathbf T\left(A(r+\cdot)+\frac14;0,n\right)
 +\mathbf T\left(G_K(r+\cdot)+\frac14;n,B\right).
 \label{eq:high-energy-split}
\end{equation}
For $n=0$, the first term is absent and the same inequality follows by
direct pointwise majorization.

Assume first $n\geq1$ and put
\[
 h(t)=A(r+t)+\frac14,
 \qquad m_j=\lfloor h(j)\rfloor.
\]
If the nonincreasing sequence $(m_j)_{j=0}^n$ leaves its first shelf, let
$p$ be its last index on that shelf; otherwise put $p=n$.  On $[0,n]$,
$h$ is decreasing, concave, and $c_\rho$-Lipschitz.  The quantitative
concave trapezoidal-floor theorem
\cite[Theorem~1.4]{FLPSannuli}, with the basic concavity bound in the
no-drop case, yields
\begin{equation}
 \mathbf T(h;0,n)
 \leq\int_0^n A(r+t)\dd t+\frac n4
 -\frac{1-c_\rho}{2}(n-p).
 \label{eq:concave-head-gain}
\end{equation}
For $n=0$, take $p=0$ and read this as an absent identity.

Let $\sigma=-A'$.  On $0<z<\mu$,
\[
 \sigma'(z)=\frac1{\pi\sqrt{\mu^2-z^2}}
 -\frac1{\pi\sqrt{K^2-z^2}}
 \geq\kappa:=\frac{1-\rho}{\pi\rho K}.
\]
Indeed, after subtracting the value at $z=0$, the remaining bracket is
decreasing in its radius parameter.  Since $\sigma(0)=0$,
\[
 A(x)-A(y)=\int_x^y\sigma(z)\dd z
 \geq\frac\kappa2(y^2-x^2)
 \geq\frac\kappa2(y-x)^2.
\]
The samples at $r$ and $r+p$ lie on the same floor shelf, so their action
drop is smaller than one.  Keeping the $r$-term in the preceding estimate
and using $r\geq1/2$ gives
\begin{equation}
 p<\sqrt{r^2+a_\rho K}-r
 \leq\sqrt{a_\rho K+\frac14}-\frac12.
 \label{eq:first-shelf-length}
\end{equation}

For the convex suffix, set
\[
 s=\max\left\{q,\frac K2\right\},
 \qquad t_*=s-r.
\]
Then $t_*\in[n,B]$, the translated ball action is
$1/2$-Lipschitz on $[n,B]$ and $1/3$-Lipschitz on $[t_*,B]$, and
\[
 G_K(s)\geq K\eta_\rho.
\]
The quantitative convex theorem \cite[Theorem~3.4]{FLPSannuli} gives
\begin{equation}
 \mathbf T\left(G_K(r+\cdot)+\frac14;n,B\right)
 \leq\int_n^B G_K(r+t)\dd t
 -\frac14\lfloor K\eta_\rho\rfloor.
 \label{eq:convex-suffix-gain}
\end{equation}

The only mismatch between the two integrals on the right of
\eqref{eq:high-energy-split} and the shell integral is
\[
 \delta=\int_q^\mu G_\mu(z)\dd z.
\]
The turning bound from \cite[Lemma~6.2]{FLPSannuli} and
$0\leq\mu-q<1$, $\mu>1/2$, give
\begin{equation}
 0\leq\delta\leq
 \frac{2(\mu-q)^{5/2}}{15\sqrt\mu}
 <\frac{2\sqrt2}{15}.
 \label{eq:interface-mismatch-high-energy}
\end{equation}
Combining \eqref{eq:high-energy-split},
\eqref{eq:concave-head-gain}, and
\eqref{eq:convex-suffix-gain} gives
\begin{equation}
 \frac{T_A(r)}2\leq\int_r^K A(z)\dd z+\delta-\frac M4,
 \quad
 M=\lfloor K\eta_\rho\rfloor+d_\rho(n-p)-p.
 \label{eq:high-energy-M}
\end{equation}
By \eqref{eq:first-shelf-length},
\[
 M>K\eta_\rho-\sqrt{a_\rho K+\frac14}-\frac12.
\]
The definition \eqref{eq:fixed-rho-threshold} is exactly the positive-root
condition
\[
 K\eta_\rho-\sqrt{a_\rho K+\frac14}\geq C_*.
\]
Thus $M>8\sqrt2/15>4\delta$, and
\eqref{eq:high-energy-M} proves the result.  All inequalities use
ordinary floors on the majorizing side, so strict action walls require no
limiting argument.
\end{proof}

The fixed-ratio threshold is deliberately uniform in the shift, but it is
not sharp as $\rho\uparrow1$.  Retaining the location of the first shelf
improves the thin-shell scale by two powers of $1-\rho$.

\begin{lemma}[Optimized thin first-shelf loss]
\label{lem:thin-first-shelf-loss}
Let
\[
 \rho=1-\varepsilon,
 \qquad 0<\varepsilon\leq\frac1{100},
 \qquad K\varepsilon^2\geq\frac{125}{8}.
\]
Use $n,p$ from the proof of
Proposition~\ref{prop:fixed-rho-high-energy}, put $m=n-p$, and set
\[
 d=d_\rho,
 \qquad \mathcal R=p-dm.
\]
If $\mathcal R>0$, then
\begin{equation}
 \mathcal R<\frac{361}{80\sqrt\varepsilon}.
 \label{eq:thin-shelf-loss}
\end{equation}
\end{lemma}

\begin{proof}
This is a scale-free consequence of the exact local slope and the
same-floor condition.  Put
\[
 y=\sqrt\varepsilon,
 \quad \varkappa=K\varepsilon^2,
 \quad P=py,
 \quad \mathfrak r=\mathcal R y,
 \quad S=(p+m)y.
\]
Elementary trigonometric bounds give
\begin{equation}
 d>\frac9{10},
 \qquad
 S=P+\frac{P-\mathfrak r}{d}<\frac{19}{9}P.
 \label{eq:thin-scaled-geometry}
\end{equation}
Write
\[
 U=\mu-r,\qquad t=\frac r\mu,\qquad
 w=\frac U\mu=1-t,
\]
and define
\[
 \mathcal Q=1+\frac{y^2}{\varkappa}
 +\frac{yS}{\varkappa},
 \qquad
 \widehat q=\frac{y^2}{\rho}(\mathcal Q-1).
\]
The unconditional first-shelf bound
$p^2<2\pi\rho K/\varepsilon$ and
\eqref{eq:thin-scaled-geometry} give
\[
 w<\frac{1+p+m}{\rho K}
 =\frac{y^4+y^3S}{\varkappa\rho}
 =\widehat q
 <\frac{y^4}{\varkappa\rho}
 +\frac{19}{9}y\sqrt{\frac{2\pi}{\varkappa\rho}}.
\]
Here the first term is at most $8/1237500<1/1000$, while the
square of the factor following $19/9$ is at most
$352/86625<1/225$.  Hence
\begin{equation}
 w<\widehat q<\frac1{1000}+\frac{19}{135}<\frac17,
 \qquad t>1-\widehat q>\frac67.
 \label{eq:thin-qhat-localization}
\end{equation}
Using
\[
 -A'(z)=\frac1\pi
 \left(\arccos\frac zK-\arccos\frac z\mu\right)
\]
in the same-floor inequality $A(r)-A(r+p)<1$ gives, after scaling,
\begin{equation}
 P^2<\mathcal H(P,\mathfrak r)
 :=\frac{2\pi^2\mathcal Q}{(1-\widehat q)^2}.
 \label{eq:thin-synthetic-bound}
\end{equation}
On the preceding domain, $\mathcal H$ decreases in $\mathfrak r$ and
\[
 \partial_P\mathcal H
 <\frac{673816}{1366875}<\frac12.
\]
Let $B_*=361/80$.  If $\mathfrak r\geq B_*$, then $P\geq\mathfrak r$
and $P^2-\mathcal H(P,B_*)$ is increasing for $P\geq B_*$.  Indeed,
every intermediate $P'\in[B_*,P]$ retains the unconditional shelf
ceiling and satisfies $S(P',B_*)<19P'/9$, so
\eqref{eq:thin-qhat-localization} and the displayed derivative bound hold
throughout the comparison interval.  Its minimum over $y\leq1/10$ and
$\varkappa\geq125/8$ is at
$P=B_*$, $y=1/10$, and $\varkappa=125/8$.  There
\[
 \mathcal Q=\frac{12869}{12500},
 \qquad
 \widehat q=\frac{41}{137500},
\]
and, using $\pi<22/7$,
\[
 B_*^2-2\left(\frac{22}{7}\right)^2
 \frac{12869/12500}{(1-41/137500)^2}
 =\frac{72581986185449}{5925464687161600}>0.
\]
This contradicts \eqref{eq:thin-synthetic-bound} and proves
\eqref{eq:thin-shelf-loss}.  All displayed comparisons are rational once
the elementary bound for $\pi$ is inserted; the dependency-free script
\path{computations/round9_verify_thin_plateau_constants.py} replays the
finite arithmetic.
\end{proof}

\begin{proposition}[Thin high-energy shifted tails]
\label{prop:thin-high-energy}
Let $\rho=1-\varepsilon$ with
$0<\varepsilon\leq1/100$.  If
\begin{equation}
 K\geq\frac{125}{8\varepsilon^2},
 \label{eq:thin-high-energy-threshold}
\end{equation}
then Conjecture~\ref{conj:shifted-tail} holds for every
$r\in\tfrac12\N$.
\end{proposition}

\begin{proof}
The outer tails are covered by Proposition~\ref{prop:easy-tails}.  For an
inner tail, retain the notation of
Proposition~\ref{prop:fixed-rho-high-energy}.  Before its final estimate,
\eqref{eq:high-energy-M} gives
\[
 \frac{T_A(r)}2
 \leq\int_r^KA(z)\dd z+\delta-\frac M4,
 \qquad
 M=\lfloor KG_1(\rho)\rfloor-\mathcal R,
\]
with $4\delta<8\sqrt2/15$.  If $\mathcal R\leq0$, the convex gain already
gives $M>1>4\delta$.  If $\mathcal R>0$, Lemma
~\ref{lem:thin-first-shelf-loss},
\[
 G_1(1-\varepsilon)
 \geq\frac{2\sqrt2}{3\pi}\varepsilon^{3/2},
 \qquad
 \sqrt2>\frac{140}{99},
 \qquad
 \sqrt2<\frac{99}{70},
 \qquad
 \pi<\frac{1571}{500},
\]
and $\lfloor x\rfloor>x-1$ give
\[
 M>
 \frac{2829397}{3732696}
 >\frac{132}{175}
 >\frac{8\sqrt2}{15}>4\delta.
\]
Thus \eqref{eq:high-energy-M} proves the claim.  The proof uses only a
positive shift and a unit-spaced sampling grid; it is unchanged on the
integer and half-integer grids, and every threshold equality is owned by
the strict-floor majorant.
\end{proof}

\begin{corollary}[A uniform small-hole sector]
\label{cor:small-hole-sector}
If
\[
 0<\rho<\omega_0,
 \qquad
 K(\omega_0-\rho)\geq\frac12+\frac{8\sqrt2}{15},
\]
then Conjecture~\ref{conj:shifted-tail} holds for every
$r\in\tfrac12\N$.
\end{corollary}

\begin{proof}
For $r<\mu$, use the split \eqref{eq:high-energy-split}, basic concavity
on the head, and \eqref{eq:convex-suffix-gain} with $s=K/2$.  This gives
\[
 \frac{T_A(r)}2\leq\int_r^K A(z)\dd z+\delta
 -\frac{\lfloor\omega_0K\rfloor-n}{4}.
\]
Since $r\geq1/2$,
\[
 \lfloor\omega_0K\rfloor-n
 >K(\omega_0-\rho)-\frac12
 \geq\frac{8\sqrt2}{15}>4\delta.
\]
The outer tails were proved in Proposition~\ref{prop:easy-tails}.
\end{proof}

\section{Cell recurrence and the corrected global target}

For $x_j=r+j$, set
\[
 Q_j=\sbr{A(x_j)+\frac14},
 \qquad
 D_j=2\int_{x_j}^KA(z)\dd z-
 \left(Q_j+2\sum_{k>j}Q_k\right).
\]
Then
\begin{equation}
 D_j-D_{j+1}
 =2\int_{x_j}^{x_j+1}A(z)\dd z-Q_j-Q_{j+1}.
 \label{eq:cell-recurrence}
\end{equation}
If
\[
 e_j=A(x_j)+\frac14-Q_j\in(0,1]
\]
and
\[
 C_j=2\int_{x_j}^{x_j+1}A(z)\dd z-A(x_j)-A(x_j+1),
\]
then
\begin{equation}
 D_j-D_{j+1}=C_j-\frac12+e_j+e_{j+1}.
 \label{eq:cell-remainder}
\end{equation}
On a fully inner cell,
\[
 C_j=-\int_0^1u(1-u)A''(x_j+u)\dd u\geq0.
\]

Individual increments in \eqref{eq:cell-remainder} can nevertheless be
negative, and a negative increment need not be paid by the next floor-drop
cell.  Thus the viable block formulation must retain the terminal convex
reserve.  If $J$ is the first index with $x_J\geq\mu$, the exact target is
\begin{equation}
 \boxed{
 D_0=\sum_{j=0}^{J-1}(D_j-D_{j+1})+D_J\geq0,}
 \label{eq:corrected-block-target}
\end{equation}
where $D_J\geq0$ is the translated convex ball-tail defect.  A proof of
\eqref{eq:corrected-block-target} needs a quantitative, not merely
nonnegative, lower bound for $D_J$ or an equivalent matched shelf reserve.

The many-cell floor dynamics can in fact be compressed to one initial
shelf and one terminal one-interface tail.  This is the form used below.

\begin{proposition}[Wall-safe first-shelf reduction]
\label{prop:first-shelf-reduction}
Suppose $r<\mu$ and put
\[
 n=\lfloor\mu-r\rfloor,
 \qquad q=r+n,
 \qquad
 d_\rho=\frac{2\arcsin\rho}{\pi}.
\]
For $0\leq j\leq n$, let
\[
 F_j=\left\lfloor A(r+j)+\frac14\right\rfloor.
\]
If this sequence leaves its first shelf, let $p$ be the last index on
that shelf; otherwise put $p=n$.  Define
\begin{equation}
 R_p=2\int_r^{r+p}A(z)\dd z-2pF_0.
 \label{eq:first-shelf-reserve}
\end{equation}
Then, at ordinary and strict floor walls alike,
\begin{equation}
 \boxed{
 D_A(r)\geq D_A(q)+R_p+\frac{d_\rho}{2}(n-p).}
 \label{eq:first-shelf-reduction}
\end{equation}
Moreover, if
\[
 \varepsilon_j=A(r+j)+\frac14-F_j\in[0,1),
 \qquad
 \sigma=-A',
\]
then
\begin{align}
 R_p
 &=\mathcal C_p+p\left(\varepsilon_0+\varepsilon_p-\frac12\right),
 \label{eq:first-shelf-exact}\\
 \mathcal C_p
 &:=-\int_0^p u(p-u)A''(r+u)\dd u\geq0,
 \label{eq:first-shelf-curvature}
\end{align}
and
\begin{equation}
 R_p\geq-\frac p2+2\int_0^p u\sigma(r+u)\dd u.
 \label{eq:first-shelf-lower}
\end{equation}
In particular, with
\[
 \kappa=\frac{1-\rho}{\pi\rho K},
 \qquad
 a_\rho=\frac{2\pi\rho}{1-\rho},
\]
one has
\begin{equation}
 p<\sqrt{r^2+a_\rho K}-r.
 \label{eq:first-shelf-sharp-cap}
\end{equation}
\end{proposition}

\begin{proof}
Put $h(t)=A(r+t)+1/4$.  At half-tail scale, the strict trapezoidal
sum
\[
 \mathbf T_{<}(h;a,b)
 :=\frac12\sbr{h(a)}+
 \sum_{j=a+1}^{b-1}\sbr{h(j)}+\frac12\sbr{h(b)}
\]
on $[0,p]$ is at most $pF_0$.  If $p<n$, the ordinary floor of $h$
drops immediately after the left endpoint of $[p,n]$.  The quantitative
concave trapezoidal-floor theorem \cite[Theorem~1.4]{FLPSannuli} gives
\[
 \mathbf T_{<}(h;p,n)
 \leq\int_p^n A(r+t)\dd t-\frac{d_\rho}{4}(n-p).
\]
If $p=n$, this block is absent.  Additivity at $p$ and $n$ leaves exactly
one half of the strict tail beginning at $q$, and hence
\[
 \frac{T_A(r)}2
 \leq pF_0+\int_{r+p}^qA(z)\dd z
 -\frac{d_\rho}{4}(n-p)+\frac{T_A(q)}2.
\]
Subtracting this from $\int_r^KA$ proves
\eqref{eq:first-shelf-reduction}.  At an ordinary wall the strict bracket
is smaller than the ordinary floor used in the majorant, so no limiting
argument is needed.

Since $F_0=F_p$, the trapezoidal identity gives
\eqref{eq:first-shelf-exact}--\eqref{eq:first-shelf-curvature}.  Also
$F_0\leq A(r+p)+1/4$, and therefore
\[
 R_p\geq-\frac p2
 +2\int_0^p\bigl(A(r+t)-A(r+p)\bigr)\dd t,
\]
which is \eqref{eq:first-shelf-lower} after reversing the order of
integration.  Finally, on the inner branch
\[
 \sigma'(z)=\frac1{\pi\sqrt{\mu^2-z^2}}
 -\frac1{\pi\sqrt{K^2-z^2}}\geq\kappa,
 \qquad \sigma(0)=0.
\]
Equality of the two endpoint floors implies
$A(r)-A(r+p)<1$.  Consequently
\[
 1>\int_r^{r+p}\sigma(z)\dd z
 \geq\frac\kappa2\bigl((r+p)^2-r^2\bigr),
\]
which is exactly \eqref{eq:first-shelf-sharp-cap}.
\end{proof}

The proposition identifies the sharp scalar which remains.  The point
$q$ is a one-interface start, so Theorem~\ref{thm:one-interface-complete}
gives $D_A(q)\geq0$ and, more precisely,
\begin{equation}
 D_A(q)=D_K(q)+\lambda_q-2\int_q^\mu G_\mu(z)\dd z,
 \qquad
 \lambda_q=
 \sbr{G_K(q)+\tfrac14}-\sbr{A(q)+\tfrac14}.
 \label{eq:first-shelf-terminal-identity}
\end{equation}
Thus every tail satisfying
\begin{equation}
 d_\rho(n-p)\geq p
 \label{eq:first-shelf-long-head}
\end{equation}
is already proved by
\eqref{eq:first-shelf-reduction} and $R_p\geq-p/2$.  Any residual tail
must satisfy the explicit interface-band restriction
\begin{equation}
 n<\frac{1+d_\rho}{d_\rho}
 \left(\sqrt{r^2+a_\rho K}-r\right),
 \label{eq:first-shelf-residual-band}
\end{equation}
and the sole remaining inequality is
\begin{equation}
 \boxed{
 D_A(q)+R_p+\frac{d_\rho}{2}(n-p)\geq0.}
 \label{eq:terminal-shelf-scalar}
\end{equation}
Unlike a local bad-cell pairing, this scalar retains the terminal convex
reserve forced by the counterexamples above.

\begin{lemma}[Scale-free shelf curvature]
\label{lem:scale-free-shelf-curvature}
In the notation of Proposition~\ref{prop:first-shelf-reduction}, suppose
$p\geq1$ and put
\[
 \Delta=A(r)-A(r+p).
\]
Then
\begin{equation}
 \mathcal C_p\geq
 \frac{p^2}{3(2r+p)}\,\Delta.
 \label{eq:scale-free-shelf-curvature}
\end{equation}
Consequently, since
$\Delta=\varepsilon_0-\varepsilon_p$ on the first ordinary-floor shelf,
\begin{equation}
 R_p\geq
 \frac{p^2}{3(2r+p)}
  (\varepsilon_0-\varepsilon_p)
 +p\left(\varepsilon_0+\varepsilon_p-\frac12\right).
 \label{eq:scale-free-shelf-reserve}
\end{equation}
Moreover,
\begin{equation}
 A(0)-A(r)\leq\frac{r\Delta}{2p}.
 \label{eq:pre-shelf-drop}
\end{equation}
\end{lemma}

\begin{proof}
On the inner branch $\sigma=-A'$ is nonnegative, increasing, and convex,
with $\sigma(0)=0$.  Indeed,
\[
 \sigma'(z)=\frac1{\pi\sqrt{\mu^2-z^2}}
 -\frac1{\pi\sqrt{K^2-z^2}},
\]
and differentiating once more preserves the sign.  Use the positive hinge
representation
\[
 \sigma(z)=az+\int_{[0,\infty)}(z-s)_+\,\dd\nu(s),
 \qquad a\geq0,\quad \nu\geq0.
\]
Both sides of \eqref{eq:scale-free-shelf-curvature} are linear in
$\sigma$.  For the generator $\sigma(z)=z$ their ratio is
\[
 \frac{p^3/6}{p(r+p/2)}
 =\frac{p^2}{3(2r+p)}.
\]
For a hinge with $s\leq r$, the same ratio is
\[
 \frac{p^3/6}{p(r-s+p/2)}
 \geq\frac{p^2}{3(2r+p)}.
\]
If $r<s<r+p$ and $L=r+p-s$, the ratio is
$p-2L/3\geq p/3$, which is again at least the required constant.
Hinges to the right contribute zero.  This proves
\eqref{eq:scale-free-shelf-curvature}, and
\eqref{eq:scale-free-shelf-reserve} follows from
\eqref{eq:first-shelf-exact}.

Finally, convexity and $\sigma(0)=0$ imply that $\sigma(z)/z$ is
nondecreasing.  Hence
\[
 \int_0^r\sigma(z)\dd z\leq\frac r2\sigma(r),
 \qquad
 \Delta\geq p\sigma(r),
\]
which proves \eqref{eq:pre-shelf-drop}.
\end{proof}

\begin{lemma}[Fractional-top convex reserve]
\label{lem:convex-level-reserve}
Let $g$ be nonnegative, strictly decreasing, convex, and $C^1$ on
$[0,L)$, with $g(L)=0$, and extend it by zero.  Set
\[
 v=g(0),\qquad c_v=-g'(0)>0,
\]
and
\[
 D_g=2\int_0^Lg(s)\dd s-
 \left(
 \sbr{g(0)+\tfrac14}
 +2\sum_{j\geq1}\sbr{g(j)+\tfrac14}
 \right),
 \qquad B=\sbr{g(0)+\tfrac14}.
\]
For $1\leq k\leq B$, let $y_k$ be the unique point satisfying
$g(y_k)=k-1/4$, and put $c_k=-g'(y_k)$.  Then
\begin{equation}
 \boxed{
 D_g\geq\sum_{k=1}^{B}
 \left(\frac1{2c_k}-1\right)
 +\frac{(v-B)_+^2}{c_v}.}
 \label{eq:convex-level-reserve}
\end{equation}
\end{lemma}

\begin{proof}
Let $y(t)$ be the decreasing inverse of $g$, extended by zero for
$t\geq g(0)$.  Convexity of $g$ makes $y$ convex; the zero extension
remains convex because its right derivative is zero.  Layer cake gives
\[
 2\int_0^Lg=2\int_0^{g(0)}y(t)\dd t.
\]
The level $k-1/4$ contributes exactly
$2\lceil y_k\rceil-1\leq2y_k+1$ to the strict tail, including when
$y_k$ is integral.  The tangent to $y$ at $k-1/4$ gives
\[
 \int_{k-1}^{k}y(t)\dd t
 \geq y_k-\frac14y'(k-\tfrac14)
 =y_k+\frac1{4c_k}.
\]
These disjoint unit intervals cover $[0,B]$.  If $v>B$, the supporting
tangent at the top endpoint is
\[
 y(t)\geq\frac{v-t}{c_v}\qquad(B\leq t\leq v),
\]
and therefore the unused top interval contributes
\[
 2\int_B^vy(t)\dd t\geq\frac{(v-B)^2}{c_v}.
\]
If $v\leq B$ this interval is empty.  Adding this contribution proves
\eqref{eq:convex-level-reserve}.  At a top action wall the equality level
is absent from $B$, exactly as required by the strict convention.
\end{proof}

\begin{corollary}[Exact-angle terminal reserve]
\label{cor:exact-angle-terminal-reserve}
At the one-interface point $q$, put
\[
 B=\sbr{G_K(q)+\tfrac14},
 \qquad Q=\sbr{A(q)+\tfrac14},
 \qquad v_q=G_K(q),
 \qquad \beta_q=\frac{\arccos(q/K)}\pi.
\]
For $1\leq k\leq B$, let $\theta_k\in(0,\pi/2)$ be the unique solution
of
\[
 \frac K\pi\bigl(\sin\theta_k-\theta_k\cos\theta_k\bigr)
 =k-\frac14.
\]
Then
\begin{equation}
 \boxed{
 D_A(q)\geq
 \max\left\{0,\,
 \frac\pi2\sum_{k=1}^{B}\frac1{\theta_k}
 -Q-2\int_q^\mu G_\mu(z)\dd z
 +\frac{(v_q-B)_+^2}{\beta_q}
 \right\}.}
 \label{eq:exact-angle-terminal-reserve}
\end{equation}
\end{corollary}

\begin{proof}
Apply Lemma~\ref{lem:convex-level-reserve} to
$g(s)=G_K(q+s)$.  At the level $k-1/4$ the ball slope is
$c_k=\theta_k/\pi$, while the top slope is $\beta_q$, so
\[
 D_K(q)\geq\sum_{k=1}^{B}
 \left(\frac\pi{2\theta_k}-1\right)
 +\frac{(v_q-B)_+^2}{\beta_q}.
\]
Insert this in \eqref{eq:first-shelf-terminal-identity}; the $-B$ from
the sum cancels the $+B$ in $\lambda_q=B-Q$.  The other entry in the
maximum is the completed one-interface theorem.
\end{proof}

\begin{corollary}[Zero-level terminal triangle]
\label{cor:zero-level-terminal-triangle}
Suppose $q\leq\mu<q+1$ and
$B=\sbr{G_K(q)+1/4}=0$.  Put
\[
 v=G_K(\mu),\qquad c=\frac{\arccos\rho}{\pi}.
\]
Then the shifted shell tail beginning at $q$ has zero discrete count and
\begin{equation}
 \boxed{
 D_A(q)=2\int_q^K A(z)\dd z
 \geq2\int_\mu^KG_K(z)\dd z
 \geq\frac{v^2}{c}.}
 \label{eq:zero-level-terminal-triangle}
\end{equation}
Consequently the scalar \eqref{eq:terminal-shelf-scalar} is nonnegative
whenever
\begin{equation}
 \frac{v^2}{c}\geq\frac{p-d_\rho(n-p)}2.
 \label{eq:zero-level-shelf-payment}
\end{equation}
\end{corollary}

\begin{proof}
The condition $B=0$ is $G_K(q)\leq3/4$.  Since $A\leq G_K$ and both
actions decrease, every strict bracket in the shell tail is zero, including
at the equality wall.  On $[\mu,K]$ one has $A=G_K$.  Moreover $G_K$ is
convex there and lies above its tangent at $\mu$,
\[
 G_K(z)\geq v-c(z-\mu).
\]
The tangent reaches zero no later than $K$, because $G_K(K)=0$; integrating
its positive triangle gives $2\int_\mu^KG_K\geq v^2/c$.  Finally use
$R_p\geq-p/2$ in \eqref{eq:first-shelf-reduction}.
\end{proof}

\subsection{A certified first-floor large ray}

We record a computer-assisted result for the ordinary first-floor branch
$F_0=1$, $p<n$.  Put
\[
 x=r+p,\qquad y=x+1,\qquad s=n-p-1,\qquad q=x+s+1,
\]
so that $q\leq\mu<q+1$.  At the first-shelf wall $A(x)=3/4$, set
\[
 v=A(y),\qquad
 \lambda=\frac{K-\mu}{\pi\mu},\qquad
 c=\frac1\pi\arccos\frac\mu K,\qquad
 S_z=\sqrt{\mu^2-z^2},
\]
and define the terminal envelope
\begin{equation}
 \mathcal E(y,v)=
 2\int_y^\mu\bigl[v-\lambda(S_y-S_z)\bigr]_+\dd z
 +\frac{(v-\lambda S_y)_+^2}{c}.
 \label{eq:first-floor-terminal-envelope}
\end{equation}
Before specializing to the wall, concavity on the first shelf and its first
drop cell proves \eqref{eq:terminal-shelf-scalar} whenever
\[
 p\bigl(A(r)+A(x)\bigr)+A(x)+A(y)+\mathcal E(y,A(y))-(1+2p)\geq0.
\]
At $A(x)=3/4$ this becomes the weakened wall certificate
\begin{equation}
 \mathcal W:=p\left(A(r)-\frac54\right)+v-\frac14
 +\mathcal E(y,v)
 \label{eq:first-floor-weakened-certificate}
\end{equation}
For fixed $(\mu,r,p)$, every term before subtracting the constant in the
unspecialized certificate is strictly increasing in $K$.  Indeed, the
sampled actions increase, while $v/\lambda$ and $\lambda$ increase in the
integral in \eqref{eq:first-floor-terminal-envelope}; on the positive outer
branch, $(v-\lambda S_y)/\arccos(\mu/K)$ also increases.  Thus the minimum
in a fixed first-floor chamber is its post-crossing value at the unique wall
$A(x)=3/4$.

\begin{theorem}[First-floor large-ray certificate]
\label{thm:first-floor-large-ray}
On the implicit wall $A(x)=3/4$, one has
\begin{equation}
 \boxed{\mathcal W>0}
 \qquad\text{if}\qquad
 r\geq\frac32,\quad s\in\{0,1\},\quad x\geq100.
 \label{eq:first-floor-large-ray}
\end{equation}
The assertion holds on the continuous relaxation of $r,p$ and
$\alpha=\mu-q\in[0,1]$, and therefore on the required half-integral and
integral chamber parameters.
\end{theorem}

\begin{proof}
Write
\[
 \delta=K-\mu,\qquad t=x^{-1/3},\qquad
 \kappa=t\delta,\qquad P=tp,
 \qquad \beta=\mu-x\in[s+1,s+2].
\]
Changing variables $R=x+u$ in the wall equation gives the exact identity
\begin{equation}
 \frac{3\pi}{4}=\int_\beta^{\beta+\delta}
 \frac{\sqrt{u(2x+u)}}{x+u}\dd u.
 \label{eq:first-floor-compact-wall}
\end{equation}
The elementary bounds
\[
 \sqrt{\frac{u}{2x}}
 \leq\frac{\sqrt{u(2x+u)}}{x+u}
 \leq\sqrt{\frac{2u}{x}}
\]
on the wall interval, with a preliminary contradiction if that interval
reaches $u=x$, give
\begin{equation}
 \frac{22}{25}<\kappa<\frac{731}{250}.
 \label{eq:first-floor-kappa-box}
\end{equation}
Moreover, a possible failure must have
$0<A(r)-3/4<1/2$.  Comparing the radius derivatives at $r$ and $x$ then
gives
\[
 A(r)-\frac34\geq
 \frac{(22/25)P}
 {2\pi(583/500)\sqrt{(1083/500)(3571/1000+P)}}.
\]
At $P=41$ the right side is strictly larger than $1/2$; hence every
possible failure lies in $0\leq P<41$.

It remains to certify a compact four-variable box while retaining the
implicit wall.  For $0\leq\xi\leq1$ put
$w_\xi=\beta t+\xi\kappa$ and define
\begin{align*}
 \Phi_x(\xi)&=\frac{\sqrt{w_\xi(2+t^2w_\xi)}}{1+t^2w_\xi},\\
 \Phi_r(\xi)&=
 \frac{\sqrt{(w_\xi+P)(2+t^2(w_\xi-P))}}{1+t^2w_\xi},\\
 \Phi_y(\xi)&=
 \frac{\sqrt{((\beta-1)t+\xi\kappa)
 (2+t^2((\beta+1)t+\xi\kappa))}}
 {1+t^2(\beta t+\xi\kappa)}.
\end{align*}
At $R=\mu+\xi\delta$ these are the exact scaled radius profiles for
$x,r,y$.  For a profile $F$, put
\[
 T_2[F]=\frac{F(0)+2F(1/2)+F(1)}4,
 \qquad M_2[F]=\frac{F(1/4)+F(3/4)}2.
\]
The radius profile is increasing and strictly concave, so two-panel
trapezoid and midpoint quadrature enclose every wall root by
\begin{equation}
 \frac\kappa\pi T_2[\Phi_x]\leq\frac34
 \leq\frac\kappa\pi M_2[\Phi_x].
 \label{eq:first-floor-wall-enclosure}
\end{equation}
The same lower quadrature gives
\[
 \underline A_r=\frac\kappa\pi T_2[\Phi_r],
 \qquad \underline v=\frac\kappa\pi T_2[\Phi_y].
\]

Put $L=\beta-1$ and
\[
 H=\frac34-
 \frac{\kappa\sqrt t\sqrt{\beta(2+\beta t^3)}}
 {\pi(1+\beta t^3)},\qquad
 \vartheta=\frac\kappa{1+\beta t^3},\qquad
 \mathcal Q=\frac{\sqrt{\vartheta(2+t^2\vartheta)}}\pi.
\]
The anchored slope envelope gives $v-\lambda S_y\geq H$.  When $H>0$,
\eqref{eq:first-floor-terminal-envelope} and
$c/t\leq\mathcal Q$ therefore yield the stable lower certificate
\begin{equation}
 t\mathcal W\geq
 \mathcal C_*:=P\left(\underline A_r-\frac54\right)
 +t\left(\underline v-\frac14+2LH\right)
 +\frac{H^2}{\mathcal Q}.
 \label{eq:first-floor-stable-certificate}
\end{equation}
This expression has a finite continuous limit at $t=0$.

The directed-rounding subdivision covers
\[
 0\leq t\leq\frac{53861}{250000},\qquad
 s+1\leq\beta\leq s+2,\qquad0\leq P\leq41,
\]
encloses every compatible $\kappa$ through
\eqref{eq:first-floor-wall-enclosure} and
\eqref{eq:first-floor-kappa-box}, and retains the exact domain condition
\[
 1-Pt^2-\frac32t^3\geq0
\]
equivalent to $r\geq3/2$.  Each terminal box is proved either outside this
domain, automatic from $\underline A_r\geq5/4$, or strictly positive from
$H>0$ and \eqref{eq:first-floor-stable-certificate}.  At 384-bit Arb
precision the two runs process, respectively,
\[
 17697\quad(s=0),\qquad25367\quad(s=1)
\]
boxes, to maximum depths $21$ and $22$.  The identical certificate passes
at 512 bits, and a 768-bit replay with tighter wall-root bisection also
passes.  The verifier is
\path{scripts/general_d_first_floor_ray_verify.py}.  All acceptance tests
use directed interval signs; floating point only orders the subdivision
queue.  This proves \eqref{eq:first-floor-large-ray}.
\end{proof}

The stronger certificate needed on the compact part of the wall keeps the
actual interface geometry instead of the relaxed terminal envelope.  Put
\[
 q=x+s+1,\qquad \alpha=\mu-q\in[0,1],
\]
and, on $A(x)=3/4$, write
\[
 a_r=A(r),\qquad u=A(q),\qquad h=G_K(\mu),\qquad
 c=\frac1\pi\arccos\frac\mu K.
\]
Concavity of $A$ on $[r,\mu]$ and the outer tangent triangle give
\begin{equation}
 \boxed{
 \mathcal C_{\rm ff}:=
 p\left(a_r+\frac34\right)
 +(s+1)\left(\frac34+u\right)
 +\alpha(u+h)+\frac{h^2}{c}-(1+2p)
 \leq D_A(r).}
 \label{eq:first-floor-true-certificate}
\end{equation}
The right-hand side is interpreted as the post-crossing value.  At the
wall itself the strict bracket can only lower the count.

\begin{theorem}[Finite first-floor wall certificate]
\label{thm:first-floor-finite}
On the implicit wall $A(x)=3/4$,
\begin{equation}
 \boxed{\mathcal C_{\rm ff}>0}
 \qquad\text{if}\qquad
 r\geq\frac32,\quad s\in\{0,1\},\quad x<100.
 \label{eq:first-floor-finite}
\end{equation}
Consequently Theorem~\ref{thm:first-floor-large-ray} and
\eqref{eq:first-floor-finite} close the complete sector
$r\geq3/2$, $s\in\{0,1\}$.
\end{theorem}

\begin{proof}
For fixed $x$, the wall has the radius-average form
\[
 \int_\mu^{K(\mu)}\sqrt{1-\frac{x^2}{R^2}}\dd R
 =\frac{3\pi}{4},
\]
and hence
\[
 K'(\mu)=
 \frac{\sqrt{1-x^2/\mu^2}}{\sqrt{1-x^2/K^2}}\in(0,1).
\]
Along this wall $A(r)$ is nonincreasing, $A(q)$ is increasing, and both
$h$ and $c$ are decreasing.  These signs follow by differentiating the
radius average; for example
\[
 \frac{\dd}{\dd\mu}A(z)
 =\frac{\sqrt{1-x^2/\mu^2}}\pi
 \left(
 \frac{\sqrt{1-z^2/K^2}}{\sqrt{1-x^2/K^2}}
 -\frac{\sqrt{1-z^2/\mu^2}}{\sqrt{1-x^2/\mu^2}}
 \right),
\]
and the squared ratio in parentheses has derivative with the sign of
$z^2-x^2$.  Also $K'<1$ makes $\mu/K$ increase, while direct
differentiation gives $h'<0$.

Thus, on an exact rational panel
$\alpha_0\leq\alpha\leq\alpha_1$, strict rational wall brackets
$K_i^-<K(q+\alpha_i)<K_i^+$ give a one-sided lower bound for every term
in \eqref{eq:first-floor-true-certificate}: evaluate $a_r,h$ at the
upper endpoint with $K_1^-$, evaluate $u$ at the lower endpoint with
$K_0^-$, and bound $c$ from above at the lower endpoint with $K_0^+$.
The wall roots are isolated by exact rational bisection, and all
transcendental signs are then decided with directed-rounding Arb
arithmetic.

The exact enumeration is
\[
 r=\frac m2,\quad3\leq m\leq199,\qquad
 p\geq0,\quad r+p<100,\qquad s\in\{0,1\}.
\]
It contains $9801$ chambers for each $s$, hence $19602$ in total.
At both $384$ and $512$ bits the verifier accepts $19606$ panels after
processing $19610$ panels and isolating $39208$ endpoint roots.  Only
$(r,p,s)=(3/2,4,0),\ldots,(3/2,7,0)$ require one bisection in $\alpha$;
no deeper split occurs.  A $768$-bit replay with $160$ root bisections per
endpoint gives the identical counts.  The checked script is
\path{scripts/general_d_first_floor_finite_verify.py}.  Since every leaf
has a strictly positive Arb lower endpoint, this proves
\eqref{eq:first-floor-finite}.
\end{proof}

\begin{theorem}[Complete small-start certificate]
\label{thm:first-floor-small-starts}
On the implicit wall $A(x)=3/4$, one has
\begin{equation}
 \boxed{\mathcal C_{\rm ff}>0}
 \qquad\text{if}\qquad
 r\in\left\{\frac12,1\right\},\quad
 s\in\{0,1\},\quad p\geq0.
 \label{eq:first-floor-small-starts}
\end{equation}
Together with Theorems~\ref{thm:first-floor-large-ray} and
\ref{thm:first-floor-finite}, this closes the whole ordinary first-floor
first-drop sector $s\in\{0,1\}$.
\end{theorem}

\begin{proof}
For $x<100$, the correlated wall-panel construction in the proof of
Theorem~\ref{thm:first-floor-finite} applies without change.  The exact
enumeration
\[
 r\in\left\{\frac12,1\right\},\qquad
 s\in\{0,1\},\qquad p\geq0,\qquad r+p<100
\]
contains $398$ chambers.  Directed-rounding replays at $384$ and $512$
bits process $420$ panels, accept $409$ leaves, and isolate $807$ endpoint
roots, with maximum depth one.  A $1024$-bit replay with $200$ root steps
has identical counts.  The verifier is
\path{scripts/general_d_first_floor_small_starts_verify.py}.

It remains to treat $x\geq100$.  Put
$\beta=\mu-x=s+1+\alpha\in[1,3]$ and $\delta=K-\mu$.  The wall equation is
\begin{equation}
 \frac{3\pi}{4}=\int_\beta^{\beta+\delta}
 \frac{\sqrt{v(2x+v)}}{x+v}\dd v.
 \label{eq:small-start-wall-integral}
\end{equation}
It first implies $\delta<x-\beta$: otherwise the integration interval
contains $[x/2,x]$, where the integrand is at least $\sqrt5/3$, already a
contradiction at $x=100$.  Hence
\[
 \delta\leq
 \left(\frac{9\pi\sqrt2}{8}\right)^{2/3}x^{1/3}
 <\frac{731}{250}x^{1/3}.
\]
With $T_*=53861/250000>100^{-1/3}$ this gives, exactly,
\[
 \frac Kx<1+3T_*^3+\frac{731}{250}T_*^2<\frac65.
\]

For $f_z(R)=\sqrt{1-z^2/R^2}$, the ratio $f_r/f_x$ decreases with $R$.
Since $r\leq1$, $x\geq100$, and $K<6x/5$,
\[
 \left(\frac{f_r(R)}{f_x(R)}\right)^2
 \geq\frac{K^2-r^2}{K^2-x^2}
 >\frac{36x^2-25}{11x^2}
 \geq\frac{14399}{4400}>\frac{25}{9}.
\]
Integration over $R\in[\mu,K]$ therefore yields
\[
 A(r)>\frac53A(x)=\frac54.
\]
Finally, if $v=A(x+1)$, the one-cell action loss is smaller than
$\arccos(\mu/K)/\pi<1/2$, so $v>1/4$.  The weaker wall certificate
\[
 p\left(A(r)-\frac54\right)+v-\frac14+\mathcal E
\]
is thus strictly positive, since $\mathcal E\geq0$.  This proves the
analytic ray and completes \eqref{eq:first-floor-small-starts}.
\end{proof}

The last possible first-floor escape requires a separate limiting estimate,
because no complete terminal level is present below the first action wall.

\begin{lemma}[Uniform critical first-floor gap]
\label{lem:critical-first-floor-gap}
Let
\[
 H_\kappa(X)=\frac{c_0}{\sqrt\kappa}
 \bigl((\kappa+X)^{3/2}-X^{3/2}\bigr),
 \qquad w=H_\kappa(0)=c_0\kappa,
\]
where $c_0=2\sqrt2/(3\pi)$.  For $0<w<3/4$, let
$H_\kappa(M)=3/4$ and $H_\kappa(N)=1$, and set
\[
 \mathcal L(w)=-2\int_M^N(1-H_\kappa(X))\dd X+\frac M2,
 \qquad \mathcal Z(w)=\frac\pi{\sqrt2}w^2.
\]
Then
\begin{equation}
 \boxed{\mathcal L(w)+\mathcal Z(w)>
 \frac\pi{8\sqrt2}.}
 \label{eq:critical-first-floor-gap}
\end{equation}
At $w=3/4$ the same conclusion holds with the smaller wall-safe terminal
value $\pi/(2\sqrt2)$.
\end{lemma}

\begin{proof}
The function $1-H_\kappa$ is convex between $M$ and $N$, so its endpoint
chord gives $\mathcal L\geq(3M-N)/4$.  Write
\[
 X=\kappa u(t),\qquad
 u(t)=\frac{(1-t^2)^2}{4t^2},\qquad
 \frac{H_\kappa(X)}w=\frac{3/t+t^3}{4}.
\]
If $\xi,\eta$ are the parameters at $M,N$, then
\[
 w=\frac{3\xi}{3+\xi^4}=\frac{4\eta}{3+\eta^4}.
\]
Using that $u$ decreases, the ranges
$0<w\leq1/2$, $1/2\leq w\leq2/3$, and $2/3\leq w<3/4$ give,
respectively,
\[
 \frac{\sqrt2}{\pi}(\mathcal L+\mathcal Z)
 \geq\frac18+\frac{85469}{51518208},\quad
 \frac18+\frac{41923}{1281424},\quad
 \frac18+\frac{743}{4608}.
\]
Here one uses $\xi\leq(259/243)w$ in the first range,
$\xi\leq(283/256)w$ in the second, and $\eta\geq1/2$ in the third.
At the wall the one-sided strict convention can lower the normalized
terminal payment by only $1/16$, leaving $1/8+455/4608$.  This proves
\eqref{eq:critical-first-floor-gap}.
\end{proof}

\begin{theorem}[Global exhaustion of the ordinary first floor]
\label{thm:first-floor-global-exhaustion}
In the universal active region $K>\pi/(1-\rho)$, every negative
ordinary first-floor first-drop scalar
\[
 F_0=1,\qquad p<n,\qquad
 D_A(q)+R_p+\frac{d_\rho}{2}(n-p)<0
\]
belongs to a relatively compact subset of the exact shell parameter
space.  Consequently the possible discrete labels $(r,n,p,Q,B)$ are
bounded.  The case $p=0$ is automatic, and the conclusion includes the
whole cone $s=n-p-1\geq2$.
\end{theorem}

\begin{proof}
Put $\tau=\sqrt{1-\rho}$, $a=K-\mu$, and $\kappa=\tau a$.
The fixed-ratio, compact-optical, and thin high-energy results already
proved above show that an escaping negative sequence must have
\[
 \tau\to0,\qquad a\to\infty,\qquad
 \tau\kappa=(1-\rho)a<\frac{125}{8},
\]
while the first-drop interface bound gives $c_0\kappa<3/4$.  The exact
noncancelling turning profile is
\begin{equation}
 \mathcal H_{\tau,\kappa}(X)
 =\frac1\pi\int_0^\kappa
 \frac{\sqrt{(\kappa+X-u)
 [2\kappa-(\kappa+X+u)\tau^2]}}
 {\kappa-u\tau^2}\dd u.
 \label{eq:first-floor-noncancelling-profile}
\end{equation}
On the complete possible negative head its ratio to $H_\kappa(X)$ tends
to one uniformly.

After passage to a subsequence, $\kappa\to\kappa_0\in[0,3/(4c_0)]$.
If $\kappa_0>0$, exact level convergence, the first-shelf reduction, and
the wall-safe terminal estimate give
\[
 \liminf\tau\left(D_A(q)+R_p+\frac{d_\rho}{2}(n-p)\right)
 \geq\mathcal Z(c_0\kappa_0)+\mathcal L(c_0\kappa_0)>0
\]
by Lemma~\ref{lem:critical-first-floor-gap}.

If $\kappa_0=0$, put $Y=\kappa X$.  Uniformly on the physical level
window,
\[
 \mathcal H_{\tau,\kappa}(Y/\kappa)
 \longrightarrow H_0(Y):=\frac{\sqrt2}{\pi}\sqrt Y.
\]
The $3/4$- and $1$-levels converge to
$Y_{3/4}=9\pi^2/32$ and $Y_1=\pi^2/2$.  The exact rescaled prefix then
satisfies
\[
 \begin{aligned}
 \liminf\kappa\tau
 \left(R_p+\frac{d_\rho}{2}(n-p)\right)
 &\geq-2\int_{Y_{3/4}}^{Y_1}(1-H_0(Y))\dd Y
 +\frac{Y_{3/4}}2\\
 &=\frac{17\pi^2}{192}>0.
 \end{aligned}
\]
The nonnegative one-interface reserve is enough on this ray.  Both
possible limits contradict negativity, so no negative sequence escapes.
The two shift grids have $r\geq1/2$, and $r<\rho K$; the fixed-ratio
bound therefore also prevents an escape through $\rho\downarrow0$.
Relative compactness follows.  It bounds the discrete labels, although it
does not assert finiteness of the connected components of the remaining
real-analytic wall set.
\end{proof}

\begin{remark}[First-floor sectors before the cone certificate]
\label{rem:first-floor-large-ray-scope}
Theorems~\ref{thm:first-floor-large-ray} and
\ref{thm:first-floor-finite} close $r\geq3/2$, $s\in\{0,1\}$, while
Theorem~\ref{thm:first-floor-small-starts} closes the remaining two starts.
Thus the whole sector $s\in\{0,1\}$ is proved.
Theorem~\ref{thm:first-floor-global-exhaustion} independently removes every
noncompact part of the complementary cone $s\geq2$; the next theorem proves
that full cone directly in the universal active region.
\end{remark}

\begin{theorem}[Active first-floor cone certificate]
\label{thm:first-floor-s-cone}
On the true ordinary first-floor interface $A(x)=3/4$, let
\[
 x=r+p,\qquad q=x+s+1,\qquad
 \alpha=\mu-q\in[0,1],\qquad
 \beta=s+1+\alpha=\mu-x.
\]
If
\begin{equation}
 r\geq\frac12,\qquad p\geq0,\qquad s\geq2,
 \qquad K-\mu\geq\pi,
 \label{eq:first-floor-s-cone-domain}
\end{equation}
then the true-interface certificate \eqref{eq:first-floor-true-certificate}
satisfies
\begin{equation}
 \boxed{\mathcal C_{\rm ff}>0.}
 \label{eq:first-floor-s-cone-positive}
\end{equation}
Equivalently, \eqref{eq:first-floor-s-cone-positive} holds throughout the
continuous cone $\beta\geq3$ in the closed universal active region.
Together with the no-mode region \eqref{eq:no-mode}, this closes the whole
ordinary first-floor first-drop branch needed for the general-dimensional spectral
theorem.  It does not assert the standalone shifted-tail conjecture below
the active wall.
\end{theorem}

\begin{proof}
Write
\[
 a_r=A(r),\qquad u=A(q),\qquad h=G_K(\mu),\qquad
 c=\frac1\pi\arccos\frac\mu K.
\]
Concavity of $A$ on $[x,\mu]$ gives
\[
 \beta u\geq\alpha\frac34+(s+1)h.
\]
Compressing the two middle trapezoids in
\eqref{eq:first-floor-true-certificate} therefore yields the wall-safe
lower bound
\begin{equation}
 \mathcal C_{\rm ff}\geq
 p\left(a_r-\frac54\right)
 +\beta\left(\frac34+h\right)+\frac{h^2}{c}-1.
 \label{eq:first-floor-beta-compression}
\end{equation}
This retains the true interface and absorbs all dependence on $s$ and
$\alpha$ into $\beta\geq3$.

Define the scale-free ball action and wall variables
\[
 g(z)=\frac{\sqrt{1-z^2}-z\arccos z}{\pi},\qquad
 \rho=\frac\mu K,\quad \xi=\frac xK,\quad \eta=\frac rK,
\]
\[
 w=g(\xi)-\rho g(\xi/\rho),\qquad
 e=\rho-\xi,\qquad \epsilon=1-\rho,
 \qquad g_\rho=g(\rho).
\]
The radius-integral formula
\begin{equation}
 w=\frac1\pi\int_\rho^1
 \sqrt{1-\frac{\xi^2}{R^2}}\dd R
 \label{eq:first-floor-s-cone-wall-integral}
\end{equation}
and the wall equation give
\begin{equation}
 Kw=\frac34,\qquad
 \beta=\frac{3e}{4w},\qquad
 p=\frac{3(\xi-\eta)}{4w},\qquad
 h=\frac{3g_\rho}{4w}.
 \label{eq:first-floor-s-cone-scaling}
\end{equation}
Thus the cone and active conditions are exactly
\begin{equation}
 e\geq4w,\qquad 3\epsilon\geq4\pi w.
 \label{eq:first-floor-s-cone-constraints}
\end{equation}

The ratio of the radius profiles at $r$ and $x$ decreases with the radius.
Averaging on the wall gives
\[
 a_r\geq\frac34Q,\qquad
 Q=\sqrt{\frac{1-\eta^2}{1-\xi^2}}.
\]
Exact one-variable minimization gives
\begin{align}
 (\xi-\eta)(3Q-5)&\geq-2(1-\xi),
 \label{eq:first-floor-s-cone-global-eta}\\
 (\xi-\eta)(3Q-5)&\geq-(1-\xi)
 \quad(\xi\geq9/10).
 \label{eq:first-floor-s-cone-local-eta}
\end{align}
For completeness, put $D=(\xi-\eta)/(1-\xi)$.  The global comparison is
immediate for $D\leq2/5$ and otherwise reduces, after squaring, to
\[
 9D^4-32D^3+24D^2+12D-4>0.
\]
For \eqref{eq:first-floor-s-cone-local-eta}, the only nontrivial case has
$Q<5/3$, whence $Q^2\geq1+4D/5$; after $D\geq1/5$ the comparison reduces to
\[
 \frac{36}{5}D^3-16D^2+10D-1>0.
\]
The first polynomial has exact positive Bernstein coefficients on
\[
 [2/5,1/2],\ [1/2,1],\ [1,3/2],\ [3/2,2],\ [2,3],\ [3,4],
\]
and is positive term by term for $D\geq4$.  The second has the same
certificate on
\[
 [1/5,1/2],\ [1/2,1],\ [1,6/5],\ [6/5,2],\ [2,3],
\]
with a positive termwise tail for $D\geq3$.

Multiplying \eqref{eq:first-floor-beta-compression} by $16w^2/9$ and using
\eqref{eq:first-floor-s-cone-global-eta} gives
\begin{equation}
 \frac{16w^2}{9}\mathcal C_{\rm ff}\geq J_2,
 \qquad
 J_2=w\left(\frac e3-\frac{2\epsilon}{3}\right)
 +eg_\rho+\frac{g_\rho^2}{c}-\frac{16}{9}w^2.
 \label{eq:first-floor-s-cone-J2}
\end{equation}
In the near-tip range $\xi\geq9/10$,
\eqref{eq:first-floor-s-cone-local-eta} improves this to
\begin{equation}
 \frac{16w^2}{9}\mathcal C_{\rm ff}\geq J_1,
 \qquad
 J_1=w\left(\frac{2e}{3}-\frac\epsilon3\right)
 +eg_\rho+\frac{g_\rho^2}{c}-\frac{16}{9}w^2.
 \label{eq:first-floor-s-cone-J1}
\end{equation}

Put $t=e/\epsilon$.  At fixed $\rho$, $w(t)$ is increasing and concave,
and $w(t)/t$ decreases.  At the physical endpoint $\xi=0$ this ratio is
$\epsilon^2/(\pi\rho)$.  If $\rho\leq11/20$, it already exceeds
$\epsilon/4$, contradicting $e\geq4w$.  On
$11/20\leq\rho\leq39/50$, a $4096$-panel directed check at $t_0=3/\pi$
proves
\[
 w(t_0)>\frac{3\epsilon}{4\pi}.
\]
For $t\leq t_0$ the cone constraint fails, and for $t\geq t_0$ the active
constraint fails.  Hence every admissible point has
\begin{equation}
 \rho>\frac{39}{50},\qquad
 \frac\xi\rho\geq\frac{\sqrt7}{4}>\frac{33}{50}.
 \label{eq:first-floor-s-cone-domain-reduction}
\end{equation}
If $t\geq4$, \eqref{eq:first-floor-s-cone-J2} and the active upper bound on
$w$ give directly
\begin{equation}
 J_2\geq\epsilon w
 \left(\frac23-\frac4{3\pi}\right)>0.
 \label{eq:first-floor-s-cone-large-t}
\end{equation}

The cap $\rho>49/50$, $0\leq t<4$, is also analytic.  With
\[
 \rho_0=\frac{49}{50},\quad
 G_0=\frac{2\sqrt{2\rho_0}}{3\pi},\quad
 H_0=\frac{2\rho_0}{3},\quad
 B_0=\frac{2\sqrt2}{3\pi\sqrt{\rho_0}},
 \quad L(t)=(1+t)^{3/2}-t^{3/2},
\]
the radius integrals imply
\begin{equation}
 g_\rho\geq G_0\epsilon^{3/2},\qquad
 \frac{g_\rho}{c}\geq H_0\epsilon,\qquad
 w\leq B_0L(t)\epsilon^{3/2}.
 \label{eq:first-floor-s-cone-cap-bounds}
\end{equation}
Here $\xi>9/10$, so $J_1$ applies.  For $0\leq t\leq1/2$, use
$(1-2t)L(t)\leq1$; for $1/2\leq t<4$, discard the nonnegative linear
$w$-term.  The resulting normalized lower margins are, respectively,
\begin{equation}
 \frac{J_1}{\epsilon^{5/2}}>0.0421933580845,
 \qquad
 \frac{J_1}{\epsilon^{5/2}}>0.108941570080.
 \label{eq:first-floor-s-cone-cap-margins}
\end{equation}

It remains only the compact strip
$39/50\leq\rho\leq49/50$, $0\leq t\leq4$, subject to
\eqref{eq:first-floor-s-cone-constraints}.  The interval calculation follows
the moving boundary $e=4w$ rather than enclosing it by an independent box.
At fixed $\rho$,
\[
 w_e=\frac{\arccos\xi-\arccos(\xi/\rho)}\pi>0,
 \qquad
 w_{ee}=\frac1{\pi\sqrt{1-\xi^2}}
 -\frac1{\pi\sqrt{\rho^2-\xi^2}}<0.
\]
Thus $F(e)=e-4w(e)$ is strictly convex and $F(0)=-4g_\rho<0$; when the
admissible interval is nonempty it has one beta-entry root.  Substitution
of $w=e/4$ at that root removes interval dependency and gives
\begin{equation}
 J_b=\frac{g_\rho^2}{c}+eg_\rho
 -\frac{\epsilon e}{6}-\frac{e^2}{36}.
 \label{eq:first-floor-s-cone-boundary}
\end{equation}
Moreover, with $z=g_\rho/w$ and
\[
 X=\frac13-\frac{2\epsilon}{3e}-\frac{32w}{9e},
\]
concavity, $0<w_e\leq(w-g_\rho)/e$, and differentiation of $J_2$ give
\begin{equation}
 \frac{J_2'}w\geq
 \frac13+z+(1-z)\min(0,X).
 \label{eq:first-floor-s-cone-derivative}
\end{equation}
The directed-rounding certificate proves that both
\eqref{eq:first-floor-s-cone-boundary} and
\eqref{eq:first-floor-s-cone-derivative} are strictly positive on every
accepted panel, so positivity propagates from the moving beta boundary to
the active endpoint.

The verifier is
\path{scripts/general_d_first_floor_s_cone_verify.py}.  Independent replays
at $256$, $384$, and $512$ bits have identical topology: $507$
$\rho$-panels, $253$ $\rho$-splits, $254$ positive panels and no empty
panel, with $2422$ derivative boxes, maximum $\rho$-depth $12$, and maximum
$t$-depth $5$.
The smallest certified boundary lower endpoints are
$5.57124699231\times10^{-9}$, $5.57124666798\times10^{-9}$, and
$5.57124693441\times10^{-9}$ at the three precisions; the smallest
normalized derivative lower endpoint is $0.0010973528439$ in each replay.
All sign decisions use exact rational arithmetic or directed Arb intervals.

The low-rho exclusion, \eqref{eq:first-floor-s-cone-large-t}, the analytic
cap, and the moving-boundary certificate cover the complete domain
\eqref{eq:first-floor-s-cone-domain}, proving
\eqref{eq:first-floor-s-cone-positive}.  Theorems
\ref{thm:first-floor-large-ray}--\ref{thm:first-floor-small-starts} cover
$s\in\{0,1\}$, while the present theorem covers $s\geq2$ on the active
side.  On the complementary side the spectral counting function vanishes
by \eqref{eq:no-mode}; that final observation is used only for the spectral
theorem and not as a below-active shifted-tail assertion.
\end{proof}

\begin{proposition}[Exact no-drop identities]
\label{prop:no-drop-identities}
Suppose $p=n\geq1$, and write
$F_0=\cdots=F_n=f$.  Put
\[
 \varepsilon_q=A(q)+\frac14-f\in[0,1),
 \qquad
 \chi_q=\mathbf 1_{\{A(q)+1/4=f\}}.
\]
Then
\begin{align}
 D_A(r)&=D_A(q)+R_n+\chi_q,
 \label{eq:no-drop-wall-slack}\\
 R_n&=-\frac n2+2n\varepsilon_q
 +2\int_0^n u\,\sigma(r+u)\dd u.
 \label{eq:no-drop-head-exact}
\end{align}
Thus the no-drop branch is automatic when
$f=0$, or $\varepsilon_q\geq1/4$, or
\[
 2\int_0^n u\,\sigma(r+u)\dd u
 \geq n\left(\frac12-2\varepsilon_q\right).
\]
Away from the endpoint wall, the exact shifted-tail condition is the
quantitative payment
\begin{equation}
 D_A(q)\geq
 \frac n2-2n\varepsilon_q
 -2\int_0^n u\,\sigma(r+u)\dd u>0.
 \label{eq:no-drop-terminal-payment}
\end{equation}
At the endpoint wall the right side is reduced by the favorable unit
$\chi_q=1$ in \eqref{eq:no-drop-wall-slack}; proving the wall-safe scalar
\eqref{eq:terminal-shelf-scalar} directly asks for the stronger displayed
payment.
\end{proposition}

\begin{proof}
Strict decrease shows that a lower ordinary-floor wall cannot occur at an
inner sample before $q$ without ending the shelf.  Hence the strict
brackets at $r,\ldots,q$ equal $f$, except that the bracket at $q$ is
$f-1$ when $\chi_q=1$.  Therefore
\[
 T_A(r)-T_A(q)=2nf-\chi_q,
\]
which proves \eqref{eq:no-drop-wall-slack}.  Since
$f=A(q)+1/4-\varepsilon_q$,
\[
 \begin{aligned}
 R_n
 &=2\int_0^n\bigl(A(r+t)-A(q)\bigr)\dd t
   +2n\left(\varepsilon_q-\frac14\right)\\
 &=2\int_0^n\int_t^n\sigma(r+u)\dd u\,\dd t
   +2n\left(\varepsilon_q-\frac14\right),
 \end{aligned}
\]
and reversing the order proves \eqref{eq:no-drop-head-exact}.
The remaining statements follow from $D_A(q)\geq0$.
\end{proof}

\subsection{High floors: an automatic thin sector and the critical ray}

Assume that the first ordinary floor is $F_0=f\geq2$, and write
\begin{equation}
 \varepsilon=1-\rho,\qquad a=\varepsilon K=K-\mu,
 \qquad h=f-\frac14,\qquad m=n-p.
 \label{eq:high-floor-scaled-data}
\end{equation}
Since $A(0)=a/\pi>A(r)\geq h$, every active high-floor configuration
satisfies
\begin{equation}
 a>\pi h.
 \label{eq:high-floor-active-width}
\end{equation}
The reduced scalar is
\[
 \mathcal S=D_A(q)+R_p+\frac{d_\rho}{2}m.
\]

For $b>0$, put $B_b(y)=\pi^{-1}\sqrt{b^2-y^2}$.  The radius-average
formula for the shell action gives, for $0\leq y\leq a\rho$,
\begin{equation}
 \frac12B_a(y)+\frac1{2\rho}B_{a\rho}(y)
 \leq \mathcal A_{\varepsilon,a}(y)
 :=A_{1-\varepsilon,a/\varepsilon}(y/\varepsilon)
 \leq B_a(y).
 \label{eq:product-action-sandwich}
\end{equation}

\begin{lemma}[Quantitative product-shelf margin]
\label{lem:product-shelf-margin}
Let $f\geq3$, $h=f-1/4$, $a>\pi h$, and
$y_b=\sqrt{a^2-\pi^2h^2}$.  If $0\leq y_0\leq s\leq y_b$, then
\begin{equation}
 \boxed{
 2\int_{y_0}^{s}B_a(y)\dd y-2f(s-y_0)+\frac{a-s}{2}
 \geq\frac{f^2}{4a}.}
 \label{eq:product-shelf-margin}
\end{equation}
\end{lemma}

\begin{proof}
For fixed $y_0$ the left side is concave in $s$, so it is enough to use
$s=y_0$ and $s=y_b$.  The first value is at least
$(a-y_b)/2\geq\pi^2h^2/(4a)>f^2/(4a)$.
At $s=y_b$, put $R=a/\pi$.  Differentiation with respect to $y_0$ gives
$2(f-B_a(y_0))$; hence the minimum occurs where $B_a(y_0)=f$ if
$R\geq f$, and at $y_0=0$ if $h\leq R<f$.

In the first case set
$u=\sqrt{R^2-h^2}$ and $v=\sqrt{R^2-f^2}$.  The chord estimate gives
\[
 \frac{\mathcal L}{\pi}>
 \frac{2R-3u+v}{4}
 \geq\frac{3h^2/2-f^2}{4R}
 >\frac{f^2}{4\pi^2R}.
\]
The last comparison follows from $\pi^2>9$ and
$112f^2-216f+27>0$.  In the second case write $R=h+t$,
$0\leq t<1/4$.  Rationalizing the same chord estimate gives
\[
 \frac{\mathcal L}{\pi}\geq\frac{P(h,t)}{6R},\qquad
 P(h,t)=(h+t)^2-4(1-t)^2t(2h+t)
 \geq h^2-2h-\frac14.
\]
Here $18(h^2-2h-1/4)-3f^2=15f^2-45f+45/8>0$ for $f\geq3$.
Again $\mathcal L/\pi>f^2/(4\pi^2R)$, which is
\eqref{eq:product-shelf-margin}.
\end{proof}

\begin{proposition}[An automatic finite-shell sector]
\label{prop:ultrathin-high-floor-sector}
In the high-floor notation above, suppose
\begin{equation}
 f\geq3,\qquad \varepsilon a^4\leq\frac{f^3}{144}.
 \label{eq:ultrathin-sector}
\end{equation}
Then
\begin{equation}
 \boxed{
 \varepsilon\left(R_p+\frac{d_\rho}{2}m\right)
 \geq\frac{f^2}{8a}>0.}
 \label{eq:ultrathin-positive-prefix}
\end{equation}
In particular, \eqref{eq:terminal-shelf-scalar} holds without using a
terminal reserve.
\end{proposition}

\begin{proof}
Put $y_0=\varepsilon r$, $s=\varepsilon(r+p)$, and
$y_q=\varepsilon q$.  Since $A(r+p)\geq h$ and the upper half of
\eqref{eq:product-action-sandwich} holds, $s\leq y_b$.  Conditions
\eqref{eq:high-floor-active-width}--\eqref{eq:ultrathin-sector} imply
$\rho>99/100$ and $y_b\leq a\rho$: indeed
\[
 1-\rho^2<2\varepsilon\leq\frac{f^3}{72a^4}
 <\frac{\pi^2h^2}{a^2}.
\]
Thus the lower half of \eqref{eq:product-action-sandwich} applies on the
whole shelf.  With
\[
 E=\int_0^{y_b}\left(B_a-\rho^{-1}B_{a\rho}\right)\dd y,
\]
direct rationalization, followed by $h\geq11f/12$, gives
\begin{equation}
 E\leq\frac{a^3(\rho^{-2}-1)}{3\pi^2h}
 <\frac{\varepsilon a^3}{11f}
 \leq\frac{f^2}{1584a}.
 \label{eq:product-transfer-error}
\end{equation}
Moreover $y_q\geq a\rho-\varepsilon$, and hence
\begin{equation}
 d_\rho(y_q-s)-(a-s)
 \geq-(1-d_\rho)a-\varepsilon(a+1).
 \label{eq:terminal-coordinate-error}
\end{equation}
Writing $\rho=\cos\theta$ gives
$1-d_\rho=2\theta/\pi\leq\sqrt{2\varepsilon}$; condition
\eqref{eq:ultrathin-sector} then yields
\[
 \frac{(1-d_\rho)a}{2}\leq\frac{f^2}{24a},
 \qquad
 \frac{\varepsilon(a+1)}2\leq\frac{f^2}{24a}.
\]
Finally, the exact scaled identity is
\[
 \varepsilon\left(R_p+\frac{d_\rho}{2}m\right)
 =2\int_{y_0}^{s}\mathcal A_{\varepsilon,a}(y)\dd y
 -2f(s-y_0)+\frac{d_\rho}{2}(y_q-s).
\]
Combining the last four displays with Lemma~\ref{lem:product-shelf-margin}
gives
\[
 \varepsilon\left(R_p+\frac{d_\rho}{2}m\right)
 \geq\frac{f^2}{4a}-\frac{f^2}{1584a}-\frac{f^2}{12a}
 >\frac{f^2}{8a}.
\]
Ordinary shelf walls and strict terminal walls only improve these bounds.
\end{proof}

The complementary unbounded scaling is the grazing ray
\begin{equation}
 \varepsilon\downarrow0,\qquad
 K\varepsilon^{3/2}\longrightarrow\kappa\in(0,\infty),
 \qquad n\sqrt\varepsilon=O(1).
 \label{eq:critical-grazing-ray}
\end{equation}
Set $c_0=2\sqrt2/(3\pi)$.  At the coordinate
$X=(\mu-z)\sqrt\varepsilon$, the ball formula gives, locally uniformly,
\begin{equation}
 A(z)\longrightarrow
 H_\kappa(X):=\frac{c_0}{\sqrt\kappa}
 \left((\kappa+X)^{3/2}-X^{3/2}\right),
 \qquad H_\kappa(0)=w:=c_0\kappa.
 \label{eq:critical-action-profile}
\end{equation}
If $B=\sbr{w+1/4}$, the exact-angle reserve gives the wall-safe limit
\begin{equation}
 \liminf_{\varepsilon\downarrow0}
 \sqrt\varepsilon\,D_A(q)
 \geq\mathcal T(w):=
 \frac\pi{2\sqrt2}w^{1/3}
 \sum_{k=1}^{B}(k-\tfrac14)^{-1/3}.
 \label{eq:critical-terminal-reserve}
\end{equation}
Indeed, at the level $\tau_k=k-1/4$,
\[
 \sqrt\varepsilon\,\frac\pi{2\theta_k}
 \longrightarrow\frac\pi{2\sqrt2}w^{1/3}\tau_k^{-1/3},
\]
whereas the finite $Q$-term and the interface cap vanish after multiplication
by $\sqrt\varepsilon$.  The fractional-top term in
\eqref{eq:exact-angle-terminal-reserve} is nonnegative and is not needed
below.  At a limiting action wall, the strict value of $B$ is the smaller
one-sided complete-level sum and is therefore the safe convention.

\begin{theorem}[No negative leading obstruction on the critical ray]
\label{thm:critical-ray}
For every fixed limiting first-floor level $f\geq2$, each no-drop or
first-drop configuration on \eqref{eq:critical-grazing-ray} has a
nonnegative leading scalar after \eqref{eq:critical-terminal-reserve} is
included.  If its limiting shelf/head prefix is negative, the terminal
reserve exceeds the magnitude of that prefix strictly.
\end{theorem}

\begin{proof}
Write $X=\kappa u$ and
$t=\sqrt{1+u}-\sqrt u\in(0,1]$.  Direct algebra gives
\begin{equation}
 \frac{H_\kappa(X)}w=F(t):=\frac{3/t+t^3}{4},
 \qquad u=\frac{(1-t^2)^2}{4t^2},
 \qquad H_\kappa'(X)=\frac{3c_0}{2}t.
 \label{eq:critical-t-coordinate}
\end{equation}
Let $M$ be the point $H_\kappa(M)=h=f-1/4$ and $N_f\geq M$ the point
$H_\kappa(N_f)=f$.  Scaling the first-shelf reduction and minimizing over
its starting point shows that the worst first-drop prefix is
\begin{equation}
 \mathcal L=-2\int_M^{N_f}(f-H_\kappa(X))\dd X+\frac M2.
 \label{eq:critical-prefix}
\end{equation}
Indeed, the first shelf gives the displayed integral, the remaining
post-drop length gives $M/2$ because $d_\rho\to1$, and moving the starting
point beyond the $f$-level adds only nonnegative shelf area.
Since $0\leq f-H_\kappa\leq1/4$ on this interval,
$\mathcal L<0$ forces $N_f>2M$.

This last inequality forces
\begin{equation}
 w>\frac34,\qquad t_f>\frac14,
 \label{eq:critical-forced-level}
\end{equation}
where $t_f$ is the parameter at $N_f$.  To see this without an asymptotic
estimate, use $3/(4t)\leq F(t)\leq1/t$ and $h/f\geq7/8$.
If $w\leq3/4$, then $t_f<3/10$ and
\[
 \frac{u_f}{u_h}\leq
 \left(\frac76\right)^2\left(\frac{400}{351}\right)^2<2,
\]
contrary to $N_f>2M$.  If $t_f\leq1/4$, the same computation gives
\[
 \frac{u_f}{u_h}<
 \left(\frac65\right)^2\left(\frac{100}{91}\right)^2<2.
\]
Thus \eqref{eq:critical-forced-level} holds.  Equations
\eqref{eq:critical-t-coordinate}--\eqref{eq:critical-prefix} now give
\[
 -\mathcal L<\frac{N_f-M}{2}
 <\frac1{3c_0}=\frac\pi{2\sqrt2}.
\]
On the other hand, $B\geq1$ and the first term of
\eqref{eq:critical-terminal-reserve} gives
\[
 \mathcal T(w)\geq\frac\pi{2\sqrt2}
 w^{1/3}(3/4)^{-1/3}>\frac\pi{2\sqrt2}.
\]
This strictly pays a negative first-drop prefix.  In the no-drop branch
$M=0$ and $w\geq h$.  If $w\geq f$ its prefix is already nonnegative;
otherwise $F(t_f)=f/w\leq f/h\leq8/7$, so $t_f>1/4$ and the same
comparison applies.
\end{proof}

\begin{corollary}[Uniform critical no-drop gap for $f=2,3$]
\label{cor:critical-low-floor-gap}
For a no-drop sequence on \eqref{eq:critical-grazing-ray} with fixed
$f\in\{2,3\}$,
\begin{equation}
 \boxed{
 \liminf_{\varepsilon\downarrow0}\sqrt\varepsilon\,\mathcal S
 \geq\gamma_0:=\frac\pi{2\sqrt2}
 \left[\left(\frac73\right)^{1/3}-1\right]>\frac13.}
 \label{eq:critical-low-floor-gap}
\end{equation}
\end{corollary}

\begin{proof}
No-drop gives $w\geq f-1/4\geq7/4$.  If the prefix is negative,
Theorem~\ref{thm:critical-ray} bounds its magnitude by
$\pi/(2\sqrt2)$, while the first complete terminal level is at least
$\pi(7/3)^{1/3}/(2\sqrt2)$.  If the prefix is nonnegative the estimate
is stronger.  Finally, $\pi/(2\sqrt2)>15/14$ and
$(7/3)^{1/3}>59/45$.
\end{proof}

\begin{theorem}[Global exhaustion of high-floor first drops]
\label{thm:high-floor-first-drop-exhaustion}
The set of negative high-floor first-drop configurations
\[
 F_0=f\geq2,\qquad p<n,\qquad
 D_A(q)+R_p+\frac{d_\rho}{2}(n-p)<0
\]
is relatively compact in the exact shell parameter space.  In particular,
the possible discrete labels $(f,n,p,Q,B)$ are bounded, although the
bounds furnished by this argument are not numerical.
\end{theorem}

\begin{proof}
Write $\epsilon=1-\rho$, $s=\sqrt\epsilon$, $a=K-\mu$,
$\kappa=sa$, $m=n-p$, $x=r+p$, and $h=f-1/4$.  If
$H(t)=A(\mu-t)$, the first drop gives
\[
 W:=H(0)=G_K(\mu)<h,
 \qquad c_0\kappa<h.
\]
Let $H(M)=h$ and let $T$ be the first point at which $H(T)=f$, or put
$T=\mu$ if that level is absent.  The shelf inequalities and the bound
$0\leq f-H\leq1/4$ on its negative part give
\begin{equation}
 R_p+\frac{d_\rho}{2}m
 \geq\frac{(1+d_\rho)M-T-d_\rho}{2}.
 \label{eq:high-floor-exact-prefix-geometry}
\end{equation}
Thus negativity forces $T>(1+d_\rho)M-d_\rho$.

The required uniform interface control follows from the shell-increment
elasticity inequality
\begin{equation}
 \frac{H(t_2)-W}{H(t_1)-W}
 \geq\sqrt{\frac{t_2(2\mu-t_2)}{t_1(2\mu-t_1)}}
 \qquad(0<t_1<t_2\leq\mu).
 \label{eq:shell-increment-elasticity}
\end{equation}
Indeed, for each $R\in[\mu,K]$ the quotient of
\[
 \sqrt{1-\mu^2/R^2+t(2\mu-t)/R^2}
 -\sqrt{1-\mu^2/R^2}
\]
by $\sqrt{t(2\mu-t)}$ is nondecreasing; integration in $R$ proves
\eqref{eq:shell-increment-elasticity}.  Applying it at $M,T$ along a
negative sequence with $\rho\to1$ yields
\begin{equation}
 f-W\leq1+\frac{\sqrt3}{2}+o(1).
 \label{eq:high-floor-interface-pinning}
\end{equation}

As in the first-floor exhaustion, the fixed-ratio, compact-optical, and
thin high-energy theorems leave only
\[
 s\to0,\qquad a\to\infty,\qquad s\kappa<\frac{125}{8}.
\]
The turning estimate $W=c_0\kappa+o(1)$ and
\eqref{eq:high-floor-interface-pinning} then pin
\begin{equation}
 f-1-\frac{\sqrt3}{2}-o(1)
 \leq c_0\kappa\leq f-\frac14+o(1).
 \label{eq:high-floor-kappa-pinning}
\end{equation}
They also confine the dangerous level band to a bounded turning window.

If $f\to\infty$, the root-free exact-angle reserve contains
$B\geq f-O(1)$ complete levels.  Since
$sK^{1/3}=\kappa^{1/3}$ and
$\sum_{k\leq B}(k-1/4)^{-1/3}\gg B^{2/3}$, it gives
$sD_A(q)\to+\infty$, whereas
\eqref{eq:high-floor-exact-prefix-geometry} and the bounded turning
window give
$s(R_p+d_\rho m/2)\geq-O(1)$.  Hence an escaping negative sequence has
bounded $f$.

Pass finally to a subsequence with fixed $f$ and
$\kappa\to\kappa_0>0$.  Exact profile and level convergence give
\[
 \liminf s\left(R_p+\frac{d_\rho}{2}m\right)
 \geq\mathcal L_f:=-2\int_{M_0}^{N_0}
 (f-H_{\kappa_0}(X))\dd X+\frac{M_0}{2}.
\]
The wall-safe complete-level reserve has limit
\[
 \mathcal T(w)=\frac\pi{2\sqrt2}w^{1/3}
 \sum_{k=1}^{\sbr{w+1/4}}(k-\tfrac14)^{-1/3},
 \qquad w=c_0\kappa_0.
\]
If $\mathcal L_f<0$, the exact prefix geometry forces $N_0>2M_0$.
The $t$-parameterization in
\eqref{eq:critical-t-coordinate} then gives, uniformly for $f\geq2$,
\[
 -\mathcal L_f<\frac{3\pi}{8\sqrt2},
 \qquad \mathcal T(w)>\frac\pi{2\sqrt2},
\]
and hence
\begin{equation}
 \mathcal L_f+\mathcal T(w)>\frac\pi{8\sqrt2}.
 \label{eq:uniform-high-floor-critical-gap}
\end{equation}
If $\mathcal L_f\geq0$, the same comparison forces a complete terminal
level, so the sum is again positive.  This excludes the last escape.
Relative compactness follows, and compactness bounds the discrete labels
without asserting a finite number of analytic wall components.
\end{proof}

\begin{theorem}[Explicit high-floor compact reduction]
\label{thm:high-floor-explicit-compact}
In the high-floor first-drop branch, put
\[
 \mathcal S=D_A(q)+R_p+\frac{d_\rho}{2}(n-p),\qquad
 \epsilon=1-\rho,\qquad s=\sqrt\epsilon.
\]
If $f\geq2$, $p<n$, and $0<s\leq10^{-4}$, then
\begin{equation}
 \boxed{\mathcal S\geq0.}
 \label{eq:explicit-high-floor-no-negative}
\end{equation}
More precisely, the contradictory assumption $\mathcal S<0$ forces
$s\mathcal S>1/2$; no unconditional half-unit gap is asserted.

Consequently, if a high-floor first-drop scalar is negative, then, with
\[
 K_*=156250000000000000,\qquad
 N_*=156249999999999999,\qquad
 F_*=52083333333333333,
\]
its continuous parameters satisfy
\begin{equation}
 \boxed{
 \frac1{312500000000000000}<\rho
 <\frac{99999999}{100000000},\qquad
 \frac{21}{4}<K<K_* ,}
 \label{eq:explicit-high-floor-continuous-box}
\end{equation}
and its discrete data satisfy the exact finite ranges
\begin{equation}
 \boxed{
 \begin{gathered}
 r\in\tfrac12\N,\qquad \tfrac12\leq r<K_*,\qquad
 2\leq n\leq N_*,\qquad 1\leq p\leq n-1,\\
 2\leq f\leq F_*,\qquad
 0\leq Q\leq f-1,\qquad Q\leq B\leq F_*.
 \end{gathered}}
 \label{eq:explicit-high-floor-discrete-box}
\end{equation}
Here the strict terminal counts are
\[
 Q=\sbr{A(q)+\tfrac14},\qquad
 B=\sbr{G_K(q)+\tfrac14}.
\]
The residual continuous chambers also retain
\begin{equation}
 \mu=\rho K,\qquad q=r+n\leq\mu<q+1,\qquad
 K-\mu>\frac{7\pi}{4},
 \label{eq:explicit-high-floor-chamber-one}
\end{equation}
\begin{equation}
 \left\lfloor A(r+j)+\frac14\right\rfloor=f
 \quad(0\leq j\leq p),\qquad
 A(r+p+1)<f-\frac14.
 \label{eq:explicit-high-floor-chamber-two}
\end{equation}
These bounds make the integer data finite; they do not certify the
remaining compact analytic walls.
\end{theorem}

\begin{proof}
We give the quantitative transfer because its contradiction quantifier is
essential.  Write
\[
 a=K-\mu,\qquad \kappa=s a=s^3K,\qquad
 m=n-p,\qquad x=r+p,\qquad h=f-\frac14,
\]
and set $\widehat H(t)=A(\mu-t)$.  Let $M$ be the unique point with
$\widehat H(M)=h$ and let $T$ be the point with
$\widehat H(T)=f$, using $T=\mu$ provisionally if that level is absent.
Under the standing assumption $\mathcal S<0$, the completed
one-interface theorem gives $D_A(q)\geq0$, while the exact shelf geometry
gives
\begin{equation}
 T>(1+d_\rho)M-d_\rho,\qquad
 d_\rho=1-2c_\rho,\qquad
 c_\rho=\frac{\arccos\rho}{\pi}<s.
 \label{eq:explicit-high-floor-prefix-geometry}
\end{equation}
The last inequality follows from
$\arccos(1-s^2)=2\arcsin(s/\sqrt2)<\pi s$.

Put $W=\widehat H(0)=G_K(\mu)$.  If $M\leq100$, the slope bound
$\widehat H'\leq c_\rho$ gives $W>h-1/100$.  If $M>100$, then
\[
 \frac TM>
 1+(1-2s)\frac{99}{100}\geq
 r_0:=1+\frac{4999}{5000}\frac{99}{100}.
\]
The exact shell-increment elasticity estimate
\[
 \frac{\widehat H(T)-W}{\widehat H(M)-W}
 \geq\sqrt{\frac{T(2\mu-T)}{M(2\mu-M)}}
 \geq\frac{T/M}{\sqrt{2T/M-1}}
\]
and the exact rational inequality
$r_0/\sqrt{2r_0-1}>23/20$ imply
\[
 \frac{f-W}{f-W-1/4}>\frac{23}{20}.
\]
This remains valid under the provisional absent-level convention, because
$f-W$ then strictly dominates $\widehat H(T)-W$.  Thus in every case
\begin{equation}
 \boxed{W>f-\frac{23}{12}\geq\frac1{12}.}
 \label{eq:explicit-high-floor-interface-pin}
\end{equation}

Introduce the exact and critical scaled profiles
\[
 \mathcal H_{s,\kappa}(X)=A\left(\mu-\frac Xs\right),\qquad
 H_\kappa(X)=\frac{c_0}{\sqrt\kappa}
 \bigl((\kappa+X)^{3/2}-X^{3/2}\bigr),\qquad
 c_0=\frac{2\sqrt2}{3\pi},
\]
and $w=H_\kappa(0)=c_0\kappa$.  The noncancelling integral formula gives
\begin{equation}
 \sqrt{1-s^2}\,w\leq W\leq\frac{w}{1-s^2}.
 \label{eq:explicit-high-floor-origin-comparison}
\end{equation}
Since \eqref{eq:explicit-high-floor-interface-pin} implies $W>f/24$,
$c_0\pi=2\sqrt2/3<1$, and $(1-s^2)/(24s)>1$,
\[
 A(0)=\frac{w}{c_0\pi s}>
 \frac{(1-s^2)f}{24s}>f.
\]
Thus the level $T$ really exists.

For $X=\kappa u$ the exact shell-envelope comparison yields
\begin{equation}
 \mathcal H_{s,\kappa}(X)\geq
 \frac{2}{3\sqrt2}H_\kappa(X).
 \label{eq:explicit-high-floor-envelope}
\end{equation}
Writing
\[
 t=\sqrt{1+u}-\sqrt u,\qquad
 F(t)=\frac{3/t+t^3}{4},\qquad
 \frac{H_\kappa(\kappa u)}w=F(t),
\]
equations \eqref{eq:explicit-high-floor-interface-pin}--
\eqref{eq:explicit-high-floor-envelope} give, at $X_T=sT$,
\begin{equation}
 F(t_T)<54,\qquad t_T>\frac1{72},\qquad
 u_T=\frac{(1-t_T^2)^2}{4t_T^2}<1296.
 \label{eq:explicit-high-floor-turning-window}
\end{equation}
On this entire window the ratio of the exact integrand to the critical
integrand is
\[
 \frac{\sqrt{1-\mathfrak a s^2}}{1-\mathfrak b s^2},\qquad
 0\leq\mathfrak a<649,\qquad0\leq\mathfrak b\leq1.
\]
Hence, by direct rational estimates at $s\leq10^{-4}$,
\begin{equation}
 \frac{999}{1000}<
 \frac{\mathcal H_{s,\kappa}(X)}{H_\kappa(X)}
 <\frac{1001}{1000}\qquad(0\leq X\leq X_T).
 \label{eq:explicit-high-floor-relative-profile}
\end{equation}

We next record the forced-slope step.  If $M\leq100$, the small-$M$
pinning and \eqref{eq:explicit-high-floor-relative-profile} give
\[
 F(t_T)<\frac{100}{87}\left(\frac{1000}{999}\right)^2
 <\frac65<F(1/3).
\]
If $M>100$, put $X_M=sM$ and $X_x=s(\mu-x)$.  Since
$D_A(q)\geq0$, negativity forces
$P:=R_p+d_\rho m/2<0$.  Concavity of the exact profile gives
\[
 sR_p\geq-\frac{X_T-X_x}{4},\qquad
 sm=X_x-s\alpha\geq X_x-s,\qquad X_x\geq X_M,
\]
and therefore $X_T/X_M>5/2$.  If $t_T\leq1/3$, then
\eqref{eq:explicit-high-floor-relative-profile} and $h/f\geq7/8$ imply
\[
 F(t_M)>\frac67F(t_T)\geq\frac{27}{14}>F(2/5),
 \qquad
 \frac{t_M}{t_T}<
 \frac87\frac{1001}{999}\frac{1891}{1875}<\frac76.
\]
The exact formula for $u(t)$ would then give
\[
 \frac{X_T}{X_M}<\left(\frac76\right)^2
 \left(\frac{25}{21}\right)^2<2,
\]
a contradiction.  Thus in both cases
\begin{equation}
 \boxed{t_T>\frac13.}
 \label{eq:explicit-high-floor-forced-slope}
\end{equation}

The exact derivative formula
\[
 \mathcal H'_{s,\kappa}(X)=\frac1{\pi s}
 \left[\arccos(1-s^2(1+u))-
 \arccos\left(1-\frac{s^2u}{1-s^2}\right)\right]
\]
together with
\[
 \sqrt{2y}\leq\arccos(1-y)
 \leq\frac{\sqrt{2y}}{\sqrt{1-y/2}}
\]
and \eqref{eq:explicit-high-floor-turning-window}--
\eqref{eq:explicit-high-floor-forced-slope} gives the exact bound
\begin{equation}
 \mathcal H'_{s,\kappa}(X)>\frac7{50}
 \quad(X_M\leq X\leq X_T),\qquad
 X_T-X_M<\frac{25}{14}.
 \label{eq:explicit-high-floor-band-width}
\end{equation}
The rational and directed-rounding comparisons suppressed in this
derivative transfer, as well as the other numerical comparisons in the
present proof, are replayed by
\path{scripts/general_d_high_floor_compact_verify.py}; independent runs at
512 and 1024 bits both pass.
The preceding prefix estimate now yields
\begin{equation}
 sP> -\frac{25}{56}-\frac s2.
 \label{eq:explicit-high-floor-prefix-loss}
\end{equation}

Finally, \eqref{eq:explicit-high-floor-relative-profile} and
\eqref{eq:explicit-high-floor-forced-slope} imply
\[
 f<\frac{1001}{1000}wF(t_T)
 <\frac{1001}{1000}w\frac{61}{27},\qquad
 w>\frac{54000}{61061}>\left(\frac{19}{20}\right)^3.
\]
Equations \eqref{eq:explicit-high-floor-origin-comparison} and $s\leq10^{-4}$
then give $W>3/4$, so the strict terminal bracket contains the first ball
level.  If $\theta_1$ is its exact angle, the exact terminal reserve is
wall-safe and gives
\[
 D_A(q)\geq\frac\pi{2\theta_1}-1
 -2\int_q^\mu G_\mu(z)\dd z.
\]
The cap is smaller than $2/5$, while the angle inequality
$\sin\theta-\theta\cos\theta\geq2\theta^3/(3\pi)$ gives
\[
 s\frac\pi{2\theta_1}
 \geq\left(\frac{\pi^2w}{6\sqrt2}\right)^{1/3}
 >w^{1/3}>\frac{19}{20}.
\]
Consequently
\begin{equation}
 sD_A(q)>\frac{19}{20}-\frac{7s}{5}.
 \label{eq:explicit-high-floor-terminal-payment}
\end{equation}
Adding \eqref{eq:explicit-high-floor-prefix-loss} and
\eqref{eq:explicit-high-floor-terminal-payment} gives, still under the
assumption $\mathcal S<0$,
\[
 s\mathcal S>
 \frac{19}{20}-\frac{25}{56}-\frac{19}{100000}>\frac12,
\]
which is impossible.  This proves
\eqref{eq:explicit-high-floor-no-negative}, including every ordinary-floor
and strict-bracket wall.

It remains to derive the displayed box.  For $\rho\leq99/100$, the
fixed-ratio scalar theorem proved above, with
\[
 \eta_\rho>\frac{49}{165000},\qquad a_\rho<623,\qquad
 C_*<\frac{13}{10},
\]
excludes negativity once
\[
 K\geq K_0(\rho)<
 \frac{624}{(49/165000)^2}<7.1\mathbin{\cdot}10^9.
\]
For $\rho>99/100$, the thin high-energy theorem excludes
$K\epsilon^2\geq125/8$, while the result just proved forces
$\epsilon>10^{-8}$ for a negative candidate.  Hence
\[
 K<\frac{125}{8\epsilon^2}<K_*.
\]
The larger value $K_*$ is therefore a global cutoff.  Negativity excludes
$p=0$, so $p\geq1$, $n\geq2$, and $r<\mu<K$, which yields
$\rho>1/(2K_*)$ and $n\leq N_*$.  Moreover
$A(r)\geq7/4$ and $A(0)=(K-\mu)/\pi>A(r)$ give
$K-\mu>7\pi/4$ and $K>21/4$.  Finally, using $\pi>3$,
\[
 f\leq A(r)+\frac14<\frac K3+\frac14,\qquad
 B<G_K(0)+\frac14<\frac K3+\frac14,
\]
so $f,B\leq F_*$.  The shelf drop gives $Q\leq f-1$, and the remaining
ranges and chamber conditions follow directly from their ordinary-floor
and strict-bracket definitions.  This proves
\eqref{eq:explicit-high-floor-continuous-box}--
\eqref{eq:explicit-high-floor-chamber-two}.
\end{proof}

\begin{theorem}[Global high-floor scalar exclusion]
\label{thm:high-floor-global-scalar}
In the high-floor first-drop branch, suppose
\[
 f=F_0\geq2,\qquad p<n,\qquad \rho>\frac{99}{100}.
\]
Then the reduced scalar is strictly positive:
\begin{equation}
 \boxed{
 \mathcal S=D_A(q)+R_p+\frac{d_\rho}{2}(n-p)>0.}
 \label{eq:high-floor-global-positive}
\end{equation}
There is no restriction on $K$, on the terminal counts, or on an
ordinary- or strict-floor wall.

Consequently every negative high-floor first-drop candidate satisfies,
with
\[
 K_7:=\frac{16988400000000}{2401},\qquad
 N_7:=7075551852,\qquad F_7:=2358517284,
\]
the improved continuous bounds
\begin{equation}
 \boxed{
 \frac{2401}{33976800000000}<\rho\leq\frac{99}{100},
 \qquad \frac{21}{4}<K<K_7<7.1\mathbin{\cdot}10^9,}
 \label{eq:high-floor-round-seven-continuous}
\end{equation}
and the exact finite integer ranges
\begin{equation}
 \boxed{
 \begin{gathered}
 r\in\tfrac12\N,\qquad 1\leq2r\leq14151103706,\qquad
 2\leq n\leq N_7,\qquad 1\leq p\leq n-1,\\
 2\leq f\leq F_7,\qquad
 0\leq Q\leq f-1,\qquad Q\leq B\leq F_7.
 \end{gathered}}
 \label{eq:high-floor-round-seven-discrete}
\end{equation}
The chamber still obeys
\begin{equation}
 \mu=\rho K,\qquad q=r+n\leq\mu<q+1,\qquad
 K-\mu>\frac{7\pi}{4},
 \label{eq:high-floor-round-seven-chamber-one}
\end{equation}
\begin{equation}
 \left\lfloor A(r+j)+\frac14\right\rfloor=f
 \quad(0\leq j\leq p),\qquad
 A(r+p+1)<f-\frac14.
 \label{eq:high-floor-round-seven-chamber-two}
\end{equation}
These are improved finite bounds, not a certificate of the residual
$\rho\leq99/100$ analytic walls.
\end{theorem}

\begin{proof}
Put
\[
 c=\frac{\arccos\rho}{\pi},\qquad d=d_\rho=1-2c,
 \qquad \widehat H(t)=A(\mu-t),\qquad
 W=\widehat H(0)=G_K(\mu).
\]
The profile $\widehat H$ is increasing and concave on $[0,\mu]$, with
\begin{equation}
 0\leq\widehat H'(t)\leq c.
 \label{eq:high-floor-global-slope}
\end{equation}
Write
\[
 h=f-\frac14,\qquad x=r+p,\qquad m=n-p,
 \qquad t_x=\mu-x.
\]
The first shelf and its first drop give $W<h$ and determine a unique
$M$ by $\widehat H(M)=h$, with
\begin{equation}
 M\leq t_x<M+1.
 \label{eq:high-floor-global-M-window}
\end{equation}
If the level $f$ exists, let $\widehat H(T)=f$ and put $\zeta=0$.
If it is absent, put $T=\mu$ and
\[
 \zeta=f-\widehat H(T)\in(0,1/4].
\]
Thus the proof keeps the present- and absent-level faces separate.

Set $P=R_p+dm/2$.  Since
$t_x=m+\alpha$ with $\alpha=\mu-q\in[0,1)$,
\eqref{eq:high-floor-global-M-window} gives $m\geq M-1$.
Changing variables in the exact shelf integral gives
\[
 R_p=2\int_{t_x}^{t_x+p}(\widehat H(t)-f)\dd t
 \geq-2\int_M^T(f-\widehat H(t))\dd t.
\]
The function $f-\widehat H$ is convex on $[M,T]$ and has endpoint
heights $1/4$ and $\zeta$.  Its endpoint chord therefore proves the
wall-safe prefix estimate
\begin{equation}
 \boxed{
 P\geq-\left(\frac14+\zeta\right)(T-M)
       +\frac d2(M-1).}
 \label{eq:high-floor-endpoint-chord}
\end{equation}

Put $\delta=h-W>0$.  The ratio assumption implies the strict bounds
\begin{equation}
 c<\frac1{20},\qquad d>\frac9{10}.
 \label{eq:high-floor-ratio-constants}
\end{equation}
Suppose first that $\delta\geq1$.  By
\eqref{eq:high-floor-global-slope},
\begin{equation}
 M\geq\frac\delta c>20.
 \label{eq:high-floor-global-M-lower}
\end{equation}
The shell-increment elasticity
\eqref{eq:shell-increment-elasticity}, together with
$V(t)=t(2\mu-t)$, gives, for $0<t_1<t_2\leq\mu$,
\begin{equation}
 \frac{\widehat H(t_2)-W}{\widehat H(t_1)-W}
 \geq\sqrt{\frac{V(t_2)}{V(t_1)}}
 \geq\frac{t_2/t_1}{\sqrt{2t_2/t_1-1}}.
 \label{eq:high-floor-global-elasticity}
\end{equation}
If the level $f$ exists, applying this at $M,T$ yields
\[
 \frac{T/M}{\sqrt{2T/M-1}}
 \leq1+\frac1{4\delta}\leq\frac54,
\]
and hence $T/M\leq5/2$.  Equations
\eqref{eq:high-floor-endpoint-chord} and
\eqref{eq:high-floor-global-M-lower} now give
\begin{equation}
 P\geq\left(\frac d2-\frac38\right)M-\frac d2>1.
 \label{eq:high-floor-present-level-margin}
\end{equation}

If the level is absent, put $a=1/4-\zeta\in[0,1/4]$.  Then
\eqref{eq:high-floor-global-elasticity} gives
\[
 \frac{\mu/M}{\sqrt{2\mu/M-1}}
 \leq1+\frac a\delta\leq1+a.
\]
For $R\geq1$, the upper solution of
$r/\sqrt{2r-1}=R$ is
\[
 r_*(R)=R^2+R\sqrt{R^2-1}.
\]
On $0\leq a\leq1/4$ one has the strict elementary estimate
\begin{equation}
 \left(\frac12-a\right)(r_*(1+a)-1)<\frac{19}{50}.
 \label{eq:high-floor-absent-algebra}
\end{equation}
Indeed, with $x=\sqrt a$ and
\[
 U(x)=\left(\frac12-x^2\right)
 \left(2x^2+x^4+\frac32x(1+x^2)\right),
\]
the bound $\sqrt{a(2+a)}\leq3\sqrt a/2$ reduces the left side to
$U(x)$.  After $y=2x$, all degree-$32$ Bernstein coefficients of
$19/50-U(y/2)$ are positive; the smallest is
$309/396800$.  Thus, on this absent-level face,
\begin{equation}
 P>\left(\frac d2-\frac{19}{50}\right)M-\frac d2
 >\frac9{10}>0.
 \label{eq:high-floor-absent-level-margin}
\end{equation}
This completes both level cases when $\delta\geq1$.

It remains to consider $0<\delta<1$.  The strict interface count is now
fixed without a wall limit:
\begin{equation}
 B_0:=\sbr{W+\tfrac14}=f-1,\qquad
 W=f-\frac14-\delta>f-\frac54\geq\frac34.
 \label{eq:high-floor-global-B-zero}
\end{equation}
Put
\[
 y=f-W=\delta+\frac14\in\left(\frac14,\frac54\right).
\]
We first prove the scale-free level-distance estimate
\begin{equation}
 \boxed{T<\frac{3y}{c},}
 \label{eq:high-floor-global-level-distance}
\end{equation}
which in particular proves that the level $f$ is present.

For fixed $\rho$, the increment
$E_K(t)=\widehat H(t)-W$ has the scaling form
$E_K(t)=K E_1(t/K)$.  Since $E_1$ is concave and $E_1(0)=0$,
\[
 \partial_KE_K(t)=E_1(u)-uE_1'(u)\geq0,
 \qquad u=t/K.
\]
At fixed $y$, the distance to the level $W+y$ is therefore largest at
the smallest admissible $W$.  Since $W=f-y\geq2-y$, it suffices to take
$f=2$ and $W=2-y$.

Let
\[
 \theta=\arccos\rho,\qquad \epsilon=1-\rho,
 \qquad \Phi(\theta)=\sin\theta-\theta\cos\theta,
 \qquad W=K\frac{\Phi(\theta)}\pi,
\]
and set
\[
 t_0=\frac{3y}{c}=\frac{3\pi y}{\theta},
 \qquad u_0=\frac{t_0}{\mu}.
\]
The elementary trigonometric bounds valid for
$\rho>99/100$ are
\begin{equation}
 \theta^2<\frac{21}{1000},\qquad
 \frac{499}{1000}\theta^2<\epsilon\leq\frac12\theta^2,
 \qquad
 \frac{9979}{30000}\theta^3<\Phi(\theta)<\frac13\theta^3.
 \label{eq:high-floor-global-trig}
\end{equation}
They imply
\begin{equation}
 u_0<\frac9{250},\qquad
 \frac{u_0}{\epsilon}>
 \frac{199}{100}\frac yW.
 \label{eq:high-floor-global-u-zero}
\end{equation}
In particular $t_0<\mu$.  Directly from the ball actions,
\begin{equation}
 E_K(t_0)=\frac\mu\pi\int_0^{u_0}
 \left[
 \arccos\bigl(\rho(1-u)\bigr)-\arccos(1-u)
 \right]\dd u.
 \label{eq:high-floor-global-increment}
\end{equation}
Using
\[
 \sqrt{2z}\leq\arccos(1-z)
 \leq\frac{\sqrt{2z}}{\sqrt{1-z/2}},
\]
then setting $u=\epsilon v$ and retaining the positive interval
$0\leq v\leq V$, where
\[
 V=\frac{199}{100}\frac yW,
\]
gives
\begin{equation}
 E_K(t_0)>
 \frac{147}{100}W\int_0^V
 \left[
 \sqrt{1+\frac{99}{100}v}-\frac{101}{100}\sqrt v
 \right]\dd v.
 \label{eq:high-floor-global-integral-lower}
\end{equation}
The integrand is positive on the retained interval.  Define
\[
 I(V)=\frac{200}{297}
 \left(\left(1+\frac{99}{100}V\right)^{3/2}-1\right)
 -\frac{101}{150}V^{3/2}.
\]
With $x=y/W=y/(2-y)$, one has $1/7<x<5/3$.  The function
\[
 J(x)=\frac{147}{100}I\left(\frac{199}{100}x\right)-x
\]
is concave, because $I''<0$, and directed endpoint evaluation gives
\begin{equation}
 J(1/7)>\frac3{20},\qquad J(5/3)>\frac7{50}.
 \label{eq:high-floor-global-J-margins}
\end{equation}
Thus \eqref{eq:high-floor-global-integral-lower} gives
$E_K(t_0)>Wx=y$, proving
\eqref{eq:high-floor-global-level-distance}.

Here $\zeta=0$.  Multiplying
\eqref{eq:high-floor-endpoint-chord} by $c$, and using
$cM\geq\delta$, $cT<3(\delta+1/4)$,
\eqref{eq:high-floor-ratio-constants}, and $d<1$, yields
\begin{equation}
 cP>
 -\frac34\left(\delta+\frac14\right)
 +\frac7{10}\delta-\frac1{40}
 >-\frac{21}{80}.
 \label{eq:high-floor-global-prefix-margin}
\end{equation}

At the terminal start put
\[
 B=\sbr{G_K(q)+\tfrac14},\qquad
 Q=\sbr{A(q)+\tfrac14}.
\]
Since $q\leq\mu$ and the ordinary shelf has dropped by $x+1\leq q$,
the strict counts satisfy
\begin{equation}
 B\geq B_0=f-1,\qquad Q\leq f-1.
 \label{eq:high-floor-global-terminal-counts}
\end{equation}
For $1\leq k\leq f-1$, the strict inequality in
\eqref{eq:high-floor-global-B-zero} gives
\[
 k-\frac14\leq f-\frac54<W=G_K(\mu).
\]
Hence every retained ball angle in
Corollary~\ref{cor:exact-angle-terminal-reserve} satisfies
$\theta_k<\theta=\pi c$.  The exact-angle reserve and the cap bound
\eqref{eq:cap-bound} therefore give
\begin{equation}
 D_A(q)>
 \frac{f-1}{2c}-(f-1)-\frac25.
 \label{eq:high-floor-global-terminal-payment}
\end{equation}
After multiplication by $c$, the right side is minimized at $f=2$;
thus
\begin{equation}
 cD_A(q)>\frac12-\frac1{20}\frac75=\frac{43}{100}.
 \label{eq:high-floor-global-terminal-margin}
\end{equation}
Adding \eqref{eq:high-floor-global-prefix-margin} and
\eqref{eq:high-floor-global-terminal-margin} proves the strict margin
\begin{equation}
 \boxed{c\mathcal S>\frac{43}{100}-\frac{21}{80}
 =\frac{67}{400}>0.}
 \label{eq:high-floor-global-final-margin}
\end{equation}
Together with \eqref{eq:high-floor-present-level-margin} and
\eqref{eq:high-floor-absent-level-margin}, this proves
\eqref{eq:high-floor-global-positive} on every wall.

For completeness, the finite algebraic and transcendental comparisons
in \eqref{eq:high-floor-absent-algebra} and
\eqref{eq:high-floor-global-trig}--
\eqref{eq:high-floor-global-J-margins} are replayed by
\path{scripts/general_d_high_floor_global_scalar_verify.py}.  Independent
directed runs at $384$, $768$, and $1024$ bits give the same strict
margins.

It remains to derive the improved residual box.  The fixed-ratio scalar
estimate already used in the preceding compactness argument has the
threshold \eqref{eq:fixed-rho-threshold}.  Uniformly for
$\rho\leq99/100$, its audited constants
\[
 \eta_\rho>\frac{49}{165000},\qquad
 a_\rho<623,\qquad C_*<\frac{13}{10},
 \qquad \eta_\rho<\frac13
\]
give
\begin{equation}
 K_0(\rho)<\frac{624}{(49/165000)^2}
 =\frac{16988400000000}{2401}=K_7.
 \label{eq:high-floor-round-seven-cutoff}
\end{equation}
Thus a negative scalar has $\rho\leq99/100$ and $K<K_7$.

Negativity excludes $p=0$, so $p\geq1$ and $n\geq2$.  Since
$r<\mu<K<K_7$ and $r\geq1/2$,
\[
 \rho>\frac1{2K_7}=\frac{2401}{33976800000000},
 \qquad
 2r\leq\lfloor2K_7\rfloor=14151103706,
\]
and, from $r+n\leq\mu<K_7$,
\[
 n\leq\left\lfloor K_7-\frac12\right\rfloor=N_7.
\]
Moreover, $A(r)\geq7/4$ and $A(0)>A(r)$ imply
$K-\mu>7\pi/4$ and $K>21/4$.  Finally, using $\pi>3$,
\[
 f\leq A(r)+\frac14<\frac K3+\frac14,
 \qquad
 B<G_K(0)+\frac14<\frac K3+\frac14,
\]
and exact integer division gives
\[
 \left\lfloor\frac{K_7}{3}+\frac14\right\rfloor
 =2358517284=F_7.
\]
The shelf drop gives $Q\leq f-1$, while $A\leq G_K$ gives $Q\leq B$.
The remaining bounds and chamber conditions are their defining
ordinary-floor and strict-bracket relations.  This proves
\eqref{eq:high-floor-round-seven-continuous}--
\eqref{eq:high-floor-round-seven-chamber-two}.
\end{proof}

\begin{theorem}[Extended global high-floor scalar exclusion]
\label{thm:high-floor-round-eight-scalar}
In the high-floor first-drop branch, suppose
\[
 f=F_0\geq2,\qquad p<n,\qquad
 \frac{603}{625}\leq\rho<1.
\]
Then the reduced scalar is strictly positive:
\begin{equation}
 \boxed{
 \mathcal S=D_A(q)+R_p+\frac{d_\rho}{2}(n-p)>0.}
 \label{eq:high-floor-round-eight-positive}
\end{equation}
There is no restriction on $K$, on the terminal counts, or on an
ordinary- or strict-floor wall.

Consequently every remaining negative high-floor first-drop candidate
satisfies the continuous bounds
\begin{equation}
 \boxed{
 \frac1{87889536}<\rho<\frac{603}{625},
 \qquad \frac{21}{4}<K<43944768,}
 \label{eq:high-floor-round-eight-continuous}
\end{equation}
and the exact finite integer ranges
\begin{equation}
 \boxed{
 \begin{gathered}
 r\in\tfrac12\N,\qquad 1\leq2r\leq87889535,\\
 2\leq n\leq43944767,\qquad 1\leq p\leq n-1,\\
 2\leq f\leq14648256,\\
 0\leq Q\leq f-1,\qquad Q\leq B\leq14648256.
 \end{gathered}}
 \label{eq:high-floor-round-eight-discrete}
\end{equation}
The chamber also retains
\begin{equation}
 \mu=\rho K,\qquad q=r+n\leq\mu<q+1,\qquad
 K-\mu>\frac{7\pi}{4},
 \label{eq:high-floor-round-eight-chamber-one}
\end{equation}
\begin{equation}
 \left\lfloor A(r+j)+\frac14\right\rfloor=f
 \quad(0\leq j\leq p),\qquad
 A(r+p+1)<f-\frac14.
 \label{eq:high-floor-round-eight-chamber-two}
\end{equation}
This theorem sharpens, but does not erase,
Theorem~\ref{thm:high-floor-global-scalar}.  It closes a real ratio
interval; it does not certify the residual compact walls in
\eqref{eq:high-floor-round-eight-continuous}--%
\eqref{eq:high-floor-round-eight-discrete}, and no closure is inferred
from their finiteness.
\end{theorem}

\begin{proof}
It is enough to treat
\[
 \frac{603}{625}\leq\rho\leq\frac{99}{100},
\]
because Theorem~\ref{thm:high-floor-global-scalar} owns the remaining
interval $\rho>99/100$.  Put
\[
 c=\frac{\arccos\rho}{\pi},\qquad d=d_\rho=1-2c,
 \qquad \widehat H(t)=A(\mu-t),\qquad
 W=\widehat H(0)=G_K(\mu),
\]
and set
\[
 h=f-\frac14,\qquad \xi=r+p,\qquad m=n-p,
 \qquad t_\xi=\mu-\xi.
\]
The first shelf and its first drop determine $M$ by
\[
 \widehat H(M)=h,\qquad M\leq t_\xi<M+1.
\]
If level $f$ exists, let $\widehat H(T)=f$ and put $\zeta=0$.
If it is absent, put $T=\mu$ and
\[
 \zeta=f-\widehat H(T)\in(0,1/4].
\]
Thus the two level faces remain separate.  Since
$t_\xi=m+\alpha$ with $\alpha=\mu-q\in[0,1)$, one has $m\geq M-1$.
Changing variables in the exact shelf integral and using the endpoint
chord of the convex function $f-\widehat H$ gives
\begin{equation}
 \boxed{
 P:=R_p+\frac d2m
 \geq-\left(\frac14+\zeta\right)(T-M)
       +\frac d2(M-1).}
 \label{eq:high-floor-round-eight-chord}
\end{equation}
This is the same wall-safe chord used in
\eqref{eq:high-floor-endpoint-chord}, now without replacing $c,d$ by
$1/20,9/10$.

Let $\delta=h-W>0$.  The slope bound
$0\leq\widehat H'\leq c$ gives
\begin{equation}
 M\geq\frac{\delta}{c}.
 \label{eq:high-floor-round-eight-M-lower}
\end{equation}
At the rational endpoint, directed evaluation gives
\begin{equation}
 \rho\geq\frac{603}{625}
 \quad\Longrightarrow\quad
 \theta^2<\frac{71}{1000},\qquad
 c<\bar c:=\frac{339}{4000},
 \qquad \theta=\arccos\rho.
 \label{eq:high-floor-round-eight-ratio}
\end{equation}

Suppose first that $\delta\geq1$.  The shell-increment elasticity
\eqref{eq:shell-increment-elasticity} gives
\begin{equation}
 \frac{\widehat H(t_2)-W}{\widehat H(t_1)-W}
 \geq\frac{t_2/t_1}{\sqrt{2t_2/t_1-1}},
 \qquad 0<t_1<t_2\leq\mu.
 \label{eq:high-floor-round-eight-elasticity}
\end{equation}
If level $f$ is present, applying this at $M,T$ gives $T/M\leq5/2$.
Equations \eqref{eq:high-floor-round-eight-chord} and
\eqref{eq:high-floor-round-eight-M-lower} therefore yield
\[
\begin{aligned}
 P
 &\geq\left(\frac d2-\frac38\right)M-\frac d2\\
 &\geq\frac{1/8-c}{c}-\left(\frac12-c\right)
 =\frac1{8c}-\frac32+c.
\end{aligned}
\]
The final expression decreases on the present range.  Hence
\eqref{eq:high-floor-round-eight-ratio} gives the exact positive ledger
\begin{equation}
 \boxed{P>\frac{80921}{1356000}.}
 \label{eq:high-floor-round-eight-present}
\end{equation}

If level $f$ is absent, put $a=1/4-\zeta\in[0,1/4]$.
Equation \eqref{eq:high-floor-round-eight-elasticity} gives
\[
 \frac{\mu/M}{\sqrt{2\mu/M-1}}
 \leq1+\frac a\delta\leq1+a.
\]
With
\[
 r_*(R)=R^2+R\sqrt{R^2-1},
\]
the degree-$32$ Bernstein certificate already used in
\eqref{eq:high-floor-absent-algebra} proves
\[
 \left(\frac12-a\right)(r_*(1+a)-1)<\frac{19}{50},
 \qquad 0\leq a\leq\frac14.
\]
Its exact smallest Bernstein coefficient is $309/396800$.  Consequently
\[
\begin{aligned}
 P
 &>\left(\frac d2-\frac{19}{50}\right)M-\frac d2\\
 &\geq\frac{3/25-c}{c}-\left(\frac12-c\right)
 =\frac3{25c}-\frac32+c.
\end{aligned}
\]
This expression also decreases on the relevant interval, so
\begin{equation}
 \boxed{P>\frac{307}{452000}>0.}
 \label{eq:high-floor-round-eight-absent}
\end{equation}
This smaller margin is the limiting ledger in the present argument.

It remains to treat $0<\delta<1$.  Strict bracketing fixes the initial
terminal count without a wall limit:
\begin{equation}
 B_0=\sbr{W+\tfrac14}=f-1,\qquad
 W=f-\frac14-\delta>f-\frac54\geq\frac34.
 \label{eq:high-floor-round-eight-B-zero}
\end{equation}
Set
\[
 y=f-W=\delta+\frac14\in\left(\frac14,\frac54\right).
\]
We first prove the refined level-distance estimate
\begin{equation}
 \boxed{T<\frac{3y}{c}.}
 \label{eq:high-floor-round-eight-level-distance}
\end{equation}
It will also show that level $f$ is present.

For fixed $\rho$, write
\[
 E_K(t)=\widehat H(t)-W=K E_1(t/K).
\]
Since $E_1$ is concave and $E_1(0)=0$,
\[
 \partial_KE_K(t)=E_1(u)-uE_1'(u)\geq0,
 \qquad u=t/K.
\]
Thus, at fixed $y$, the distance to $W+y$ is largest at the smallest
admissible $W$.  Since $W=f-y\geq2-y$, it suffices to take
\[
 f=2,\qquad W=2-y,\qquad
 x=\frac yW\in\left[\frac17,\frac53\right].
\]
Put
\[
 g_\rho=G_1(\rho),\qquad
 e_\rho(s)=G_1(\rho-s)-\rho G_1(1-s/\rho)-g_\rho.
\]
The concavity needed below is exact, rather than numerical.  Indeed,
$G_1''(z)=1/[\pi\sqrt{1-z^2}]$, and with $b=1-s/\rho$,
\[
\begin{aligned}
 e_\rho''(s)
 &=G_1''(\rho-s)-\frac1\rho G_1''(1-s/\rho)\\
 &=\frac1{\pi\sqrt{1-\rho^2b^2}}
   -\frac1{\pi\rho\sqrt{1-b^2}}\leq0,
\end{aligned}
\]
because $\rho^2(1-b^2)\leq1-\rho^2b^2$.  Since $K=W/g_\rho$,
\begin{equation}
 \mathcal F_\rho(x)
 :=\frac{E_K(3y/c)}W-x
 =\frac{e_\rho(3xg_\rho/c)}{g_\rho}-x
 \label{eq:high-floor-round-eight-concavity}
\end{equation}
is concave in $x$.  It is therefore enough to prove
$\mathcal F_\rho(x)>0$ at $x=1/7,5/3$.

For the endpoint certificate, let
\[
 \epsilon=1-\rho,\qquad
 \Phi=\sin\theta-\theta\cos\theta,\qquad
 W=K\frac{\Phi}{\pi}.
\]
For $t_0=3y/c$, put $u_0=t_0/\mu$.  Then
\begin{equation}
 u_0=\frac{3x\Phi}{\theta\rho}.
 \label{eq:high-floor-round-eight-u-zero}
\end{equation}
The bounds $\Phi<\theta^3/3$ and
\eqref{eq:high-floor-round-eight-ratio} give
\begin{equation}
 u_0<\frac53\frac{71/1000}{603/625}<\frac18.
 \label{eq:high-floor-round-eight-u-upper}
\end{equation}
There is also the strict lower estimate
\begin{equation}
 \Phi>\frac{\rho\theta^3}{3}.
 \label{eq:high-floor-round-eight-Phi-lower}
\end{equation}
Indeed, the difference between the two sides vanishes at zero and has
derivative
\[
 \theta\Phi+\frac{\theta^3\sin\theta}{3}>0.
\]
Since $\epsilon<\theta^2/2$,
\eqref{eq:high-floor-round-eight-u-zero} and
\eqref{eq:high-floor-round-eight-Phi-lower} imply
\begin{equation}
 \frac{u_0}{\epsilon}>2x.
 \label{eq:high-floor-round-eight-retained-range}
\end{equation}

We now use the noncancelling positive series
\begin{equation}
 \arccos(1-z)=\sqrt{2z}\,S(z),\qquad
 S(z)=\sum_{j\geq0}
 \frac{\binom{2j}{j}}{8^j(2j+1)}z^j.
 \label{eq:high-floor-round-eight-arccos-series}
\end{equation}
If $a_j$ denotes its coefficient, then
\[
 a_0=1,\qquad a_1=\frac1{12},\qquad a_2=\frac3{160},
 \qquad
 \frac{a_{j+1}}{a_j}
 =\frac{(2j+1)^2}{4(j+1)(2j+3)}<1.
\]
Thus, for $0\leq z\leq1/8$,
\begin{equation}
 1+\frac z{12}+\frac{3z^2}{160}
 \leq S(z)
 \leq1+\frac z{12}+\frac{3z^2}{140}.
 \label{eq:high-floor-round-eight-series-bounds}
\end{equation}
The upper coefficient is exact in direction: monotonicity of $a_j$
gives
\[
 \sum_{j\geq2}a_jz^j
 \leq\frac{3z^2}{160}\sum_{k\geq0}z^k
 \leq\frac{3z^2}{160}\frac87
 =\frac{3z^2}{140}.
\]

The ball-action identity
\[
 E_K(t_0)=\frac{\mu}{\pi}\int_0^{u_0}
 \left[
 \arccos\bigl(\rho(1-u)\bigr)-\arccos(1-u)
 \right]\dd u
\]
and the substitution $u=\epsilon v$ now give the exact normalized formula
\begin{equation}
 \frac{E_K(t_0)}W
 =C_\rho\int_0^{u_0/\epsilon}
 \left[
 \sqrt{1+\rho v}\,S\bigl(\epsilon(1+\rho v)\bigr)
 -\sqrt v\,S(\epsilon v)
 \right]\dd v,
 \quad
 C_\rho=\frac{\rho\sqrt2\,\epsilon^{3/2}}{\Phi}.
 \label{eq:high-floor-round-eight-noncancelling}
\end{equation}
Alternating Taylor bounds give
\[
 \epsilon>\theta^2\left(\frac12-\frac{\theta^2}{24}\right),
 \qquad
 \Phi<\theta^3\left(
 \frac13-\frac{\theta^2}{30}+\frac{\theta^4}{840}
 \right).
\]
The function
\[
 t\longmapsto
 \frac{(1/2-t/24)^{3/2}}{1/3-t/30+t^2/840}
\]
is decreasing for $0\leq t\leq71/1000$: after logarithmic
differentiation, its sign reduces to
\[
 84+10t-\frac{t^2}{2}>0.
\]
Directed evaluation at $t=71/1000$ and $\rho=603/625$ therefore proves
\begin{equation}
 \boxed{C_\rho>\frac{36}{25}.}
 \label{eq:high-floor-round-eight-coefficient}
\end{equation}

The integrand in
\eqref{eq:high-floor-round-eight-noncancelling} is positive.  By
\eqref{eq:high-floor-round-eight-retained-range} we may retain only
$0\leq v\leq V:=2x$.  On this interval,
\[
 \epsilon v\leq\frac{22}{625}\frac{10}{3}
 =\frac{44}{375}<\frac18.
\]
Keeping the first three positive terms in the first copy of $S$ and using
the upper bound in
\eqref{eq:high-floor-round-eight-series-bounds} in the subtracted copy
produces
\begin{equation}
\begin{aligned}
 \mathcal I_\rho(V)
={}&\frac{2}{3\rho}
 \left((1+\rho V)^{3/2}-1\right)\\
&+\frac{\epsilon}{30\rho}
 \left((1+\rho V)^{5/2}-1\right)\\
&+\frac{3\epsilon^2}{560\rho}
 \left((1+\rho V)^{7/2}-1\right)\\
&-\frac23V^{3/2}-\frac{\epsilon}{30}V^{5/2}
 -\frac{3\epsilon^2}{490}V^{7/2}.
\end{aligned}
 \label{eq:high-floor-round-eight-integral}
\end{equation}
Partitioning
$[603/625,99/100]$ into $16$ equal rational panels, a directed Arb
evaluation on each entire panel, at each of $x=1/7,5/3$, proves
\begin{equation}
 \mathcal I_\rho(2x)>0,\qquad
 \frac{36}{25}\mathcal I_\rho(2x)-x>0.
 \label{eq:high-floor-round-eight-panels}
\end{equation}
These are interval enclosures of all ratios in each panel, not point
samples.  At $384$, $768$, and $1024$ bits, the smallest directed lower
margin in the second inequality is $0.122748159345\ldots$, and the
smallest retained-integral lower bound is $0.202999800557\ldots$.
Equations \eqref{eq:high-floor-round-eight-noncancelling}--%
\eqref{eq:high-floor-round-eight-panels}, followed by concavity in
\eqref{eq:high-floor-round-eight-concavity}, prove
\eqref{eq:high-floor-round-eight-level-distance}.

Level $f$ is therefore present, so $\zeta=0$.  From
\eqref{eq:high-floor-round-eight-chord},
\eqref{eq:high-floor-round-eight-M-lower}, and
\eqref{eq:high-floor-round-eight-level-distance},
\begin{equation}
\begin{aligned}
 cP
 &>-\frac34\left(\delta+\frac14\right)
   +\left(\frac14+\frac d2\right)\delta
   -\frac{cd}{2}\\
 &=-c\delta-\frac3{16}-\frac{cd}{2}
 >-c-\frac3{16}-\frac{cd}{2}.
\end{aligned}
 \label{eq:high-floor-round-eight-prefix}
\end{equation}
At the terminal start, strict bracketing gives
\[
 B=\sbr{G_K(q)+\tfrac14}\geq f-1,\qquad
 Q=\sbr{A(q)+\tfrac14}\leq f-1.
\]
For every $1\leq k\leq f-1$,
\[
 k-\frac14\leq f-\frac54<W=G_K(\mu)
\]
strictly.  Hence every retained ball angle in
Corollary~\ref{cor:exact-angle-terminal-reserve} is smaller than
$\theta=\pi c$.  That reserve and the cap bound \eqref{eq:cap-bound}
give
\begin{equation}
 cD_A(q)>\frac12-\frac75c.
 \label{eq:high-floor-round-eight-terminal}
\end{equation}
Adding this to \eqref{eq:high-floor-round-eight-prefix}, and using
$d=1-2c$, yields
\[
 c\mathcal S>\frac5{16}-\frac{29}{10}c+c^2.
\]
The right side decreases on the present range.  At $c=\bar c$,
\begin{equation}
 \boxed{
 c\mathcal S>\frac{1182521}{16000000}>0.}
 \label{eq:high-floor-round-eight-final}
\end{equation}
Together with \eqref{eq:high-floor-round-eight-present} and
\eqref{eq:high-floor-round-eight-absent}, this proves
\eqref{eq:high-floor-round-eight-positive}.

It remains to derive the residual box.  For
$\rho<603/625$, monotonicity of $G_1$ and directed evaluation at the
rational endpoint give
\begin{equation}
 \eta_\rho>\frac1{504},\qquad \eta_\rho<\frac19.
 \label{eq:high-floor-round-eight-eta}
\end{equation}
The second inequality also follows from the previously used bound
$G_1(1/2)<1/9$.  Moreover, using $\pi<355/113$,
\[
 a_\rho<2\frac{355}{113}\frac{603}{22}
 =\frac{214065}{1243}.
\]
Since $C_*<13/10$, \eqref{eq:high-floor-round-eight-eta} implies
\begin{equation}
 a_\rho+2\eta_\rho C_*
 <\frac{214065}{1243}+\frac{13}{45}<173.
 \label{eq:high-floor-round-eight-numerator}
\end{equation}
Because $C_*>1/2$,
\[
 \sqrt{a_\rho^2+4a_\rho\eta_\rho C_*+\eta_\rho^2}
 <a_\rho+2\eta_\rho C_*.
\]
The fixed-ratio threshold \eqref{eq:fixed-rho-threshold} therefore obeys
\begin{equation}
 \boxed{
 K_0(\rho)<173\mathbin{\cdot}504^2=43944768.}
 \label{eq:high-floor-round-eight-cutoff}
\end{equation}

Negativity excludes $p=0$, so $p\geq1$ and $n\geq2$.  Since
$r<\mu<K<43944768$ and $r\geq1/2$,
\[
 \rho>\frac1{87889536},\qquad
 2r\leq87889535,\qquad n\leq43944767.
\]
As before, $A(r)\geq7/4$ and $A(0)>A(r)$ give
$K-\mu>7\pi/4$ and $K>21/4$.  Finally, $\pi>3$ gives
\[
 f,B<\frac K3+\frac14<14648256+\frac14,
\]
so the integer bounds in
\eqref{eq:high-floor-round-eight-discrete} follow.  The shelf drop gives
$Q\leq f-1$, while $A\leq G_K$ gives $Q\leq B$; the remaining chamber
conditions are their ordinary-floor and strict-bracket definitions.

All finite comparisons just used are replayed by the Round~8A verifier
\begin{center}
 \small\path{scripts/general_d_high_floor_fixed_ratio_verify.py}.
\end{center}
Independent directed runs at $384$, $768$, and $1024$ bits give the same
exact ledgers and cover the full rational panels.  This proves
\eqref{eq:high-floor-round-eight-continuous}--%
\eqref{eq:high-floor-round-eight-chamber-two}.
\end{proof}

\begin{remark}[Limiting ledger in the extended ratio argument]
\label{rem:high-floor-round-eight-obstruction}
The present-level bound remains positive until
\[
 c<\frac{3-\sqrt7}{4}=0.08856\ldots,
\]
and \eqref{eq:high-floor-round-eight-final} remains positive farther
still.  The first obstruction is the absent-level estimate: its lower
bound changes sign at
\[
 c_{\mathrm{abs}}
 =\frac{\frac32-\sqrt{177/100}}2
 =0.0847932652\ldots,
\qquad
 \rho_{\mathrm{abs}}=\cos(\pi c_{\mathrm{abs}})
 =0.9647285943\ldots.
\]
The endpoint $603/625=0.9648$ lies only
$7.14\mathbin{\cdot}10^{-5}$ above this barrier.  Pushing lower requires
an improvement of the absent-level chord loss $19/50$ or additional
positive terminal information on that face.  This is an obstruction to
the present proof ledger, not a counterexample to the shell scalar.
\end{remark}

\begin{theorem}[Terminally strengthened high-floor exclusion]
\label{thm:high-floor-round-eight-b-scalar}
In the high-floor first-drop branch, suppose
\[
 f=F_0\geq2,\qquad p<n,\qquad
 \frac{477}{500}\leq\rho<1.
\]
Then the reduced scalar is strictly positive:
\begin{equation}
 \boxed{
 \mathcal S=D_A(q)+R_p+\frac{d_\rho}{2}(n-p)>0.}
 \label{eq:high-floor-round-eight-b-positive}
\end{equation}
Consequently every remaining negative high-floor first-drop candidate
satisfies
\begin{equation}
 \boxed{
 \frac1{29931928}<\rho<\frac{477}{500},\qquad
 \frac{21}{4}<K<14965964,}
 \label{eq:high-floor-round-eight-b-continuous}
\end{equation}
and the exact finite integer ranges
\begin{equation}
 \boxed{
 \begin{gathered}
 r\in\tfrac12\N,\qquad 1\leq2r\leq29931927,\\
 2\leq n\leq14965963,\qquad 1\leq p\leq n-1,\\
 2\leq f\leq4988654,\\
 0\leq Q\leq f-1,\qquad Q\leq B\leq4988654.
 \end{gathered}}
 \label{eq:high-floor-round-eight-b-discrete}
\end{equation}
The chamber conditions
\begin{equation}
 \mu=\rho K,\qquad q=r+n\leq\mu<q+1,\qquad
 K-\mu>\frac{7\pi}{4},
 \label{eq:high-floor-round-eight-b-chamber-one}
\end{equation}
\begin{equation}
 \left\lfloor A(r+j)+\frac14\right\rfloor=f
 \quad(0\leq j\leq p),\qquad
 A(r+p+1)<f-\frac14
 \label{eq:high-floor-round-eight-b-chamber-two}
\end{equation}
remain in force.  This theorem supersedes the residual box in
Theorem~\ref{thm:high-floor-round-eight-scalar}; it does not certify the
compact walls left by \eqref{eq:high-floor-round-eight-b-continuous}.
\end{theorem}

\begin{proof}
It is enough to treat
\[
 \frac{477}{500}\leq\rho\leq\frac{99}{100},
\]
because Theorem~\ref{thm:high-floor-global-scalar} owns
$\rho>99/100$.  Retain the notation
\[
 c=\frac{\arccos\rho}{\pi},\qquad d=d_\rho=1-2c,
 \qquad \widehat H(t)=A(\mu-t),\qquad
 W=\widehat H(0)=G_K(\mu),
\]
\[
 h=f-\frac14,\qquad \xi=r+p,\qquad m=n-p,
 \qquad \delta=h-W>0.
\]
The endpoint-chord construction in
\eqref{eq:high-floor-round-eight-chord} and
\eqref{eq:high-floor-round-eight-M-lower} remains unchanged:
\begin{equation}
 P:=R_p+\frac d2m
 \geq-\left(\frac14+\zeta\right)(T-M)
       +\frac d2(M-1),\qquad
 M\geq\frac\delta c.
 \label{eq:high-floor-round-eight-b-chord}
\end{equation}
Here $\widehat H(M)=h$; if level $f$ exists then
$\widehat H(T)=f$ and $\zeta=0$, while on the absent face
$T=\mu$ and $\zeta=f-\widehat H(T)\in(0,1/4]$.
Directed endpoint evaluation gives
\begin{equation}
 \rho\geq\frac{477}{500}
 \quad\Longrightarrow\quad
 (\arccos\rho)^2<\frac{93}{1000},\qquad
 c<\bar c:=\frac{97}{1000}.
 \label{eq:high-floor-round-eight-b-ratio}
\end{equation}

Suppose first that $\delta\geq1$ and put
\[
 \gamma=\frac3{25}-c>0.
\]
On the present-level face, the elasticity estimate
\eqref{eq:high-floor-round-eight-elasticity} gives $T/M\leq5/2$.
Thus \eqref{eq:high-floor-round-eight-b-chord} yields
\[
 cP\geq\left(\frac18-c\right)\delta-\frac{cd}{2}
 >\gamma\delta-\frac{cd}{2},
\]
where $1/8-3/25=1/200$.  On the absent-level face, the degree-$32$
Bernstein comparison with smallest coefficient $309/396800$ gives
\[
 \left(\frac12-a\right)(r_*(1+a)-1)<\frac{19}{50},
 \qquad 0\leq a\leq\frac14.
\]
The same calculation as in
\eqref{eq:high-floor-round-eight-absent}, now retaining $\delta$, gives
the common strict prefix
\begin{equation}
 \boxed{cP>\gamma\delta-\frac{cd}{2}.}
 \label{eq:high-floor-round-eight-b-prefix}
\end{equation}

At the terminal start define
\[
 B_0=\sbr{W+\tfrac14},\qquad
 B=\sbr{G_K(q)+\tfrac14},\qquad
 Q=\sbr{A(q)+\tfrac14}.
\]
Since $q=\mu-\alpha$ with $0\leq\alpha<1$, the slope bound gives
$0\leq A(q)-W<c$.  Hence, on all strict walls,
\begin{equation}
 B\geq B_0,\qquad Q\leq B_0+1.
 \label{eq:high-floor-round-eight-b-counts}
\end{equation}
If $B_0\geq1$, the first $B_0$ ball angles are all smaller than
$\pi c$.  The exact-angle reserve and the cap bound
\eqref{eq:cap-bound} give
\[
 cD_A(q)>
 B_0\left(\frac12-c\right)-\frac{7c}{5}.
\]
Combining this with
\eqref{eq:high-floor-round-eight-b-prefix}, $\delta\geq1$, and
$B_0\geq1$ yields
\begin{equation}
 c\mathcal S>\frac{31}{50}-\frac{39}{10}c+c^2
 \geq\frac{251109}{1000000}>0.
 \label{eq:high-floor-round-eight-b-positive-level}
\end{equation}

It remains in the $\delta\geq1$ sector to consider $B_0=0$.  Strict
bracketing says
\[
 0<W\leq\frac34,\qquad Q\in\{0,1\}.
\]
Moreover $q+1>\mu$, so every sample after $q$ lies on the outer ball
tail and has strict count zero.  If $Q=0$, the shell discrete tail is
empty.  The outer tangent triangle gives
\[
 D_A(q)\geq2\int_\mu^K G_K(z)\dd z\geq\frac{W^2}{c}.
\]
Since $\delta\geq7/4-W$, minimizing the resulting convex quadratic at
$W=\gamma/2$ gives
\begin{equation}
 c\mathcal S>
 \frac{129}{625}-\frac{219}{100}c+\frac34c^2
 \geq\frac{4107}{4000000}>0.
 \label{eq:high-floor-round-eight-b-zero-count}
\end{equation}
This is the limiting ledger of the theorem.

If $Q=1$, then $A(q)>3/4$.  Since $A(q)-W<c\alpha$,
\[
 \alpha>\frac{3/4-W}{c},\qquad
 \frac34-c<W\leq\frac34.
\]
The inner payment $A\geq W$ on $[q,\mu]$, together with the outer
tangent triangle, gives
\[
 cD_A(q)>\frac32W-W^2-c.
\]
After \eqref{eq:high-floor-round-eight-b-prefix} and
$\delta\geq7/4-W$ are added, the expression is concave in $W$.
Its two endpoint lower bounds are
\begin{equation}
 \frac{273}{400}-\frac52c+c^2
 \geq\frac{449409}{1000000},\qquad
 \frac{273}{400}-\frac{119}{50}c-c^2
 \geq\frac{442231}{1000000}.
 \label{eq:high-floor-round-eight-b-one-count}
\end{equation}
Thus every $\delta\geq1$ face is positive.

Now suppose $0<\delta<1$.  Then
\[
 B_0=f-1,\qquad W=f-\frac14-\delta>f-\frac54\geq\frac34,
 \qquad y=f-W=\delta+\frac14.
\]
We need the enlarged level-distance estimate
\begin{equation}
 \boxed{T<\frac{3y}{c}.}
 \label{eq:high-floor-round-eight-b-level-distance}
\end{equation}
The scale reduction and concavity argument in
\eqref{eq:high-floor-round-eight-concavity} is unchanged: it reduces the
claim to $f=2$, $W=2-y$, and the two endpoint values
\[
 x=\frac yW\in\left[\frac17,\frac53\right].
\]
For $t_0=3y/c$ and $u_0=t_0/\mu$, equations
\eqref{eq:high-floor-round-eight-u-zero} and
\eqref{eq:high-floor-round-eight-Phi-lower} now give
\[
 u_0<\frac53\frac{93/1000}{477/500}<\frac16,
 \qquad \frac{u_0}{1-\rho}>2x.
\]

Write $\epsilon=1-\rho$ and use the positive series
\eqref{eq:high-floor-round-eight-arccos-series}.  On the enlarged
subtracted range
\[
 0\leq z\leq\frac{23}{150}
\]
its decreasing coefficients give
\begin{equation}
 1+\frac z{12}+\frac{3z^2}{160}
 \leq S(z)\leq
 1+\frac z{12}+\frac{45z^2}{2032}.
 \label{eq:high-floor-round-eight-b-series}
\end{equation}
The positive first copy of the series is used only for
$z\leq299/1500<1$.  The coefficient in the noncancelling identity
\eqref{eq:high-floor-round-eight-noncancelling} satisfies
\begin{equation}
 C_\rho=\frac{\rho\sqrt2\,\epsilon^{3/2}}
 {\sin\theta-\theta\cos\theta}>\frac75.
 \label{eq:high-floor-round-eight-b-coefficient}
\end{equation}
Indeed, the same decreasing Taylor quotient as in
\eqref{eq:high-floor-round-eight-coefficient}, evaluated at
$\theta^2=93/1000$ and $\rho=477/500$, has this directed lower bound.

Retaining $0\leq v\leq V=2x$ in the positive integral produces
\begin{equation}
\begin{aligned}
 \mathcal I^{\mathrm B}_\rho(V)
={}&\frac{2}{3\rho}
 \left((1+\rho V)^{3/2}-1\right)
 +\frac{\epsilon}{30\rho}
 \left((1+\rho V)^{5/2}-1\right)\\
&+\frac{3\epsilon^2}{560\rho}
 \left((1+\rho V)^{7/2}-1\right)
 -\frac23V^{3/2}-\frac{\epsilon}{30}V^{5/2}
 -\frac{45\epsilon^2}{7112}V^{7/2}.
\end{aligned}
 \label{eq:high-floor-round-eight-b-integral}
\end{equation}
The last coefficient is exactly $(45/2032)(2/7)$.  Partitioning
$[477/500,99/100]$ into $32$ equal rational panels and evaluating each
entire Arb interval at $x=1/7,5/3$ gives
\begin{equation}
 \mathcal I^{\mathrm B}_\rho(2x)>0,\qquad
 \frac75\mathcal I^{\mathrm B}_\rho(2x)-x>0.
 \label{eq:high-floor-round-eight-b-panels}
\end{equation}
The smallest directed lower bounds are $0.2031479283$ for the retained
integral and $0.0607609473$ for the second expression.  Concavity in
$x$ proves \eqref{eq:high-floor-round-eight-b-level-distance}.

Level $f$ is therefore present.  The same chord calculation as in
\eqref{eq:high-floor-round-eight-prefix} gives
\[
 cP>-c-\frac3{16}-\frac{cd}{2}.
\]
The shelf drop \eqref{eq:high-floor-round-eight-b-chamber-two} gives
$Q\leq f-1$, while $B_0=f-1$ supplies $f-1$ angles smaller than
$\pi c$.  Hence
\[
 cD_A(q)>\frac12-\frac75c.
\]
Adding and using $d=1-2c$ yields
\begin{equation}
 c\mathcal S>\frac5{16}-\frac{29}{10}c+c^2
 \geq\frac{40609}{1000000}>0.
 \label{eq:high-floor-round-eight-b-small-delta}
\end{equation}
Together with
\eqref{eq:high-floor-round-eight-b-positive-level}--%
\eqref{eq:high-floor-round-eight-b-one-count}, this proves
\eqref{eq:high-floor-round-eight-b-positive}.

For the new box, a residual ratio $\rho<477/500$ satisfies
\[
 \eta_\rho>\frac1{338},\qquad \eta_\rho<\frac19,
 \qquad
 a_\rho<2\frac{355}{113}\frac{477}{23}
 =\frac{338670}{2599}.
\]
Since $C_*<13/10$,
\[
 a_\rho+2\eta_\rho C_*<131.
\]
The square root in \eqref{eq:fixed-rho-threshold} is smaller than this
same quantity because $C_*>1/2$.  Therefore
\begin{equation}
 \boxed{K_0(\rho)<131\mathbin{\cdot}338^2=14965964.}
 \label{eq:high-floor-round-eight-b-cutoff}
\end{equation}
The arguments following
\eqref{eq:high-floor-round-eight-cutoff} now give
$2r\leq29931927$, $n\leq14965963$, and
$f,B\leq4988654$, together with
\eqref{eq:high-floor-round-eight-b-chamber-one}--%
\eqref{eq:high-floor-round-eight-b-chamber-two}.

All directed and exact comparisons in this proof are replayed by
\begin{center}
 \small\path{scripts/general_d_high_floor_absent_terminal_verify.py}.
\end{center}
Independent runs at $384$, $512$, $768$, and $1024$ bits give the same
rational ledgers and cover all $64$ ratio-panel endpoint evaluations.
\end{proof}

\begin{remark}[The new limiting zero-count ledger]
\label{rem:high-floor-round-eight-b-obstruction}
The weakest polynomial in the strengthened argument is
\eqref{eq:high-floor-round-eight-b-zero-count}.  Its first zero is
\[
 c_*=\frac{219-\sqrt{41769}}{150}
 =0.0975022935\ldots,\qquad
 \cos(\pi c_*)=0.9534519995\ldots.
\]
The rational choice $477/500=0.954$ remains strictly above this point.
The high-floor active-width relation contains additional information not
used in the zero-count estimate, so this sign change is an obstruction to
the present ledger, not a counterexample to the scalar.
\end{remark}

\begin{theorem}[Active-width and sharp-cap high-floor exclusion]
\label{thm:high-floor-round-eight-d-scalar}
In the high-floor first-drop branch, suppose
\[
 f=F_0\geq2,\qquad p<n,\qquad
 \frac{93}{100}\leq\rho<1.
\]
Then
\begin{equation}
 \boxed{
 \mathcal S=D_A(q)+R_p+\frac{d_\rho}{2}(n-p)>0.}
 \label{eq:high-floor-round-eight-d-positive}
\end{equation}
Consequently every remaining negative high-floor first-drop candidate
satisfies
\begin{equation}
 \boxed{
 \frac1{5443200}<\rho<\frac{93}{100},\qquad
 \frac{21}{4}<K<2721600,}
 \label{eq:high-floor-round-eight-d-continuous}
\end{equation}
and
\begin{equation}
 \boxed{
 \begin{gathered}
 r\in\tfrac12\N,\qquad 1\leq2r\leq5443199,\\
 2\leq n\leq2721599,\qquad 1\leq p\leq n-1,\\
 2\leq f\leq907200,\\
 0\leq Q\leq f-1,\qquad Q\leq B\leq907200.
 \end{gathered}}
 \label{eq:high-floor-round-eight-d-discrete}
\end{equation}
The exact chamber conditions remain
\begin{equation}
 \mu=\rho K,\qquad q=r+n\leq\mu<q+1,\qquad
 K-\mu>\frac{7\pi}{4},
 \label{eq:high-floor-round-eight-d-chamber-one}
\end{equation}
\begin{equation}
 \left\lfloor A(r+j)+\frac14\right\rfloor=f
 \quad(0\leq j\leq p),\qquad
 A(r+p+1)<f-\frac14.
 \label{eq:high-floor-round-eight-d-chamber-two}
\end{equation}
This theorem supersedes the residual box in
Theorem~\ref{thm:high-floor-round-eight-b-scalar}.  It certifies the
displayed ratio interval, not the compact walls left by
\eqref{eq:high-floor-round-eight-d-continuous}.
\end{theorem}

\begin{proof}
By Theorem~\ref{thm:high-floor-round-eight-b-scalar}, it is enough to
treat
\[
 \frac{93}{100}\leq\rho\leq\frac{477}{500}.
\]
Retain the notation
\[
 \epsilon=1-\rho,\qquad a=\epsilon K,\qquad
 h=f-\frac14,\qquad
 W=G_K(\mu),\qquad \delta=h-W,
\]
\[
 c=\frac{\arccos\rho}{\pi},\qquad d=1-2c,\qquad
 g_\rho=G_1(\rho).
\]
The active high-floor inequality \eqref{eq:high-floor-active-width} and
$W=ag_\rho/\epsilon$ give, with
\[
 \lambda_\rho=\frac{\pi g_\rho}{\epsilon},
\]
\begin{equation}
 W>\lambda_\rho(W+\delta).
 \label{eq:high-floor-round-eight-d-active}
\end{equation}
Thus $\lambda_\rho<1$.  Moreover
\[
 \pi g_\rho
 =\int_0^\epsilon\arccos(1-y)\dd y
 >\frac{2\sqrt2}{3}\epsilon^{3/2}.
\]
Since $\epsilon\geq23/500$ and
\[
 \frac89\frac{23}{500}>\frac1{25},
\]
one has $\lambda_\rho>1/5$.  As
$\delta\geq7/4-W$, \eqref{eq:high-floor-round-eight-d-active} gives
\begin{equation}
 \boxed{W>\frac74\lambda_\rho>\frac7{20}.}
 \label{eq:high-floor-round-eight-d-W-floor}
\end{equation}
Directed endpoint evaluation also gives
\begin{equation}
 (\arccos\rho)^2<\frac{71}{500},\qquad
 c<\bar c:=\frac{1199}{10000}<\frac3{25}.
 \label{eq:high-floor-round-eight-d-ratio}
\end{equation}

Suppose first that $\delta\geq1$.  The endpoint-chord and
degree-$32$ absent-level arguments used in
\eqref{eq:high-floor-round-eight-b-prefix} remain valid and give
\begin{equation}
 cP>
 \left(\frac3{25}-c\right)\delta-\frac{cd}{2},
 \qquad P=R_p+\frac d2(n-p).
 \label{eq:high-floor-round-eight-d-prefix}
\end{equation}
The exact smallest Bernstein coefficient is still $309/396800$.
For $B_0\geq1$, the Round~8B terminal reserve gives
\[
 c\mathcal S>
 \frac{31}{50}-\frac{39}{10}c+c^2
 \geq\frac{16676601}{100000000}>0.
\]
If $B_0=Q=0$, the outer tangent triangle, together with
\eqref{eq:high-floor-round-eight-d-W-floor}, gives
\[
\begin{aligned}
 c\mathcal S
 &>\left(\frac3{25}-c\right)\left(\frac74-W\right)
   -\frac c2+c^2+W^2\\
 &\geq\frac{7706601}{100000000}>0.
\end{aligned}
\]
Indeed, the derivative in $W\geq7/20$ is larger than
$7/10-3/25=29/50$.  Finally, on the face $B_0=0,Q=1$, the two endpoint
ledgers are
\[
 \frac{273}{400}-\frac52c+c^2
 \geq\frac{39712601}{100000000}>0,
\]
\[
 \frac{273}{400}-\frac{119}{50}c-c^2
 \geq\frac{38276199}{100000000}>0.
\]
This proves positivity on every $\delta\geq1$ face, including the strict
walls $\delta=1$, $W=3/4$, and $A(q)=3/4$.

Now let $0<\delta<1$.  As in Round~8B,
\[
 B_0=f-1,\qquad
 W=f-\frac14-\delta>f-\frac54\geq\frac34,\qquad
 y=f-W=\delta+\frac14.
\]
We first prove
\begin{equation}
 \boxed{T<\frac{3y}{c}.}
 \label{eq:high-floor-round-eight-d-level-distance}
\end{equation}
The scale and concavity reductions in
\eqref{eq:high-floor-round-eight-concavity} reduce the claim to $f=2$
and the endpoints
\[
 x=\frac yW\in\left\{\frac17,\frac53\right\}.
\]
For $t_0=3y/c$ and $u_0=t_0/\mu$, the estimates used in
\eqref{eq:high-floor-round-eight-retained-range} give
\[
 u_0<
 \frac53\frac{71/500}{93/100}
 =\frac{71}{279}<\frac{51}{200}<1,\qquad
 \frac{u_0}{1-\rho}>2x.
\]
Use the positive series
\[
 \arccos(1-z)=\sqrt{2z}\,S(z),\qquad
 S(z)=\sum_{j\geq0}
 \frac{\binom{2j}{j}}{8^j(2j+1)}z^j.
\]
On the subtracted retained range $z\leq7/30$, its decreasing
coefficients give
\begin{equation}
 1+\frac z{12}+\frac{3z^2}{160}
 \leq S(z)\leq
 1+\frac z{12}+\frac{9z^2}{368},
 \label{eq:high-floor-round-eight-d-series}
\end{equation}
where $(3/160)/(1-7/30)=9/368$.  The positive first copy is used only
for $z\leq91/300<1$.  The decreasing Taylor quotient from
\eqref{eq:high-floor-round-eight-b-coefficient}, evaluated at
$\rho=93/100$ and $\theta^2=71/500$, gives
\begin{equation}
 C_\rho=
 \frac{\rho\sqrt2(1-\rho)^{3/2}}
 {\sin\theta-\theta\cos\theta}>
 \frac{1389}{1000}.
 \label{eq:high-floor-round-eight-d-coefficient}
\end{equation}
Retaining $0\leq v\leq V=2x$ therefore produces
\begin{equation}
\begin{aligned}
 \mathcal I^{\mathrm D}_\rho(V)
={}&\frac{2}{3\rho}
 \left((1+\rho V)^{3/2}-1\right)
 +\frac{1-\rho}{30\rho}
 \left((1+\rho V)^{5/2}-1\right)\\
&+\frac{3(1-\rho)^2}{560\rho}
 \left((1+\rho V)^{7/2}-1\right)
 -\frac23V^{3/2}-\frac{1-\rho}{30}V^{5/2}
 -\frac{9(1-\rho)^2}{1288}V^{7/2}.
\end{aligned}
 \label{eq:high-floor-round-eight-d-integral}
\end{equation}
The last coefficient is $(9/368)(2/7)$.  On $64$ equal rational panels
covering $[93/100,477/500]$, directed evaluation at
$x=1/7,5/3$ gives
\begin{equation}
 \mathcal I^{\mathrm D}_\rho(2x)>0,\qquad
 \frac{1389}{1000}\mathcal I^{\mathrm D}_\rho(2x)-x>0.
 \label{eq:high-floor-round-eight-d-panels}
\end{equation}
Across runs at $384$, $512$, $768$, and $1024$ bits, the smallest
retained-integral lower bound exceeds $0.2037148927$.  The smallest
directed margin exceeds $0.0162868443$.
Concavity proves \eqref{eq:high-floor-round-eight-d-level-distance}.

It remains to sharpen the terminal cap.  Since $p<n$, one has $n\geq1$
and $q=r+n\geq3/2$.  If $q=\mu$, the cap vanishes.  If $q<\mu$, strict
convexity gives
\[
 2\int_q^\mu G_\mu(z)\dd z
 <G_\mu(q)(\mu-q)<G_{q+1}(q).
\]
For $H(q)=G_{q+1}(q)$ and $q/(q+1)=\cos\vartheta$,
\[
 H'(q)=\frac{\sin\vartheta-\vartheta}{\pi}<0.
\]
Directed evaluation at the endpoint gives
\begin{equation}
 \boxed{
 2\int_q^\mu G_\mu(z)\dd z
 <G_{5/2}(3/2)<\frac15.}
 \label{eq:high-floor-round-eight-d-cap}
\end{equation}
The endpoint chord and
\eqref{eq:high-floor-round-eight-d-level-distance} yield
\[
 cP>-c-\frac3{16}-\frac{cd}{2}.
\]
The exact-angle reserve, $B_0=f-1$, $Q\leq f-1$, and
\eqref{eq:high-floor-round-eight-d-cap} give
\[
 cD_A(q)>
 \frac{f-1}{2}-c(f-1)-\frac c5
 \geq\frac12-\frac{6c}{5}.
\]
Adding and using $d=1-2c$ proves
\begin{equation}
 c\mathcal S>
 \frac5{16}-\frac{27}{10}c+c^2
 \geq\frac{314601}{100000000}>0.
 \label{eq:high-floor-round-eight-d-small-delta}
\end{equation}
Together with the $\delta\geq1$ ledgers, this proves
\eqref{eq:high-floor-round-eight-d-positive}.

For a residual ratio $\rho<93/100$, monotonicity and directed endpoint
evaluation give
\[
 \eta_\rho>\frac1{180},\qquad \eta_\rho<\frac19,
\]
and
\[
 a_\rho<
 2\frac{355}{113}\frac{93}{7}
 =\frac{66030}{791}.
\]
Since $C_*<13/10$,
\[
 a_\rho+2\eta_\rho C_*
 <\frac{66030}{791}+\frac{13}{45}
 =\frac{2981633}{35595}<84.
\]
The fixed-ratio threshold therefore satisfies
\begin{equation}
 \boxed{K_0(\rho)<84\mathbin{\cdot}180^2=2721600.}
 \label{eq:high-floor-round-eight-d-cutoff}
\end{equation}
The integer arguments following
\eqref{eq:high-floor-round-eight-b-cutoff} give
\eqref{eq:high-floor-round-eight-d-continuous}--%
\eqref{eq:high-floor-round-eight-d-discrete} and the chamber conditions.

All directed and exact comparisons in this proof are replayed by
\begin{center}
 \small\path{scripts/general_d_high_floor_sharp_cap_verify.py}.
\end{center}
The helper hash, the full $128$ panel-endpoint evaluations, the cap
endpoint, every rational ledger, and the cutoff arithmetic are checked
inside the certificate.
\end{proof}

\begin{theorem}[Refined-cap high-floor exclusion]
\label{thm:high-floor-round-eight-e-scalar}
In the high-floor first-drop branch, suppose
\[
 f=F_0\geq2,\qquad p<n,\qquad
 \frac{927}{1000}\leq\rho<1.
\]
Then
\begin{equation}
 \boxed{
 \mathcal S=D_A(q)+R_p+\frac{d_\rho}{2}(n-p)>0.}
 \label{eq:high-floor-round-eight-e-positive}
\end{equation}
Consequently every remaining negative high-floor first-drop candidate
satisfies
\begin{equation}
 \boxed{
 \frac1{4626882}<\rho<\frac{927}{1000},\qquad
 \frac{21}{4}<K<2313441,}
 \label{eq:high-floor-round-eight-e-continuous}
\end{equation}
and
\begin{equation}
 \boxed{
 \begin{gathered}
 r\in\tfrac12\N,\qquad 1\leq2r\leq4626881,\\
 2\leq n\leq2313440,\qquad 1\leq p\leq n-1,\\
 2\leq f\leq771147,\\
 0\leq Q\leq f-1,\qquad Q\leq B\leq771147.
 \end{gathered}}
 \label{eq:high-floor-round-eight-e-discrete}
\end{equation}
The exact chamber conditions remain
\begin{equation}
 \mu=\rho K,\qquad q=r+n\leq\mu<q+1,\qquad
 K-\mu>\frac{7\pi}{4},
 \label{eq:high-floor-round-eight-e-chamber-one}
\end{equation}
\begin{equation}
 \left\lfloor A(r+j)+\frac14\right\rfloor=f
 \quad(0\leq j\leq p),\qquad
 A(r+p+1)<f-\frac14.
 \label{eq:high-floor-round-eight-e-chamber-two}
\end{equation}
This theorem supersedes the residual box in
Theorem~\ref{thm:high-floor-round-eight-d-scalar}.  It certifies the
displayed ratio interval, not the compact walls left by
\eqref{eq:high-floor-round-eight-e-continuous}.
\end{theorem}

\begin{proof}
By Theorem~\ref{thm:high-floor-round-eight-d-scalar}, it is enough to
treat
\[
 \frac{927}{1000}\leq\rho\leq\frac{93}{100}.
\]
If $p=0$, then $R_0=0$, and
Theorem~\ref{thm:one-interface-complete} gives
$\mathcal S=D_A(q)+d_\rho n/2>0$.  A negative candidate must therefore
have $p\geq1$.  Since $p<n$ and $r\geq1/2$, this implies
\begin{equation}
 n\geq2,\qquad q=r+n\geq\frac52.
 \label{eq:high-floor-round-eight-e-terminal-start}
\end{equation}
If $q=\mu$, the inner cap vanishes.  Otherwise strict convexity of
$G_\mu$ gives
\[
 2\int_q^\mu G_\mu(z)\dd z
 <G_\mu(q)(\mu-q)<G_{q+1}(q).
\]
For $H(q)=G_{q+1}(q)$ and $q/(q+1)=\cos\vartheta$,
\[
 H'(q)=\frac{\sin\vartheta-\vartheta}{\pi}<0.
\]
Directed evaluation at the endpoint in
\eqref{eq:high-floor-round-eight-e-terminal-start} therefore proves
\begin{equation}
 \boxed{
 2\int_q^\mu G_\mu(z)\dd z
 <G_{7/2}(5/2)<\frac16.}
 \label{eq:high-floor-round-eight-e-cap}
\end{equation}

Retain
\[
 c=\frac{\arccos\rho}{\pi},\qquad d=1-2c,\qquad
 W=G_K(\mu),\qquad h=f-\frac14,\qquad \delta=h-W.
\]
We first sharpen the absent-level endpoint chord used in Round~8D.
When level $f$ is absent, put
$a=1/4-\zeta\in[0,1/4]$ and
\[
 r_*(R)=R^2+R\sqrt{R^2-1}.
\]
The elasticity reduction for $\delta\geq1$ contains the coefficient
\[
 L(a)=\left(\frac12-a\right)\bigl(r_*(1+a)-1\bigr).
\]
Set
\[
 b(a)=\frac{377}{1000}
      -\left(\frac12-a\right)(2a+a^2)
 =\frac{377}{1000}-a+\frac32a^2+a^3.
\]
The function $b$ decreases on $[0,1/4]$ and
$b(1/4)=1891/8000>0$.  Moreover
\begin{equation}
\begin{aligned}
 &b(a)^2-\left(\frac12-a\right)^2(1+a)^2a(2+a)\\
 &\quad=
 \frac{142129}{10^6}-\frac{627}{500}a
 +\frac{2881}{1000}a^2-\frac{123}{500}a^3-a^4>0.
\end{aligned}
 \label{eq:high-floor-round-eight-e-absent-polynomial}
\end{equation}
Indeed, after $a=y/4$, all degree-$48$ Bernstein coefficients are
positive.  Their smallest member is
\[
 \frac{7422203}{77832000000}.
\]
The sign of $b$ makes the squaring step legitimate, so
\begin{equation}
 \boxed{L(a)<\frac{377}{1000}.}
 \label{eq:high-floor-round-eight-e-absent-chord}
\end{equation}

Put
\[
 \beta=\frac{123}{1000}-c.
\]
If level $f$ is present, the Round~8B elasticity estimate gives
\[
 cP\geq\left(\frac18-c\right)\delta-\frac{cd}{2}
 >\beta\delta-\frac{cd}{2}.
\]
If it is absent, \eqref{eq:high-floor-round-eight-e-absent-chord} gives
\[
 P>\left(\frac d2-\frac{377}{1000}\right)M-\frac d2
   =\beta M-\frac d2.
\]
Directed endpoint evaluation on the present ratio band gives
\begin{equation}
 (\arccos\rho)^2<\frac{37}{250},\qquad
 c<\bar c:=\frac{49}{400}<\frac{123}{1000}.
 \label{eq:high-floor-round-eight-e-ratio}
\end{equation}
Thus $\beta>0$, and $M\geq\delta/c$ yields the common prefix
\begin{equation}
 \boxed{cP>\beta\delta-\frac{cd}{2}.}
 \label{eq:high-floor-round-eight-e-prefix}
\end{equation}

The active-width relation still gives
\[
 W>\lambda_\rho(W+\delta),\qquad
 \lambda_\rho=\frac{\pi G_1(\rho)}{1-\rho}.
\]
Here $\epsilon=1-\rho\geq7/100$, and
\[
 \lambda_\rho>\frac{2\sqrt2}{3}\sqrt\epsilon
 >\frac{249}{1000},
 \qquad
 \frac89\frac7{100}-\left(\frac{249}{1000}\right)^2
 =\frac{1991}{9000000}>0.
\]
Since $W+\delta=f-1/4\geq7/4$,
\begin{equation}
 \boxed{W>\frac{1743}{4000}.}
 \label{eq:high-floor-round-eight-e-W-floor}
\end{equation}

Suppose first that $\delta\geq1$.  On the face $B_0\geq1$, the
exact-angle reserve, $Q\leq B_0+1$, and
\eqref{eq:high-floor-round-eight-e-cap} give
\[
 cD_A(q)>
 B_0\left(\frac12-c\right)-\frac{7c}{6}.
\]
Combining this with \eqref{eq:high-floor-round-eight-e-prefix},
$B_0\geq1$, and $\delta\geq1$, gives
\begin{equation}
 c\mathcal S>
 \frac{623}{1000}-\frac{11}{3}c+c^2
 \geq\frac{90643}{480000}>0.
 \label{eq:high-floor-round-eight-e-B-positive}
\end{equation}
If $B_0=Q=0$, the shell tail is empty and the outer tangent triangle
gives $cD_A(q)\geq W^2$.  Since
$\delta\geq7/4-W$, equations
\eqref{eq:high-floor-round-eight-e-prefix} and
\eqref{eq:high-floor-round-eight-e-W-floor} give
\begin{equation}
 c\mathcal S>
 \beta\left(\frac74-W\right)-\frac c2+c^2+W^2
 \geq\frac{2308663}{16000000}>0.
 \label{eq:high-floor-round-eight-e-Q-zero}
\end{equation}
The derivative in $W$ is $2W-\beta>0$, and the last expression
decreases in $c$, which justifies the displayed endpoints.  Finally,
if $B_0=0,Q=1$, then $3/4-c<W\leq3/4$ and
\[
 cD_A(q)>\frac32W-W^2-c.
\]
The resulting expression is concave in $W$.  Its two endpoint ledgers
are
\begin{equation}
 \frac{1371}{2000}-\frac52c+c^2
 \geq\frac{63081}{160000}>0,
 \qquad
 \frac{1371}{2000}-\frac{2377}{1000}c-c^2
 \geq\frac{303449}{800000}>0.
 \label{eq:high-floor-round-eight-e-Q-one}
\end{equation}
This closes every $\delta\geq1$ terminal face, including the
strict-count walls.

Now let $0<\delta<1$.  The scale and concavity reductions used in
\eqref{eq:high-floor-round-eight-d-level-distance} again reduce
\[
 T<\frac{3y}{c},\qquad y=f-W=\delta+\frac14,
\]
to $f=2$ and $x=y/W\in\{1/7,5/3\}$.  With
$t_0=3y/c$ and $u_0=t_0/\mu$, equation
\eqref{eq:high-floor-round-eight-e-ratio} gives
\[
 u_0<\frac53\frac{37/250}{927/1000}
 =\frac{740}{2781}<\frac{267}{1000}<1,\qquad
 \frac{u_0}{1-\rho}>2x.
\]
On the subtracted retained range,
$z\leq(73/1000)(10/3)=73/300$, so the positive arccos series satisfies
\begin{equation}
 1+\frac z{12}+\frac{3z^2}{160}
 \leq S(z)\leq
 1+\frac z{12}+\frac{45z^2}{1816}.
 \label{eq:high-floor-round-eight-e-series}
\end{equation}
The positive first copy is used only for $z\leq949/3000<1$.
The decreasing Taylor quotient from Round~8D, now evaluated at
$\rho=927/1000$ and $\theta^2=37/250$, gives
\[
 C_\rho=
 \frac{\rho\sqrt2(1-\rho)^{3/2}}
 {\sin\theta-\theta\cos\theta}>
 \frac{173}{125}.
\]
For $V=2x$, the retained lower integral is
\begin{equation}
\begin{aligned}
 \mathcal I^{\mathrm E}_\rho(V)
={}&\frac{2}{3\rho}
 \left((1+\rho V)^{3/2}-1\right)
 +\frac{1-\rho}{30\rho}
 \left((1+\rho V)^{5/2}-1\right)\\
&+\frac{3(1-\rho)^2}{560\rho}
 \left((1+\rho V)^{7/2}-1\right)
 -\frac23V^{3/2}-\frac{1-\rho}{30}V^{5/2}\\
&-\frac{45(1-\rho)^2}{6356}V^{7/2}.
\end{aligned}
 \label{eq:high-floor-round-eight-e-integral}
\end{equation}
On $64$ equal rational panels covering
$[927/1000,93/100]$, directed evaluation at both values of $x$ proves
\[
 \mathcal I^{\mathrm E}_\rho(2x)>0,\qquad
 \frac{173}{125}\mathcal I^{\mathrm E}_\rho(2x)-x>0.
\]
Across $384$, $512$, $768$, and $1024$ bits, the smallest directed
lower bounds are $0.2040442323$ and $0.0085053826$, respectively.
Concavity proves $T<3y/c$ on the whole band.  Consequently
\[
 cP>-c-\frac3{16}-\frac{cd}{2}.
\]
Here $B_0=f-1$ and $Q\leq f-1$, so the refined cap gives
\[
 cD_A(q)>
 \frac{f-1}{2}-c(f-1)-\frac c6
 \geq\frac12-\frac{7c}{6}.
\]
Adding these inequalities and using $d=1-2c$ gives
\begin{equation}
 c\mathcal S>
 \frac5{16}-\frac83c+c^2
 \geq\frac{403}{480000}>0.
 \label{eq:high-floor-round-eight-e-small-delta}
\end{equation}
Together with
\eqref{eq:high-floor-round-eight-e-B-positive}--%
\eqref{eq:high-floor-round-eight-e-Q-one}, this proves
\eqref{eq:high-floor-round-eight-e-positive}.

For a residual ratio $\rho<927/1000$, monotonicity and directed
endpoint evaluation give
\[
 \eta_\rho>\frac1{169},\qquad \eta_\rho<\frac19,
\]
and
\[
 a_\rho<
 2\frac{355}{113}\frac{927}{73}
 =\frac{658170}{8249}.
\]
Since $C_*<13/10$,
\[
 a_\rho+2\eta_\rho C_*
 <\frac{658170}{8249}+\frac{13}{45}
 =\frac{29724887}{371205}<81.
\]
As $C_*>1/2$, the square root in the fixed-ratio threshold is
smaller than the same left side.  Hence
\begin{equation}
 \boxed{K_0(\rho)<81\mathbin{\cdot}169^2=2313441.}
 \label{eq:high-floor-round-eight-e-cutoff}
\end{equation}
The integer arguments following
\eqref{eq:high-floor-round-eight-d-cutoff} give
\eqref{eq:high-floor-round-eight-e-continuous}--%
\eqref{eq:high-floor-round-eight-e-discrete} and the chamber conditions.

All exact and directed comparisons in this proof are replayed by
\begin{center}
 \small\path{scripts/general_d_high_floor_refined_cap_verify.py}.
\end{center}
The certificate pins the Round~8D verifier hash, checks the
degree-$48$ Bernstein ledger, and covers all $128$ panel-endpoint
evaluations at four precisions.
\end{proof}

\begin{theorem}[Signed-prefix high-floor exclusion]
\label{thm:high-floor-round-eight-f-scalar}
In the high-floor first-drop branch, suppose
\[
 f=F_0\geq2,\qquad p<n,\qquad
 \frac{1847}{2000}\leq\rho<1.
\]
Then
\begin{equation}
 \boxed{
 \mathcal S=D_A(q)+R_p+\frac{d_\rho}{2}(n-p)>0.}
 \label{eq:high-floor-round-eight-f-positive}
\end{equation}
Consequently every remaining negative high-floor first-drop candidate
satisfies
\begin{equation}
 \boxed{
 \frac1{3844456}<\rho<\frac{1847}{2000},\qquad
 \frac{21}{4}<K<1922228,}
 \label{eq:high-floor-round-eight-f-continuous}
\end{equation}
and
\begin{equation}
 \boxed{
 \begin{gathered}
 r\in\tfrac12\N,\qquad 1\leq2r\leq3844455,\\
 2\leq n\leq1922227,\qquad 1\leq p\leq n-1,\\
 2\leq f\leq640742,\\
 0\leq Q\leq f-1,\qquad Q\leq B\leq640742.
 \end{gathered}}
 \label{eq:high-floor-round-eight-f-discrete}
\end{equation}
The exact chamber conditions remain
\begin{equation}
 \mu=\rho K,\qquad q=r+n\leq\mu<q+1,\qquad
 K-\mu>\frac{7\pi}{4},
 \label{eq:high-floor-round-eight-f-chamber-one}
\end{equation}
\begin{equation}
 \left\lfloor A(r+j)+\frac14\right\rfloor=f
 \quad(0\leq j\leq p),\qquad
 A(r+p+1)<f-\frac14.
 \label{eq:high-floor-round-eight-f-chamber-two}
\end{equation}
This theorem supersedes the residual box in
Theorem~\ref{thm:high-floor-round-eight-e-scalar}.  It certifies the
displayed ratio interval, not the compact walls left by
\eqref{eq:high-floor-round-eight-f-continuous}.
\end{theorem}

\begin{proof}
By Theorem~\ref{thm:high-floor-round-eight-e-scalar}, it is enough to
treat
\begin{equation}
 \frac{1847}{2000}\leq\rho\leq\frac{927}{1000}.
 \label{eq:high-floor-round-eight-f-band}
\end{equation}
As in the preceding round, $p=0$ is automatic by
Theorem~\ref{thm:one-interface-complete}.  Thus a negative candidate has
$p\geq1$, $n\geq2$, and $q\geq5/2$.  In particular the refined terminal
cap \eqref{eq:high-floor-round-eight-e-cap} remains available:
\begin{equation}
 \boxed{2\int_q^\mu G_\mu(z)\dd z<\frac16.}
 \label{eq:high-floor-round-eight-f-cap}
\end{equation}

Retain
\[
 c=\frac{\arccos\rho}{\pi},\qquad d=1-2c,\qquad
 W=G_K(\mu),\qquad h=f-\frac14,\qquad \delta=h-W,
\]
and put
\[
 P=R_p+\frac d2(n-p),\qquad
 \beta=\frac{123}{1000}-c.
\]
Suppose first that $\delta\geq1$.  If level $f$ is absent, the refined
chord \eqref{eq:high-floor-round-eight-e-absent-chord} gives
\[
 P>\left(\frac d2-\frac{377}{1000}\right)M-\frac d2
   =\beta M-\frac d2.
\]
If the level is present, the elasticity estimate gives instead
\[
 P\geq\left(\frac18-c\right)M-\frac d2
   >\beta M-\frac d2.
\]
Thus both level faces have the common pre-substitution estimate
\begin{equation}
 \boxed{P>\beta M-\frac d2.}
 \label{eq:high-floor-round-eight-f-common-prefix}
\end{equation}

Directed endpoint evaluation on \eqref{eq:high-floor-round-eight-f-band}
gives
\begin{equation}
 (\arccos\rho)^2<\frac{311}{2000}<\frac4{25},
 \qquad c<\bar c:=\frac{627}{5000}.
 \label{eq:high-floor-round-eight-f-ratio}
\end{equation}
Write $\epsilon=1-\rho$ and
\[
 \lambda_\rho=\frac{\pi G_1(\rho)}{1-\rho}.
\]
Since $\epsilon\geq73/1000$ and
\[
 \frac89\frac{73}{1000}-\left(\frac14\right)^2
 =\frac{43}{18000}>0,
\]
the active-width estimate gives
\begin{equation}
 \lambda_\rho>\frac14,\qquad W>\frac7{16}.
 \label{eq:high-floor-round-eight-f-width}
\end{equation}
Moreover, with $\theta=\arccos\rho$,
\begin{equation}
 \frac{c\mu}{W}
 =\frac{\theta\rho}{\lambda_\rho\epsilon}
 <\frac{(2/5)(927/1000)}{(1/4)(73/1000)}
 =\frac{7416}{365}<21.
 \label{eq:high-floor-round-eight-f-radial-ratio}
\end{equation}
When $\beta\geq0$, use $M\geq\delta/c$ in
\eqref{eq:high-floor-round-eight-f-common-prefix}; when $\beta<0$, use
$M\leq\mu$, with the inequality direction reversed by the negative
coefficient.  This gives the two signed prefixes
\begin{equation}
 \boxed{cP>\beta\delta-\frac{cd}{2}}
 \quad(\beta\geq0),
 \label{eq:high-floor-round-eight-f-positive-prefix}
\end{equation}
\begin{equation}
 \boxed{cP>21\beta W-\frac{cd}{2}}
 \quad(\beta<0).
 \label{eq:high-floor-round-eight-f-negative-prefix}
\end{equation}

Define
\[
 B_0=\sbr{W+\tfrac14},\qquad
 B=\sbr{G_K(q)+\tfrac14},\qquad
 Q=\sbr{A(q)+\tfrac14}.
\]
As before, $B\geq B_0$ and $Q\leq B_0+1$.  We now close all terminal
strata for $\delta\geq1$.  If $\beta\geq0$, every lower ledger decreases
up to $c=123/1000$.  The exact-angle reserve,
\eqref{eq:high-floor-round-eight-f-cap}, and
\eqref{eq:high-floor-round-eight-f-positive-prefix} give there
\begin{equation}
 \begin{aligned}
 B_0\geq1:&\qquad
 c\mathcal S>\frac{187129}{1000000},\\
 B_0=Q=0:&\qquad
 c\mathcal S>\frac{580141}{4000000},\\
 B_0=0,\ Q=1:&\qquad
 c\mathcal S>
 \min\left\{\frac{393129}{1000000},\frac{189}{500}\right\}.
 \end{aligned}
 \label{eq:high-floor-round-eight-f-positive-ledgers}
\end{equation}
For the middle line, the outer tangent triangle gives
$cD_A(q)\geq W^2$ and the expression increases for $W\geq7/16$.
For the last line, the expression is concave on
$3/4-c<W\leq3/4$, so its minimum lies at an endpoint.

Now let $\beta<0$.  On the face $B_0\geq1$, strict bracketing gives
$W\leq B_0+3/4$.  The coefficient of $B_0$ after
\eqref{eq:high-floor-round-eight-f-negative-prefix} is added satisfies
\[
 \frac12-c+21\beta\geq\frac{1621}{5000}>0.
\]
Thus $B_0=1$ is the lower endpoint.  On $B_0=Q=0$, the derivative in
$W\geq7/16$ is bounded below by
\[
 2\frac7{16}+21\left(-\frac3{1250}\right)
 =\frac{4123}{5000}>0.
\]
Finally the $B_0=0,Q=1$ expression is concave on
$3/4-c<W\leq3/4$.  All resulting endpoint polynomials decrease in $c$,
and evaluation at $c=\bar c$ gives
\begin{equation}
 \begin{aligned}
 B_0\geq1:&\qquad
 c\mathcal S>\frac{2328129}{25000000},\\
 B_0=Q=0:&\qquad
 c\mathcal S>\frac{12238141}{100000000},\\
 B_0=0,\ Q=1,\ W=\frac34:&\qquad
 c\mathcal S>\frac{8808129}{25000000},\\
 B_0=0,\ Q=1,\ W=\frac34-c:&\qquad
 c\mathcal S>\frac{2143251}{6250000}.
 \end{aligned}
 \label{eq:high-floor-round-eight-f-negative-ledgers}
\end{equation}
This proves positivity on every $\delta\geq1$ face.

It remains to treat $0<\delta<1$.  Put
\[
 y=f-W=\delta+\frac14,
 \qquad L=\frac{149}{50}.
\]
We claim
\begin{equation}
 \boxed{T<\frac{Ly}{c}.}
 \label{eq:high-floor-round-eight-f-level-distance}
\end{equation}
The scale and concavity reductions used in
\eqref{eq:high-floor-round-eight-d-level-distance} reduce the claim to
$f=2$ and
\[
 x=\frac yW\in\left\{\frac17,\frac53\right\}.
\]
Set
\[
 \Phi=\sin\theta-\theta\cos\theta,\qquad
 t_0=\frac{Ly}{c},\qquad u_0=\frac{t_0}{\mu}.
\]
The bounds $\Phi<\theta^3/3$ and
\eqref{eq:high-floor-round-eight-f-ratio} give
\begin{equation}
 u_0<\frac{149}{50}\frac59
       \frac{311/2000}{1847/2000}
 =\frac{46339}{166230}<\frac{279}{1000}<1.
 \label{eq:high-floor-round-eight-f-u-zero}
\end{equation}
For the retained range, put
\[
 \kappa_\rho=\frac{\Phi}{\theta\rho(1-\rho)}.
\]
Alternating Taylor bounds give
\[
 \Phi>\theta^3\left(\frac13-\frac{\theta^2}{30}\right),
\]
\[
 1-\cos\theta<\theta^2\left(
 \frac12-\frac{\theta^2}{24}+\frac{\theta^4}{720}\right).
\]
The quotient
\[
 t\longmapsto
 \frac{1/3-t/30}{(927/1000)(1/2-t/24+t^2/720)}
\]
is decreasing because its derivative has the sign of
$t^2-20t-60<0$.  At $t=311/2000$ its gap above $7/10$ is
\[
 \frac{25224124159}{1464079822630}>0.
\]
Consequently
\begin{equation}
 \kappa_\rho>\frac7{10},\qquad
 \frac{u_0}{1-\rho}>\frac{1043}{500}x.
 \label{eq:high-floor-round-eight-f-retained-range}
\end{equation}

Use
\[
 \arccos(1-z)=\sqrt{2z}\,S(z).
\]
On the retained subtracted range, with $V=(1043/500)x$,
\[
 V\leq\frac{1043}{300},\qquad
 z\leq\frac{53193}{200000}.
\]
The decreasing positive coefficients of $S$ give
\begin{equation}
 1+\frac z{12}+\frac{3z^2}{160}
 \leq S(z)\leq
 1+\frac z{12}+\frac{3750z^2}{146807}.
 \label{eq:high-floor-round-eight-f-series}
\end{equation}
The integrated upper coefficient is $7500/1027649$, while the positive
first copy is used only for $z\leq68493/200000<1$.  The decreasing
coefficient quotient used in Round~8E now yields
\begin{equation}
 C_\rho=
 \frac{\rho\sqrt2(1-\rho)^{3/2}}
 {\sin\theta-\theta\cos\theta}>\frac{689}{500}.
 \label{eq:high-floor-round-eight-f-coefficient}
\end{equation}
The retained lower integral is
\begin{equation}
\begin{aligned}
 \mathcal I^{\mathrm F}_\rho(V)
={}&\frac{2}{3\rho}
 \left((1+\rho V)^{3/2}-1\right)
 +\frac{1-\rho}{30\rho}
 \left((1+\rho V)^{5/2}-1\right)\\
&+\frac{3(1-\rho)^2}{560\rho}
 \left((1+\rho V)^{7/2}-1\right)
 -\frac23V^{3/2}-\frac{1-\rho}{30}V^{5/2}\\
&-\frac{7500(1-\rho)^2}{1027649}V^{7/2}.
\end{aligned}
 \label{eq:high-floor-round-eight-f-integral}
\end{equation}
On $64$ equal rational panels covering
$[1847/2000,927/1000]$, directed evaluation at both endpoint values of
$x$ proves
\begin{equation}
 \mathcal I^{\mathrm F}_\rho(V)>0,\qquad
 \frac{689}{500}\mathcal I^{\mathrm F}_\rho(V)-x>0.
 \label{eq:high-floor-round-eight-f-panels}
\end{equation}
Across $384$, $512$, $768$, and $1024$ bits, the smallest directed lower
bounds exceed $0.2113764065$ and $0.0357944486$, respectively.
Concavity proves \eqref{eq:high-floor-round-eight-f-level-distance}.

The endpoint chord, $M\geq\delta/c$, and
\eqref{eq:high-floor-round-eight-f-level-distance} give
\[
 cP>\left(\frac{3-L}{4}-c\right)\delta
      -\frac L{16}-\frac{cd}{2}.
\]
The coefficient of $\delta$ is negative, so $\delta<1$ yields
\begin{equation}
 cP>\frac{3-L}{4}-c-\frac L{16}-\frac{cd}{2}.
 \label{eq:high-floor-round-eight-f-small-prefix}
\end{equation}
Here $B_0=f-1$ and $Q\leq f-1$.  The exact-angle reserve and
\eqref{eq:high-floor-round-eight-f-cap} therefore give
\[
 cD_A(q)>\frac{f-1}{2}-c(f-1)-\frac c6
 \geq\frac12-\frac{7c}{6}.
\]
Adding this to \eqref{eq:high-floor-round-eight-f-small-prefix}, using
$d=1-2c$ and $L=149/50$, gives
\begin{equation}
 c\mathcal S>
 \frac{51}{160}-\frac83c+c^2
 \geq\frac{1879}{25000000}>0.
 \label{eq:high-floor-round-eight-f-small-delta}
\end{equation}
Together with \eqref{eq:high-floor-round-eight-f-positive-ledgers} and
\eqref{eq:high-floor-round-eight-f-negative-ledgers}, this proves
\eqref{eq:high-floor-round-eight-f-positive}.

For a residual ratio $\rho<1847/2000$, monotonicity and directed endpoint
evaluation give
\[
 \eta_\rho>\frac1{158},\qquad \eta_\rho<\frac19,
\]
and
\[
 a_\rho<2\frac{355}{113}\frac{1847}{153}
 =\frac{1311370}{17289}.
\]
Since $C_*<13/10$,
\[
 a_\rho+2\eta_\rho C_*
 <\frac{1311370}{17289}+\frac{13}{45}
 =\frac{2193941}{28815}<77.
\]
As $C_*>1/2$, the square root in the fixed-ratio threshold is smaller
than the same left side.  Hence
\begin{equation}
 \boxed{K_0(\rho)<77\mathbin{\cdot}158^2=1922228.}
 \label{eq:high-floor-round-eight-f-cutoff}
\end{equation}
The integer arguments following
\eqref{eq:high-floor-round-eight-d-cutoff} give
\eqref{eq:high-floor-round-eight-f-continuous}--%
\eqref{eq:high-floor-round-eight-f-discrete} and the chamber conditions.

All exact and directed comparisons in this proof are replayed by
\begin{center}
 \small\path{scripts/general_d_high_floor_signed_prefix_verify.py}.
\end{center}
The certificate pins the Round~8E verifier hash and covers all $128$
panel-endpoint evaluations at four precisions.
\end{proof}

\begin{theorem}[Correlated-prefix high-floor exclusion]
\label{thm:high-floor-round-eight-g-scalar}
In the high-floor first-drop branch, suppose
\[
 f=F_0\geq2,\qquad p<n,\qquad
 \frac{23}{25}\leq\rho<1.
\]
Then
\begin{equation}
 \boxed{
 \mathcal S=D_A(q)+R_p+\frac{d_\rho}{2}(n-p)>0.}
 \label{eq:high-floor-round-eight-g-positive}
\end{equation}
Consequently every remaining negative high-floor first-drop candidate
satisfies
\begin{equation}
 \boxed{
 \frac1{3154914}<\rho<\frac{23}{25},\qquad
 \frac{21}{4}<K<1577457,}
 \label{eq:high-floor-round-eight-g-continuous}
\end{equation}
and
\begin{equation}
 \boxed{
 \begin{gathered}
 r\in\tfrac12\N,\qquad 1\leq2r\leq3154913,\\
 2\leq n\leq1577456,\qquad 1\leq p\leq n-1,\\
 2\leq f\leq525819,\\
 0\leq Q\leq f-1,\qquad Q\leq B\leq525819.
 \end{gathered}}
 \label{eq:high-floor-round-eight-g-discrete}
\end{equation}
The exact chamber conditions remain
\begin{equation}
 \mu=\rho K,\qquad q=r+n\leq\mu<q+1,\qquad
 K-\mu>\frac{7\pi}{4},
 \label{eq:high-floor-round-eight-g-chamber-one}
\end{equation}
\begin{equation}
 \left\lfloor A(r+j)+\frac14\right\rfloor=f
 \quad(0\leq j\leq p),\qquad
 A(r+p+1)<f-\frac14.
 \label{eq:high-floor-round-eight-g-chamber-two}
\end{equation}
This theorem supersedes the residual box in
Theorem~\ref{thm:high-floor-round-eight-f-scalar}.  It certifies the
displayed ratio interval; it does not certify any analytic wall in the
residual box \eqref{eq:high-floor-round-eight-g-continuous}.
\end{theorem}

\begin{proof}
By Theorem~\ref{thm:high-floor-round-eight-f-scalar}, it is enough to
treat the correlated band
\begin{equation}
 \frac{23}{25}\leq\rho\leq\frac{1847}{2000}.
 \label{eq:high-floor-round-eight-g-band}
\end{equation}
As before, $p=0$ is automatic by
Theorem~\ref{thm:one-interface-complete}.  Hence a negative candidate has
$p\geq1$, $n\geq2$, and $q\geq5/2$.  Strict convexity of $G_\mu$, the
chamber $q\leq\mu<q+1$, and monotonicity of $G_{q+1}(q)$ give the strict
terminal cap
\begin{equation}
 \boxed{2\int_q^\mu G_\mu(z)\dd z
 <G_{7/2}(5/2)<\frac16.}
 \label{eq:high-floor-round-eight-g-cap}
\end{equation}

Put $\theta=\arccos\rho$ and retain
\[
 c=\frac{\theta}{\pi},\qquad d=1-2c,\qquad
 \widehat H(t)=A(\mu-t),\qquad
 W=\widehat H(0)=G_K(\mu).
\]
Write
\[
 h=f-\frac14,\qquad \delta=h-W,\qquad
 P=R_p+\frac d2(n-p),\qquad
 \beta=\frac{123}{1000}-c.
\]
The first shelf determines the point $M$ by $\widehat H(M)=h$.
If level $f$ is present, let $\widehat H(T)=f$; if it is absent, put
$T=\mu$.  These are the same wall-safe shelf conventions used in the
preceding rounds.

Suppose first that $\delta\geq1$.  On the absent-level face the refined
chord from Theorem~\ref{thm:high-floor-round-eight-e-scalar} gives
\[
 P>\beta M-\frac d2.
\]
On the present-level face, the elasticity estimate gives
\[
 P\geq\left(\frac18-c\right)M-\frac d2
 >\beta M-\frac d2,
\]
because the coefficient gap is $1/500$ and $M>0$.  Thus both faces have
the strict common prefix
\begin{equation}
 \boxed{P>\beta M-\frac d2.}
 \label{eq:high-floor-round-eight-g-common-prefix}
\end{equation}

On \eqref{eq:high-floor-round-eight-g-band}, directed endpoint
evaluation gives
\begin{equation}
 \theta^2<\frac{163}{1000},\qquad
 c<\bar c:=\frac{641}{5000}.
 \label{eq:high-floor-round-eight-g-ratio}
\end{equation}
At $\rho=1847/2000$ one has $c>123/1000$, and $c$ increases as
$\rho$ decreases.  Hence $\beta<0$ throughout the band.  Define
\begin{equation}
 \Phi=\sin\theta-\theta\cos\theta,\qquad
 R_\rho=\frac{\theta\rho}{\Phi},\qquad
 \gamma(\theta)=\beta R_\rho.
 \label{eq:high-floor-round-eight-g-gamma}
\end{equation}
Since $M\leq\mu$, multiplication by the negative coefficient reverses
the comparison in the substitution, and
\begin{equation}
 cP>\beta c\mu-\frac{cd}{2}
 =\boxed{\gamma(\theta)W-\frac{cd}{2}.}
 \label{eq:high-floor-round-eight-g-correlated-prefix}
\end{equation}

For
\[
 R(\theta)=\frac{\theta\cos\theta}{\Phi},
\]
direct differentiation gives
\[
 R'(\theta)=
 \frac{\sin\theta\cos\theta-\theta}{\Phi^2}<0
\]
and
\begin{equation}
 \gamma'(\theta)=-\frac{R(\theta)}{\pi}
 +\left(\frac{123}{1000}-\frac{\theta}{\pi}\right)R'(\theta).
 \label{eq:high-floor-round-eight-g-gamma-derivative}
\end{equation}
Directed Arb evaluation on $64$ complete rational $\rho$-panels proves
\begin{equation}
 \boxed{\gamma'(\theta)<-\frac{49}{10}<0.}
 \label{eq:high-floor-round-eight-g-gamma-wall}
\end{equation}
Across the four replay precisions, the directed enclosure union was
\[
 -5.595800902<\gamma'(\theta)<-4.986279971.
\]
Thus $\gamma$ decreases with $\theta$, and its minimum is at
$\rho=23/25$.  Directed endpoint evaluation there gives
\[
 \gamma(\arccos(23/25))
 =-0.0897444733642\ldots>-\frac9{100},
\]
so
\begin{equation}
 \boxed{-\frac9{100}<\gamma<0.}
 \label{eq:high-floor-round-eight-g-gamma-range}
\end{equation}

The exact active-width factor
\[
 \lambda(\theta)=\frac{\Phi}{1-\cos\theta}
\]
satisfies
\[
 \lambda'(\theta)=
 \frac{\sin\theta(\theta-\sin\theta)}
 {(1-\cos\theta)^2}>0.
\]
Its minimum is at $\rho=1847/2000$, where directed evaluation gives
$\lambda>229/875$.  Since the active relation gives
$W>\lambda(W+\delta)$ and $W+\delta=f-1/4\geq7/4$,
\begin{equation}
 \boxed{W>\frac{229}{500}.}
 \label{eq:high-floor-round-eight-g-width}
\end{equation}

Introduce the strict terminal counts
\[
 B_0=\sbr{W+\tfrac14},\qquad
 B=\sbr{G_K(q)+\tfrac14},\qquad
 Q=\sbr{A(q)+\tfrac14}.
\]
One has $B\geq B_0$ and $Q\leq B_0+1$, so the following three strata
are exhaustive.

If $B_0\geq1$, strict bracketing gives $W\leq B_0+3/4$.
Because $\gamma<0$, the cap
\eqref{eq:high-floor-round-eight-g-cap}, the correlated prefix
\eqref{eq:high-floor-round-eight-g-correlated-prefix}, and the
exact-angle terminal reserve yield
\begin{equation}
 c\mathcal S>
 B_0\left(\frac12-c+\gamma\right)
 +\frac34\gamma-\frac{7c}{6}-\frac{cd}{2}.
 \label{eq:high-floor-round-eight-g-B-ledger}
\end{equation}
The coefficient is positive:
\[
 \frac12-c+\gamma>
 \frac12-\frac{641}{5000}-\frac9{100}
 =\frac{1409}{5000}>0.
\]
Thus $B_0=1$ is the lower endpoint.  The resulting ledger
\[
 L_B(c,\gamma)=\frac12-\frac83c+c^2+\frac74\gamma
\]
decreases in $c$ and increases in $\gamma$.  Hence
\begin{equation}
 \boxed{c\mathcal S>
 \frac{1280143}{75000000}>0.}
 \label{eq:high-floor-round-eight-g-B-positive}
\end{equation}

If $B_0=Q=0$, the outer tangent triangle and
\eqref{eq:high-floor-round-eight-g-correlated-prefix} give
\[
 c\mathcal S>
 F_0(W):=W^2+\gamma W-\frac c2+c^2.
\]
By \eqref{eq:high-floor-round-eight-g-gamma-range} and
\eqref{eq:high-floor-round-eight-g-width},
\[
 F_0'(W)=2W+\gamma>
 2\frac{229}{500}-\frac9{100}=\frac{413}{500}>0.
\]
At $W=229/500$ the ledger decreases in $c$ and increases in $\gamma$,
so
\begin{equation}
 \boxed{c\mathcal S>
 \frac{3021981}{25000000}>0.}
 \label{eq:high-floor-round-eight-g-Q-zero}
\end{equation}

Finally suppose $B_0=0$ and $Q=1$.  Then
\[
 \frac34-c<W\leq\frac34,
\]
and the inner payment plus the outer tangent triangle gives
\[
 c\mathcal S>
 F_1(W):=-W^2+\left(\frac32+\gamma\right)W
 -\frac32c+c^2.
\]
This function is concave in $W$, so only the two endpoints are needed.
Both endpoint expressions decrease in $c$ and increase in $\gamma$ on
the certified rectangle.  Therefore
\begin{equation}
 \boxed{
 F_1(3/4)>\frac{7978381}{25000000}>0,}
 \label{eq:high-floor-round-eight-g-Q-one-upper}
\end{equation}
\begin{equation}
 \boxed{
 F_1(3/4-c)>\frac{157119}{500000}>0.}
 \label{eq:high-floor-round-eight-g-Q-one-lower}
\end{equation}
For an independent correlated replay, $256$ finer complete panels give
the respective directed lower bounds
\[
 0.017263564377,\qquad 0.120924743146,\qquad
 0.319220429717,\qquad 0.314297954935.
\]
The corresponding $B_0$-coefficient and $Q=0$ derivative bounds are,
respectively,
\[
 0.281907873981\qquad\hbox{and}\qquad0.826096309901.
\]
Thus
\eqref{eq:high-floor-round-eight-g-B-positive}--%
\eqref{eq:high-floor-round-eight-g-Q-one-lower} close every
$\delta\geq1$ face, including the strict-count walls.

It remains to treat $0<\delta<1$.  Put
\[
 y=f-W=\delta+\frac14,\qquad L=\frac{29}{10}.
\]
The inherited scale monotonicity reduces the largest level distance at
fixed $y$ to $f=2$, $W=2-y$.  With $x=y/W$, the inherited concavity
reduction leaves only
\begin{equation}
 x\in\left\{\frac17,\frac53\right\}.
 \label{eq:high-floor-round-eight-g-x-endpoints}
\end{equation}
Let
\[
 t_0=\frac{Ly}{c},\qquad u_0=\frac{t_0}{\mu},\qquad
 \kappa_\rho=\frac{\Phi}{\theta\rho(1-\rho)}.
\]
The bound $\Phi<\theta^3/3$, together with
\eqref{eq:high-floor-round-eight-g-ratio}, gives
\begin{equation}
 u_0<
 \frac{29}{10}\frac59\frac{163/1000}{23/25}
 =\frac{4727}{16560}<\frac{143}{500}<1.
 \label{eq:high-floor-round-eight-g-u-zero}
\end{equation}

For the sharper retained range, write
\[
 \kappa(\theta)=
 \frac{\sin\theta-\theta\cos\theta}
 {\theta\cos\theta(1-\cos\theta)}.
\]
If $N=\Phi$ and $D=\theta\cos\theta(1-\cos\theta)$, then
\begin{equation}
 N'D-ND'=
 \theta\sin\theta\,\theta\rho(1-\rho)
 -\Phi\left(\rho(1-\rho)
 +\theta\sin\theta(2\rho-1)\right).
 \label{eq:high-floor-round-eight-g-kappa-numerator}
\end{equation}
Directed evaluation on $64$ complete rational panels gives
\[
 N'D-ND'>0.000190444285561.
\]
Thus $\kappa$ increases with $\theta$, and its minimum is at
$\rho=1847/2000$.  At that endpoint
\begin{equation}
 \boxed{\kappa_\rho>\frac{18}{25},}
 \qquad
 \kappa_{1847/2000}-\frac{18}{25}
 >1.6195058963\times10^{-5}.
 \label{eq:high-floor-round-eight-g-kappa}
\end{equation}
Consequently
\begin{equation}
 \frac{u_0}{1-\rho}>
 L\frac{18}{25}x=\frac{261}{125}x.
 \label{eq:high-floor-round-eight-g-retained-range}
\end{equation}

Use
\[
 \arccos(1-z)=\sqrt{2z}\,S(z).
\]
The alternating Taylor estimates
\[
 1-\cos\theta>
 \theta^2\left(\frac12-\frac{\theta^2}{24}\right),
\]
\[
 \Phi<\theta^3\left(
 \frac13-\frac{\theta^2}{30}+\frac{\theta^4}{840}\right)
\]
and the decreasing coefficient quotient give
\begin{equation}
 \boxed{
 C_\rho=\frac{\rho\sqrt2(1-\rho)^{3/2}}{\Phi}
 >\frac{137}{100}.}
 \label{eq:high-floor-round-eight-g-coefficient}
\end{equation}
Indeed, the logarithmic derivative test reduces to
\[
 84+10t-\frac{t^2}{2}>0,\qquad
 0\leq t\leq\frac{163}{1000},
\]
and evaluation at $\rho=23/25$, $t=163/1000$ gives the stronger
directed lower value $1.37424804127\ldots$.

Retain
\[
 0\leq v\leq V:=\frac{261}{125}x\leq\frac{87}{25}.
\]
Since $1-\rho\leq2/25$,
\begin{equation}
 (1-\rho)V\leq\frac{174}{625},\qquad
 (1-\rho)(1+V)\leq\frac{224}{625}<1.
 \label{eq:high-floor-round-eight-g-series-domains}
\end{equation}
The positive decreasing coefficients of $S$ give
\begin{equation}
 1+\frac z{12}+\frac{3z^2}{160}
 \leq S(z)\leq
 1+\frac z{12}+\frac{21z^2}{1000}.
 \label{eq:high-floor-round-eight-g-series}
\end{equation}
Indeed, on $z\leq174/625$ the required upper quadratic coefficient is
\[
 \frac3{160}+\frac5{896}
 \frac{174/625}{1-174/625}
 =\frac{21117}{1010240}<\frac{21}{1000},
\]
whose convenient integrated ceiling is
\[
 \frac{21}{1000}\frac27=\frac3{500}.
\]
The resulting retained lower integral is
\begin{equation}
\begin{aligned}
 \mathcal I^{\mathrm G}_\rho(V)
={}&\frac{2}{3\rho}
 \left((1+\rho V)^{3/2}-1\right)
 +\frac{1-\rho}{30\rho}
 \left((1+\rho V)^{5/2}-1\right)\\
&+\frac{3(1-\rho)^2}{560\rho}
 \left((1+\rho V)^{7/2}-1\right)
 -\frac23V^{3/2}-\frac{1-\rho}{30}V^{5/2}\\
&-\frac{3(1-\rho)^2}{500}V^{7/2}.
\end{aligned}
 \label{eq:high-floor-round-eight-g-integral}
\end{equation}
On $64$ equal rational panels covering
\eqref{eq:high-floor-round-eight-g-band}, directed Arb evaluation at
both endpoint values in
\eqref{eq:high-floor-round-eight-g-x-endpoints} gives $128$ checks and
proves
\begin{equation}
 \mathcal I^{\mathrm G}_\rho(V)>0,\qquad
 \frac{137}{100}\mathcal I^{\mathrm G}_\rho(V)-x>0.
 \label{eq:high-floor-round-eight-g-integral-panels}
\end{equation}
Across $384$, $512$, $768$, and $1024$ bits, the smallest directed
lower bounds were
\[
 \mathcal I^{\mathrm G}_\rho(V)>0.211576376293,
\qquad
 \frac{137}{100}\mathcal I^{\mathrm G}_\rho(V)-x
 >0.0214999332585.
\]
Concavity fills the full $x$-interval.  Hence level $f$ is present and
\begin{equation}
 \boxed{T<\frac{29}{10}\frac yc.}
 \label{eq:high-floor-round-eight-g-level-distance}
\end{equation}

The endpoint chord, $M\geq\delta/c$, and
\eqref{eq:high-floor-round-eight-g-level-distance} now give
\[
 cP>
 \left(\frac{3-L}{4}-c\right)\delta
 -\frac L{16}-\frac{cd}{2}.
\]
Here $(3-L)/4-c<0$, so $\delta<1$ puts the lower endpoint at
$\delta=1$.  Moreover, $B_0=f-1$, $Q\leq f-1$, and
\eqref{eq:high-floor-round-eight-g-cap} gives
\[
 cD_A(q)>
 \frac{f-1}{2}-c(f-1)-\frac c6
 \geq\frac12-\frac{7c}{6}.
\]
Adding the two bounds and using $d=1-2c$ and $L=29/10$ yields
\[
 c\mathcal S>
 \frac54-\frac{5L}{16}-\frac83c+c^2.
\]
This decreases on the certified interval; at $c=641/5000$,
\begin{equation}
 \boxed{
 c\mathcal S>\frac{1373893}{75000000}>0.}
 \label{eq:high-floor-round-eight-g-small-delta}
\end{equation}
Together with the three terminal strata, this proves
\eqref{eq:high-floor-round-eight-g-positive}.

It remains to sharpen the finite residual box.  For $\rho<23/25$, put
\[
 \eta_\rho=G_1(\max\{\rho,1/2\}),\qquad
 a_\rho=\frac{2\pi\rho}{1-\rho}.
\]
Monotonicity and directed endpoint evaluation give
\begin{equation}
 \eta_\rho>\frac1{147},\qquad \eta_\rho<\frac19,
 \label{eq:high-floor-round-eight-g-eta}
\end{equation}
while
\[
 a_\rho<23\pi<23\frac{355}{113}=\frac{8165}{113}.
\]
For the fixed-ratio constant $C_*$ used in
Proposition~\ref{prop:fixed-rho-high-energy}, the inherited walls
$1/2<C_*<13/10$ give
\[
 a_\rho+2\eta_\rho C_*
 <\frac{8165}{113}+\frac{13}{45}
 =\frac{368894}{5085}<73.
\]
The lower wall $C_*>1/2$ makes the square-root term in the
fixed-ratio threshold strictly smaller than the same left side.
Therefore
\begin{equation}
 \boxed{K_0(\rho)<73\mathbin{\cdot}147^2=1577457.}
 \label{eq:high-floor-round-eight-g-cutoff}
\end{equation}

Negativity excludes $p=0$, so $p\geq1$, $n\geq2$, and
$r<\mu<K$.  Since $r\geq1/2$,
\[
 \rho>\frac1{3154914},\qquad
 2r\leq3154913,\qquad
 n\leq1577456.
\]
Finally $A(r)\geq7/4$ gives $K-\mu>7\pi/4$ and $K>21/4$.
Since $\pi>3$,
\[
 f,B<\frac K3+\frac14<525819+\frac14.
\]
Together with $1\leq p\leq n-1$, $Q\leq f-1$, and $Q\leq B$, these
are precisely \eqref{eq:high-floor-round-eight-g-continuous}--%
\eqref{eq:high-floor-round-eight-g-chamber-two}.

All directed comparisons are replayed by
\begin{center}
 \small\path{scripts/general_d_round_08g_correlated_prefix_verify.py}.
\end{center}
The verifier uses $64$ full panels for
\eqref{eq:high-floor-round-eight-g-gamma-wall}, $256$ finer panels for
the correlated terminal ledgers, $64$ full panels for
\eqref{eq:high-floor-round-eight-g-kappa-numerator}, and $128$
panel-endpoint checks for
\eqref{eq:high-floor-round-eight-g-integral-panels}.  It pins both
frozen Round~8F artifacts and uses exact rational panel endpoints,
outward-rounded Arb enclosures, and exact rational arithmetic for the
remaining ledgers and cutoff.  The resulting finite box is only a
reduction: no wall in
\eqref{eq:high-floor-round-eight-g-continuous} is claimed certified.
\end{proof}

\subsection{Global finite exhaustion of the no-drop branch}

We first record the root-free consequence of the terminal estimate that
drives the exhaustion.  For $J\geq1$, put
\begin{equation}
 C_0^3=\frac\pi{12},\qquad c_*=\frac{4\sqrt2}{15},
 \qquad S_J=\sum_{k=1}^{J}(k-\tfrac14)^{-1/3}.
 \label{eq:no-drop-cutoff-constants}
\end{equation}

\begin{proposition}[Root-free no-drop cutoff]
\label{prop:root-free-no-drop-cutoff}
In the no-drop branch of Proposition~\ref{prop:no-drop-identities}, let
\[
 Q=\sbr{A(q)+\tfrac14}\in\{f-1,f\},\qquad
 \mathcal S_n=D_A(q)+R_n.
\]
Then
\begin{equation}
 D_A(q)\geq
 \max\{0,C_0K^{1/3}S_Q-Q-c_*\}.
 \label{eq:root-free-terminal-reserve}
\end{equation}
Consequently the no-drop scalar is nonnegative whenever
\begin{equation}
 K\geq K_{\rm nd}(n,Q):=
 \left(\frac{Q+n/2+c_*}{C_0S_Q}\right)^3.
 \label{eq:no-drop-cutoff}
\end{equation}
Every branch below this cutoff is confined, for fixed $(n,f,Q)$, to a
compact strip with finitely many action walls.
\end{proposition}

\begin{proof}
For $0\leq\theta\leq\pi/2$,
\[
 \sin\theta-\theta\cos\theta
 =\int_0^\theta t\sin t\dd t\geq\frac{2\theta^3}{3\pi}.
\]
Thus the angles in Corollary~\ref{cor:exact-angle-terminal-reserve} satisfy
\[
 \frac\pi{2\theta_k}\geq
 C_0K^{1/3}(k-\tfrac14)^{-1/3}.
\]
The cap bound \eqref{eq:cap-bound} is smaller than $c_*$, and the
fractional-top term is nonnegative.  Since the ball level count $B$ is at
least $Q$, this proves \eqref{eq:root-free-terminal-reserve}.  Under
\eqref{eq:no-drop-cutoff} it gives $D_A(q)\geq n/2$, whereas
\eqref{eq:no-drop-head-exact} gives $R_n\geq-n/2$.

Finally, $K-\mu>\pi(f-1/4)$ and $\mu\geq r+n$, so
\[
 K>\pi(f-\tfrac14)+r+n.
\]
Together with $K<K_{\rm nd}(n,Q)$ this bounds $r,K,\mu$, every terminal
level, and every sample index.  On each open chamber the exact scalar has
\[
 \partial_K\mathcal S_n=\frac2\pi\int_r^K
 \sqrt{1-z^2/K^2}\dd z>0,\qquad
 \partial_\mu\mathcal S_n=-\frac2\pi\int_r^\mu
 \sqrt{1-z^2/\mu^2}\dd z<0.
\]
It therefore reduces to the lower-$K$, upper-$\mu$, and finitely many
one-sided sample walls; the strict equality face is retained rather than
passed through by continuity.
\end{proof}

\begin{proposition}[Uniform half-unit terminal reserve]
\label{prop:no-drop-half-unit}
Let $q\leq\mu<q+1$, $q\geq3/2$, and $A(q)\geq7/4$.  Then, including at
the lower ordinary-floor wall at $q$,
\begin{equation}
 \boxed{D_A(q)>\frac12.}
 \label{eq:no-drop-half-unit}
\end{equation}
\end{proposition}

\begin{proof}
Put
\[
 v=G_K(q),\qquad h=G_\mu(q),\qquad
 f=\left\lfloor A(q)+\frac14\right\rfloor\geq2,
 \quad Q=\sbr{A(q)+\tfrac14},\quad B=\sbr{v+\tfrac14}.
\]
If $q=\mu$, the inner cap vanishes.  If $q<\mu$, strict convexity of
$G_\mu$ on $[q,\mu]$ and $\mu-q<1$ give
\[
 2\int_q^\mu G_\mu(z)\dd z<h(\mu-q)<h.
\]
Moreover $h\leq G_{q+1}(q)$, and $G_{q+1}(q)$ decreases with $q$.  Outward
arithmetic therefore gives the uniform cap bound
\begin{equation}
 2\int_q^\mu G_\mu(z)\dd z
 <G_{5/2}(3/2)<\frac15.
 \label{eq:no-drop-small-cap}
\end{equation}
For $1\leq k\leq B$, let $\theta_k$ be the exact ball angle at level
$k-1/4$.  Corollary~\ref{cor:exact-angle-terminal-reserve}, with its
nonnegative fractional-top term omitted, gives
\begin{equation}
 D_A(q)\geq\frac\pi2\sum_{k=1}^{B}\frac1{\theta_k}
 -Q-2\int_q^\mu G_\mu(z)\dd z.
 \label{eq:no-drop-half-unit-angle}
\end{equation}

If $q\geq K/2$, every retained level crosses strictly to the right of
$q$, so $\theta_k<\arccos(q/K)\leq\pi/3$.  Away from the lower wall,
$Q=f$ and $B\geq f$, which gives more than $f/2-1/5$.  At the lower wall
with $q<\mu$, one has $Q=f-1$ but $B\geq f$; with $q=\mu$, the cap is zero
and $B=Q=f-1$, so pairing all $f-1$ count units with their angle terms
gives $D_A(q)>(f-1)/2\geq1/2$.

If $q<K/2$, let $K_*$ be the unique solution of
\[
 G_{K_*}(3/2)=\frac74.
\]
Monotonicity gives $K\geq K_*$.  At that radius, a directed 512-bit Arb
calculation proves
\[
 \theta_1<\frac{503}{500},\qquad
 \theta_2<\frac{11}{8},
\]
and hence
\begin{equation}
 \frac\pi2\left(\frac1{\theta_1}+\frac1{\theta_2}\right)-2
 >\frac7{10},\qquad
 \frac\pi{2\theta_1}-1>\frac12.
 \label{eq:no-drop-two-angle-margin}
\end{equation}
The angles decrease as $K$ increases, and every later net level is
positive.  Away from the lower wall, the first inequality in
\eqref{eq:no-drop-two-angle-margin} and
\eqref{eq:no-drop-small-cap} give more than $1/2$.  At the lower wall the
same two cases $q<\mu$ and $q=\mu$ are handled, respectively, by the extra
ball level and the second inequality in
\eqref{eq:no-drop-two-angle-margin}.  The computation is reproduced by
\path{scripts/general_d_no_drop_round5_verify.py} and passes unchanged at
1024 bits.
\end{proof}

\begin{corollary}[Closure of the one-cell no-drop branch]
\label{cor:no-drop-n-one}
Every no-drop branch with $f\geq2$ and $n=1$ satisfies
\begin{equation}
 \boxed{\mathcal S_1=D_A(q)+R_1>0.}
 \label{eq:no-drop-n-one}
\end{equation}
\end{corollary}

\begin{proof}
Concavity of $A$ on $[r,q]$ gives
\[
 2\int_r^qA(z)\dd z\geq A(r)+A(q)\geq2f-\frac12.
\]
Thus $R_1\geq-1/2$, while
Proposition~\ref{prop:no-drop-half-unit} is strict.
\end{proof}

\begin{corollary}[A finite-chamber discard sector]
\label{cor:no-drop-finite-discard}
For a no-drop branch with $f\geq2$, put
\[
 \varepsilon=A(q)+\frac14-f,\qquad
 \Delta=A(r)-A(q).
\]
If
\begin{equation}
 2\varepsilon+\Delta+
 \frac{n\Delta}{3(2r+n)}\geq\frac12-\frac1{2n},
 \label{eq:no-drop-finite-discard}
\end{equation}
then $\mathcal S_n>0$.
\end{corollary}

\begin{proof}
The scale-free shelf-curvature bound gives
\[
 R_n\geq\frac{n^2\Delta}{3(2r+n)}
 +n\left(2\varepsilon+\Delta-\frac12\right)\geq-\frac12.
\]
Apply Proposition~\ref{prop:no-drop-half-unit}.
\end{proof}

\begin{theorem}[Quantitative no-drop transfer]
\label{thm:quantitative-no-drop-transfer}
Let $f\in\{2,3\}$, $b=f-1/4$, and
\[
 c_f=\frac{\sqrt{1+b^2}-1}{2},\qquad
 C_K(f)=\left(\frac{f+1/2+c_*}{C_0S_{f-1}}\right)^3,\qquad
 C_\delta(f)=\left(8\pi^2f^2C_K(f)\right)^{1/3}.
\]
Suppose a residual no-drop candidate satisfies, with $\delta=K-\mu$,
\begin{equation}
 \frac{c_f}{2}n\leq\delta\leq C_\delta(f)n,\qquad
 \frac{c_f^2}{8\pi^2(C_\delta(f)+2)}n^3
 <K<C_K(f)n^3.
 \label{eq:quantitative-no-drop-cone}
\end{equation}
Put
\[
 s=\sqrt{\delta/K},\qquad \kappa=s\delta=s^3K,\qquad N=sn.
\]
If $0<s\leq1/10$, then, including at every endpoint wall,
\begin{equation}
 \boxed{s\mathcal S_n=s\bigl(D_A(q)+R_n\bigr)>\frac1{10}.}
 \label{eq:quantitative-no-drop-gap}
\end{equation}
\end{theorem}

\begin{proof}
Put $\alpha=\mu-q\in[0,1)$ and $X_q=s\alpha$.  The residual endpoint
floors give
\begin{equation}
 b\leq A(q)<f\leq3,
 \qquad \frac{21}{8}\leq\kappa.
 \label{eq:round-five-endpoint-box}
\end{equation}
The last inequality follows from the radius representation and
$A(q)\leq2\kappa/\pi$.

Use the exact head coordinate $X=s(\mu-z)$ and set
\[
 \mathcal H_{s,\kappa}(X)=A(\mu-X/s),\qquad
 H_\kappa(X)=\frac{c_0}{\sqrt\kappa}
 \left((\kappa+X)^{3/2}-X^{3/2}\right),
 \quad c_0=\frac{2\sqrt2}{3\pi}.
\]
Unlike a two-ball subtraction, the shell has the non-cancelling exact
representation
\begin{equation}
 \mathcal H_{s,\kappa}(X)=\frac1\pi\int_0^\kappa
 \frac{\sqrt{(\kappa+X-u)
 [2\kappa-(\kappa+X+u)s^2]}}
 {\kappa-us^2}\dd u.
 \label{eq:round-five-shell-integral}
\end{equation}
The ratio of its integrand to the integrand defining $H_\kappa$ is
\[
 \mathcal R(u,X)=
 \frac{\sqrt{1-a(u,X)s^2}}{1-b_0(u)s^2},\qquad
 a(u,X)=\frac{\kappa+X+u}{2\kappa},\quad
 b_0(u)=\frac u\kappa.
\]
For $0\leq X\leq2$ at every physical point,
\[
 0\leq b_0\leq1,\qquad0\leq a\leq1+\frac1\kappa\leq\frac{29}{21}.
\]
The inequalities
\[
 1-y\leq\sqrt{1-y}\leq1,\qquad
 \frac1{1-s^2}-1\leq\frac{100}{99}s^2
\]
therefore give $|\mathcal R-1|\leq(29/21)s^2$.  At $X_q<s$, the sharper
bound $a<107/105$ gives
\[
 H_\kappa(X_q)
 \leq\frac{\mathcal H_{s,\kappa}(X_q)}
 {1-(107/105)s^2}<\frac{31}{10}.
\]
Since
\[
 0<H_\kappa'(X)=\frac{3c_0}{2\sqrt\kappa}
 \bigl(\sqrt{\kappa+X}-\sqrt X\bigr)<\frac12,
\]
one has $H_\kappa<41/10$ on $[0,2]$.  Integrating the ratio bound against
the positive limiting integrand proves
\begin{equation}
 \left|\mathcal H_{s,\kappa}(X)-H_\kappa(X)\right|<6s^2
 \qquad(0\leq X\leq2).
 \label{eq:quantitative-turning-error}
\end{equation}

Put
\[
 \zeta=\frac s2+6s^2\leq\frac{11}{100},
 \qquad w=H_\kappa(0)=c_0\kappa.
\]
The endpoint estimate, \eqref{eq:quantitative-turning-error}, and
$H_\kappa'<1/2$ give
\begin{equation}
 w>b-\zeta.
 \label{eq:round-five-w-lower}
\end{equation}
For $X=\kappa u$, set
\[
 t=\sqrt{1+u}-\sqrt u,\qquad
 F(t)=\frac{3/t+t^3}{4}.
\]
Then
\[
 \frac{H_\kappa(X)}w=F(t),
 \qquad H_\kappa'(X)=\frac{3c_0}{2}t.
\]
Let $X_*$ be the point at which $H_\kappa(X_*)=f+\zeta$.  This point is
positive: otherwise \eqref{eq:quantitative-turning-error} at $X_q$ would
contradict $A(q)<f$.  Furthermore
\[
 \frac{f+\zeta}{w}
 <\frac{f+\zeta}{f-1/4-\zeta}
 \leq\frac{211}{164}<\frac{163}{125}=F(3/5).
\]
Because $F$ decreases, $t_*>3/5$, and hence
$H_\kappa'>9c_0/10$ on $[0,X_*]$.  With
$d=f+\zeta-w$, this gives
\[
 0<d<\frac14+2\zeta\leq\frac{47}{100},
 \qquad X_*<\frac{10d}{9c_0}<\frac{19}{10}.
\]
Thus $X_q+X_*<2$, so the error estimate applies on the whole possible
negative head.  Since $f+\zeta-H_\kappa$ is convex with endpoint values
$d$ and $0$, its integral is at most $dX_*/2$.  The exact increasing
profile and the no-drop head identity therefore give
\begin{equation}
 sR_n\geq-2\int_0^{X_*}
 (f+\zeta-H_\kappa(X))\dd X
 \geq-\frac{10d^2}{9c_0}.
 \label{eq:quantitative-head-transfer}
\end{equation}

The first complete terminal level $\tau=3/4$ is present because
$G_K(q)\geq A(q)\geq7/4$, including at the strict wall.  Its angle lies in
$(0,\pi/2)$.  If
\[
 x_\tau=\left(\frac{3\pi\tau}{K}\right)^{1/3},
\]
then $x_\tau^3=3(\sin\theta-\theta\cos\theta)$.  The chord bound for sine
and \eqref{eq:round-five-endpoint-box} imply $\theta<\sqrt3s$, while
\[
 \sin\theta-\theta\cos\theta
 \geq\frac{\theta^3}{3}-\frac{\theta^5}{30}
\]
gives $x_\tau/\theta>1-3s^2/10$.  The exact-angle reserve, the cap bound,
and $Q\leq3$ now give
\begin{equation}
 sD_A(q)\geq
 \frac\pi{2\sqrt2}
 \left(\frac{b-\zeta}{3/4}\right)^{1/3}
 \left(1-\frac{3s^2}{10}\right)-\frac{17}{5}s.
 \label{eq:quantitative-terminal-transfer}
\end{equation}
Here all further exact-angle and fractional-top terms were discarded in
the favorable direction; the count and cap cost is less than
$(3+2/5)s=17s/5$.

Finally, all constants close rationally:
\[
 \frac\pi{2\sqrt2}>1,\qquad c_0>\frac{49}{165},
 \qquad
 \frac{b-\zeta}{3/4}\geq\frac{164}{75}
 >\left(\frac{129}{100}\right)^3.
\]
Adding \eqref{eq:quantitative-head-transfer} and
\eqref{eq:quantitative-terminal-transfer} yields
\[
 s\mathcal S_n>
 \frac{129}{100}\frac{997}{1000}-\frac{17}{50}
 -\frac{10}{9}\frac{165}{49}\left(\frac{47}{100}\right)^2
 =\frac{1758611}{14700000}>\frac1{10}.
\]
At the lower ordinary-floor wall the separate identity
\eqref{eq:no-drop-wall-slack} supplies one additional favorable unit for
the original defect.  The outward-rounded constants used above are replayed
by \path{scripts/general_d_no_drop_round5_verify.py}; the certificate passes
at both 512 and 1024 bits.
\end{proof}

\begin{theorem}[Global finite exhaustion of no-drop data]
\label{thm:no-drop-global-exhaustion}
Suppose $f\geq2$, $n\geq1$, and $\mathcal S_n<0$.  Then
Proposition~\ref{prop:root-free-no-drop-cutoff} must fail, and the following
properties hold:
\begin{enumerate}
\item if $r\leq K/2$, necessarily
 \begin{equation}
 2\leq f\leq49,\qquad2\leq n\leq48;
 \label{eq:no-drop-central-box}
 \end{equation}
 \item if $r>K/2$ and $f\geq4$, necessarily
 \begin{equation}
 4\leq f\leq8,\qquad2\leq n\leq39;
 \label{eq:no-drop-outer-box}
 \end{equation}
 \item in the two exceptional rows $f=2,3$, necessarily
 $2\leq n\leq379$.  Before the quantitative transfer, an unbounded
 sequence could occur only in these rows, and then
 \begin{equation}
 K-\mu\asymp n,\qquad K\asymp n^3,\qquad
 r/K\longrightarrow1,\qquad\rho\longrightarrow1.
 \label{eq:no-drop-escaping-ray}
 \end{equation}
 Theorem~\ref{thm:quantitative-no-drop-transfer} gives the explicit
 conclusion $n<380$ in both rows.  Thus the residual triples
 $(n,f,Q)$ form a completely explicit finite set.
\end{enumerate}
The certified central list has $647$ necessary pairs after the two $n=1$
pairs are removed; the outer $f\geq4$ list has $60$ necessary pairs and is
unchanged.  These counts precede certification of the continuous action
chambers.
This is a finite reduction, not a proof of the residual finite wall
chambers.
\end{theorem}

\begin{proof}
Put
\[
 \delta=K-\mu,\quad \alpha=\mu-q\in[0,1),\quad b=f-\frac14,
 \quad A(q)=b+\xi,\quad A(r)=b+\xi_0,
 \quad\Delta=\xi_0-\xi.
\]
Only $R_n<0$ can remain.  The exact shelf-curvature identity then gives
\begin{equation}
 \xi_0+\xi<\frac12,\qquad0\leq\xi<\frac14,\qquad
 0\leq\Delta<\frac12,\qquad A(q)<f,\qquad A(r)<f+\frac14.
 \label{eq:no-drop-residual-phases}
\end{equation}
In addition to $K<K_{\rm nd}(n,Q)$, three elementary action estimates are
needed.  The radius derivative and $K-q=\delta+\alpha<\delta+1$ give
\begin{equation}
 K\leq U_f(\delta):=
 \frac{2\delta^2(\delta+1)}{\pi^2b^2}.
 \label{eq:no-drop-upper-geometry}
\end{equation}
The slope representation
\[
 \sigma(z)=\frac1\pi\int_\mu^K
 \frac{z}{a\sqrt{a^2-z^2}}\dd a
 \geq\frac{\delta z}{\pi K^2}
\]
and \eqref{eq:no-drop-residual-phases} imply
\begin{equation}
 b\,n(2r+n)<K^2,\qquad b\,n(n+1)<K^2.
 \label{eq:no-drop-slope-geometry}
\end{equation}
Finally, \eqref{eq:pre-shelf-drop} with $p=n$ gives
\begin{equation}
 \frac\delta\pi<f+\frac14+\frac r{4n}.
 \label{eq:no-drop-width-geometry}
\end{equation}

If $r\leq K/2$, then $r<\delta+n+1$, and hence
\[
 \delta<D_{n,f}:=
 \frac{f+1/2+1/(4n)}{1/\pi-1/(4n)}.
\]
Therefore every residual central pair satisfies
\begin{equation}
 b\,n(n+1)<U_f(D_{n,f})^2.
 \label{eq:no-drop-central-obstruction}
\end{equation}
The left side increases with $n$ and the right side decreases.  Also
\[
 S_Q\geq\frac32\left((Q+\tfrac34)^{2/3}
 -(\tfrac34)^{2/3}\right)\geq\frac54Q^{2/3}\quad(Q\geq7).
\]
For $f\geq16$ this first forces $n>f/6$; then
$D_{n,f}<5f$, $U_f(D_{n,f})<30f$, $n<31\sqrt f$, and thus
$f<34596$.  The remaining monotone endpoint comparisons are finite.
Outward Arb arithmetic verifies, for both $Q=f-1,f$, that every
$50\leq f\leq34595$ is excluded and leaves exactly 649 necessary pairs
before the half-unit reserve is used.  Its two $n=1$ pairs are
$(f,n)=(2,1),(3,1)$; Corollary~\ref{cor:no-drop-n-one} removes them and
leaves $647$ necessary pairs inside \eqref{eq:no-drop-central-box}.

If $r>K/2$, the inequality $R_n<0$ and monotonicity of $\sigma$ give
$\sigma(r)<1/(2n)$.  The radius representation also gives
\[
 \sigma(r)>
 \frac{\delta}{2\pi\sqrt{2K(\delta+n+1)}},
\]
and therefore
\begin{equation}
 K>\frac{\delta^2n^2}{2\pi^2(\delta+n+1)},
 \qquad
 \delta+1>c_fn,\quad
 c_f=\frac{\sqrt{1+b^2}-1}{2}.
 \label{eq:no-drop-outer-geometry}
\end{equation}
With $\delta_{n,f}=\max\{\pi b,c_fn-1\}$, every residual outer pair must
satisfy
\begin{equation}
 \frac{\delta_{n,f}^2n^2}
 {2\pi^2(\delta_{n,f}+n+1)}<K_{\rm nd}(n,Q).
 \label{eq:no-drop-outer-obstruction}
\end{equation}
The same sum bound gives $f\leq66$.  Once $c_fn-1$ is active, the left
side of the normalized comparison increases with $n$, while its right side
decreases.  Its limiting coefficients are
\[
 \gamma_f=\frac{c_f^2}{2\pi^2(c_f+1)},
 \qquad \beta_J=\frac3{2\pi S_J^3}.
\]
The certified signs
$\gamma_2<\beta_2$, $\gamma_3<\beta_3$, and
$\gamma_4>\beta_3$ leave only $f=2,3$ as possible unbounded rows.
For $4\leq f\leq66$, the finite monotone replay leaves exactly 60
necessary pairs in \eqref{eq:no-drop-outer-box}; none has $n=1$.

For reproducibility, all outward-rounded comparisons just invoked are in
\path{scripts/general_d_no_drop_exhaustion_verify.py}.  At 256-bit
precision it performs 138,184 central endpoint comparisons and 980 outer
comparisons, and reports the pre-reserve counts 649 and 60 above.  The same
certificate was independently replayed through 1024 bits; the Round 5
replay confirms the updated counts 647 and 60.

It remains to rule out escape in $f=2,3$.  On the upper half of
$[\mu,K]$,
\[
 A(q)\geq\frac\delta{2\pi}\sqrt{\frac\delta{2K}}.
\]
Together with $A(q)<f$, \eqref{eq:no-drop-cutoff}, and
\eqref{eq:no-drop-outer-geometry}, this proves
$\delta\asymp n$ and $K\asymp n^3$.  Set
\[
 \eta=\frac\delta K,\qquad
 \kappa=K\eta^{3/2},\qquad N=n\sqrt\eta.
\]
This is exactly the critical scaling \eqref{eq:critical-grazing-ray}.
For an explicit transfer, define $C_K(f)$ and $C_\delta(f)$ as in
Theorem~\ref{thm:quantitative-no-drop-transfer}.  If
$n\geq\lceil2/c_f\rceil$, then
\eqref{eq:no-drop-outer-geometry}, the failed cutoff, and
$\delta^3<8\pi^2f^2K$ give precisely
\[
 \frac{c_f}{2}n\leq\delta\leq C_\delta(f)n,\qquad
 \frac{c_f^2}{8\pi^2(C_\delta(f)+2)}n^3
 <K<C_K(f)n^3.
\]
The endpoint-action bound $A(q)\geq2\kappa/(3\pi)$ and $A(q)<f$ give
$\kappa<3\pi f/2$.  Since the cone also gives
$\delta\geq c_fn/2$,
\begin{equation}
 s=\frac\kappa\delta<\frac{3\pi f}{c_fn}<\frac{38}{n}
 \qquad(f=2,3).
 \label{eq:no-drop-round-five-s-bound}
\end{equation}
Indeed, the two coefficients are
$48\pi/(\sqrt{65}-4)$ and $72\pi/(\sqrt{137}-4)$, both strictly smaller
than $38$.
Thus Theorem~\ref{thm:quantitative-no-drop-transfer} excludes every
residual candidate with $n\geq380$, since then $s<1/10$.  Candidates below
the preliminary threshold $\lceil2/c_f\rceil$ already satisfy this bound.
The two low-floor rows therefore satisfy the explicit bound $n<380$.
\end{proof}

\begin{theorem}[Certified closure of the no-drop floors $f\geq2$]
\label{thm:no-drop-high-floors-complete}
Suppose
\[
 p=n\geq1,\qquad F_0=\cdots=F_n=f\geq2.
\]
Then
\begin{equation}
 \boxed{\mathcal S_n=D_A(q)+R_n>0.}
 \label{eq:no-drop-high-floors-complete}
\end{equation}
Consequently \eqref{eq:no-drop-wall-slack} gives
$D_A(r)=\mathcal S_n+\chi_q>0$.
\end{theorem}

\begin{proof}
The case $n=1$ is Corollary~\ref{cor:no-drop-n-one}.  Let $n\geq2$.

For $f=2,3$, Theorem~\ref{thm:quantitative-no-drop-transfer} and the
geometric reductions in the proof of
Theorem~\ref{thm:no-drop-global-exhaustion} cover the sector
\[
 s:=\sqrt{\frac{K-\mu}{K}}\leq\frac1{10}.
\]
They also show that $n\geq380$ forces $s<38/n\leq1/10$.  It remains, for
$2\leq n\leq379$, to cover $s\geq1/10$.  The directed-Arb certificate
uses the exact coordinates
\[
 K=\frac\kappa{s^3},\qquad
 \mu=\frac{\kappa(1-s^2)}{s^3},\qquad
 q=\mu-\alpha,\qquad r=q-n,
\]
and retains three one-sided strict phases: the corner
$\alpha=0$, the lower endpoint wall, and the open side reduced
monotonically to that wall with its open-side count data unchanged.  It
certifies all $378\cdot3=1134$ phases on each floor.  The equality
$s=1/10$ is included in both the analytic and interval modules.

For $f\geq4$, the central exhaustion leaves $595$ pairs and the outer
exhaustion leaves $60$, with an intersection of $45$.  Thus their exact
union has $610$ pairs.  The same corner, lower, and open one-sided phases
give $610\cdot3=1830$ cases.  The eight-shard interval certificate proves
that every terminal box either violates a proved necessary condition or
has a strictly positive outward Arb lower bound for $\mathcal S_n$.  Its
aggregate checker verifies the exact pair union, phase coverage, forest
identities, shard assignments, and empty error streams.

The complete finite ledger is
\begin{center}
\small
\begin{tabular}{@{}lrrrr@{}}
\toprule
floor & phase cases & processed boxes & positive leaves & maximum depth\\
\midrule
$f=2$    & $1134$ & $19{,}887{,}149$ & $29{,}323$ & $31$\\
$f=3$    & $1134$ & $ 7{,}695{,}182$ & $10{,}921$ & $23$\\
$f\geq4$ & $1830$ & $ 8{,}232{,}800$ & $22{,}986$ & $15$\\
\midrule
total    & $4098$ & $35{,}815{,}131$ & $63{,}230$ & $31$\\
\bottomrule
\end{tabular}
\end{center}
The smallest reported positive endpoints in the three rows are,
respectively,
\[
 1.7050400932\times10^{-6},\qquad
 2.4999770759\times10^{-4},\qquad
 2.1442116136\times10^{-5}.
\]
These decimals are diagnostics only; every proof decision uses exact
rational arithmetic or a directed interval sign.

\begingroup\scriptsize
\noindent For reproducibility, the frozen low-floor source and hash are
\par\noindent
\path{scripts/general_d_no_drop_low_floor_final_verify.py}
\par\noindent
SHA-256
\path{4AD4BAD7CECE561102B84C4983BD4C50D1BC39D59DA4D1ECD3D04BC7ADA9319F}.
The high-floor replay preserves
\par\noindent
\path{scripts/general_d_no_drop_compact_verify.py}
\par\noindent
at SHA-256
\path{EEF8150E8D4D56DC1F0088E0C1F3C2EAE3D7E93D688AB52F2CF5975416FACCDE}
and is accepted by
\par\noindent
\path{scripts/_aggregate_general_d_no_drop_high_floor_shards.py}
\par\noindent
at SHA-256
\path{79999E43B6D0CBC7BF29570F720FF5AE2FBFA4A986431D79F9716966426B7996}.
The independent proof-boundary audits are
\par\noindent
\path{human/outbox/general-d-round-05-no-drop-compact-final-independent-audit.md}
\par\noindent
with SHA-256
\path{83ABBE4DB8228B06678CDE41FED4B02281C154A711ABC9D26FF204FCC47B5475},
and
\par\noindent
\path{human/outbox/general-d-round-08-no-drop-high-floor-final-note-independent-audit.md}
\par\noindent
with SHA-256
\path{7DB8123BD3577E0CF9558F8BB10D8D5591821AF3DA0FFA234B35E550982CD747}.
They pin the remaining verifier, transcript, coverage, and aggregate-audit
hashes.
\endgroup
This proves \eqref{eq:no-drop-high-floors-complete}.
\end{proof}

\section{A dimension-dependent no-mode region}

The radial form \eqref{eq:radial-operator}, the one-dimensional Dirichlet
Poincar\'e inequality, and $r^{-2}\geq1$ on $(\rho,1)$ give
\[
 \lambda_{\ell,n}\geq
 \frac{n^2\pi^2}{(1-\rho)^2}+\nu_{\ell,d}^2-\frac14.
\]
Since $\nu_{0,d}=(d-2)/2$, one obtains the strengthened wall
\begin{equation}
 N_D(A_{\rho,1}^{(d)},K^2)=0
 \quad\text{if}\quad
 K^2\leq\frac{\pi^2}{(1-\rho)^2}
 +\frac{(d-2)^2-1}{4}.
 \label{eq:no-mode}
\end{equation}
The equality face belongs to the no-mode region because the counting
function is strict.

In particular, for the general-dimensional P\'olya theorem it is enough
to prove the shifted-tail estimate in the universal active region
\begin{equation}
 K>\frac\pi{1-\rho}.
 \label{eq:universal-active-region}
\end{equation}
Indeed, the complementary region has no eigenvalues already in dimension
three, and the additional centrifugal term only enlarges it for $d>3$.
This observation does not prove the stronger standalone
Conjecture~\ref{conj:shifted-tail} below the wall, but it removes that
region from the spectral critical path.

\section{Remaining proof obligations}

The first interface-cell milestone is complete, and
Proposition~\ref{prop:fixed-rho-high-energy} proves that for each fixed
$\rho$ only a bounded frequency interval remains; in the thin sector
Proposition~\ref{prop:thin-high-energy} improves the sufficient threshold to
$K\ge 125/[8(1-\rho)^2]$.  Proposition
~\ref{prop:first-shelf-reduction} further replaces the many-cell dynamics
by the single scalar \eqref{eq:terminal-shelf-scalar}.  For the spectral
theorem this scalar is needed only in \eqref{eq:universal-active-region}.
Lemma~\ref{lem:scale-free-shelf-curvature} sharpens its head term, while
Corollary~\ref{cor:exact-angle-terminal-reserve} replaces qualitative
terminal nonnegativity by a fractional-top positive sum, and
Corollary~\ref{cor:zero-level-terminal-triangle} covers the empty-level
case.  Proposition~\ref{prop:ultrathin-high-floor-sector} closes an
explicit high-floor finite-shell sector, while
Theorem~\ref{thm:critical-ray} rules out a negative leading obstruction on
the grazing ray.  For $p=n$, Proposition~\ref{prop:no-drop-identities}
identifies the exact endpoint phase and
 Theorem~\ref{thm:no-drop-global-exhaustion} proves a global finite
 reduction for $f\geq2$, and
 Theorem~\ref{thm:no-drop-high-floors-complete} certifies every retained
 chamber.  The distinct first-floor no-drop case $f=1$, $p=n$, is not part
 of that exhaustion.
 Theorems~\ref{thm:first-floor-global-exhaustion} and
 \ref{thm:high-floor-first-drop-exhaustion} remove every noncompact
 first-drop escape.  Theorem~\ref{thm:first-floor-s-cone}, together with
 Theorems~\ref{thm:first-floor-large-ray}--
 \ref{thm:first-floor-small-starts}, closes every ordinary first-floor
 first-drop active wall; \eqref{eq:no-mode} removes the below-active region from the
 spectral critical path.  Theorem~\ref{thm:high-floor-explicit-compact}
 gives a first explicit high-floor reduction, and
 Theorem~\ref{thm:high-floor-global-scalar} supersedes its enormous cutoff:
 it first reduced every negative candidate to the Round~7 box
 \eqref{eq:high-floor-round-seven-continuous}--
 \eqref{eq:high-floor-round-seven-discrete}.  The extended
 Theorem~\ref{thm:high-floor-round-eight-scalar} next supersedes that box,
 and the terminally strengthened
 Theorem~\ref{thm:high-floor-round-eight-b-scalar} supersedes it again,
 Theorem~\ref{thm:high-floor-round-eight-d-scalar} sharpens it once more,
 Theorem~\ref{thm:high-floor-round-eight-e-scalar} improves the terminal
 cap and absent-level chord,
 Theorem~\ref{thm:high-floor-round-eight-f-scalar} crosses the common
 prefix sign wall, and
 Theorem~\ref{thm:high-floor-round-eight-g-scalar} gives the current
 sharpest box:
 every negative candidate lies in
 \eqref{eq:high-floor-round-eight-g-continuous}--%
 \eqref{eq:high-floor-round-eight-g-discrete}.  This sharper reduction still
 does not certify the residual compact walls.  The primary route now has
 the following precise steps.
\begin{enumerate}
 \item Prove the first-floor no-drop case
 $F_0=\cdots=F_n=1$, $p=n$.  This row is nonempty and is not contained in
 the ordinary first-floor first-drop certificates, whose hypothesis is
 $p<n$.
 \item In the high-floor $p<n$ branch,
 Theorem~\ref{thm:high-floor-global-scalar} historically closes
 $\rho>99/100$, and
 Theorem~\ref{thm:high-floor-round-eight-scalar} extends the exclusion to
 $\rho\geq603/625$, and
 Theorem~\ref{thm:high-floor-round-eight-b-scalar} extends it to
 $\rho\geq477/500$, and
 Theorem~\ref{thm:high-floor-round-eight-d-scalar} uses the active width
 and sharp terminal cap to extend it further to $\rho\geq93/100$, and
 Theorem~\ref{thm:high-floor-round-eight-e-scalar} combines the automatic
 $q\geq5/2$ cap with a sharper absent-level chord to reach
 $\rho\geq927/1000$, and
 Theorem~\ref{thm:high-floor-round-eight-f-scalar} crosses the sign wall of
 the common prefix by using $M\leq\mu$ and reaches
 $\rho\geq1847/2000$, and
 Theorem~\ref{thm:high-floor-round-eight-g-scalar} retains the exact
 correlated factor, sharpens the active width and small-level distance,
 and reaches $\rho\geq23/25$.
 It leaves only
 \[
  \frac1{3154914}<\rho<\frac{23}{25},\qquad
  \frac{21}{4}<K<1577457,
 \]
 with the exact integer ranges in
 \eqref{eq:high-floor-round-eight-g-discrete}.  Every fixed-ratio compact
 analytic wall inside those bounds is still uncertified; it requires
 sharper analytic localization or directed-rounding certificates, not a
 claim of closure from finiteness alone.  There is no remaining ordinary
 first-floor first-drop wall.
 \item Insert the resulting all-frequency shifted-tail theorem into
 Theorem~\ref{thm:conditional}.  No dimension-dependent angular or optical
 theorem is then needed.
\end{enumerate}
If either residual branch fails only for an explicitly controlled family of tails,
the exact identity \eqref{eq:master-defect} is the appropriate aggregate
fallback: the weighted negative part must be compared directly with
$\mathcal B_d(A)$.

\begin{thebibliography}{9}

\bibitem{FLPSballs}
N.~Filonov, M.~Levitin, I.~Polterovich, and D.~A.~Sher,
\newblock P\'olya's conjecture for Euclidean balls,
\newblock \emph{Invent. Math.} \textbf{234} (2023), 129--169.
\newblock \url{https://doi.org/10.1007/s00222-023-01198-1}.

\bibitem{FLPSphase}
N.~Filonov, M.~Levitin, I.~Polterovich, and D.~A.~Sher,
\newblock Uniform enclosures for the phase and zeros of Bessel functions
and their derivatives,
\newblock \emph{SIAM J. Math. Anal.} \textbf{56} (2024), no.~6,
7644--7682.
\newblock \url{https://doi.org/10.1137/24M1642032}.

\bibitem{FLPSannuli}
N.~Filonov, M.~Levitin, I.~Polterovich, and D.~A.~Sher,
\newblock P\'olya's conjecture for Dirichlet eigenvalues of annuli,
\newblock \emph{J. Lond. Math. Soc.} \textbf{113} (2026), no.~2, e70425.
\newblock \url{https://doi.org/10.1112/jlms.70425}.

\end{thebibliography}

\end{document}
